\chapter{实数}
\label{cha:real-numbers}

\index{real numbers|(}%
任何一个名副其实的数学基础，最终都必须能够构造出数学分析中所理解的那种实数，也就是作为一个完备的 Archimedean 有序域。
\index{ordered field}%
完备性有两种概念。Cauchy（柯西）的版本要求实数对柯西序列\index{Cauchy!sequence}的极限封闭，而更强的 Dedekind版本则要求对Dedekind 分割\index{cut!Dedekind}封闭。
这引出了两种构造实数的方法，我们将分别在\cref{sec:dedekind-reals}和\cref{sec:cauchy-reals}中研究。在\cref{RD-final-field,RC-initial-Cauchy-complete}中，我们用泛性质来刻画这两种构造：Dedekind 实数是终 Archimedean 有序域，而 Cauchy 实数是始 Cauchy 完备 Archimedean 有序域。

在传统的构造性数学\index{mathematics!constructive}中，实数似乎总是需要某些妥协。
例如，Dedekind 实数与幂集或某些其他形式的非直谓性配合得更好，而 Cauchy 实数在可数选择公理\index{axiom!of choice!countable}成立时表现良好。
然而，我们给出了 Cauchy 实数的一种新构造，将它定义为一个高阶归纳-归纳类型；这似乎提供了第三种可能，它既不需要幂集，也不需要可数选择公理。

在~\cref{sec:comp-cauchy-dedek}中，我们比较这两种实数构造。
Cauchy 实数包含在Dedekind 实数中。
若排中律或可数选择公理成立，则二者一致；但在一般情况下，这一包含可能是真包含。

在~\cref{sec:compactness-interval}中，我们考察闭区间~$[0,1]$ 的三种紧性概念。
我们首先证明 $[0,1]$ 是\emph{度量紧}\indexdef{metricly compact}\indexdef{compactness!metric}的，意思是它既完备又完全有界，并且度量紧空间上的一致连续映射\index{uniformly continuous}的行为符合预期。
与之形成对比的是，Bolzano（波尔查诺）--Weierstra\ss{}（魏尔斯特拉斯）性质——即每个序列都有收敛子列——蕴涵有限全知原理，而后者是排中律的一个特例。
最后，我们讨论 Heine（海涅）--Borel（博雷尔）紧性。
有限子覆盖性质的朴素表述行不通，但一种证明相关的归纳覆盖\index{inductive cover}概念则可行。
本节内容基本属于标准的构造性分析。

在同伦类型论中发展起来的实数与分析理论，可以轻松地与经典数学兼容。
只要假设排中律与选择公理，我们便得到标准的经典分析\index{mathematics!classical}\index{classical!analysis}：此时Dedekind 实数与 Cauchy 实数一致，关于Dedekind 实数非直谓本质的基础问题随之消失，区间也尽可能地紧了。

在 \cref{sec:surreals} 中，我们通过将 Conway（康威）的超现实数构造为高阶归纳-归纳类型来结束本章；
这一构造在泛等类型论中比在经典集合论中更为自然。

除了 \cref{cha:basics,cha:logic} 的基本理论之外，如上所述，我们对 Cauchy 实数和超现实数使用了``高阶归纳-归纳类型''：它们结合了 \cref{cha:hits} 的思想与 \cref{sec:generalizations} 中提到的归纳-归纳类型概念。
我们也将频繁使用 \cref{subsec:prop-trunc} 中描述的传统逻辑记号，以及（在 \cref{sec:piw-pretopos} 中证明的）事实：我们的``集合''的行为方式符合预期。

注意，圆周的万有覆叠的全空间（它在 \cref{subsec:pi1s1-homotopy-theory} 中扮演了类似于经典代数拓扑中``实数''的角色）\emph{不是}我们要找的实数类型。
该类型是可缩的，因此等价于单点类型，所以它不能配备非平凡的代数结构。

\section{有理数域}
\label{sec:field-rati-numb}

\indexdef{rational numbers}%
\indexsee{number!rational}{rational numbers}%
我们首先构造有理数 \Q，因为实数可以看作 \Q 的完备化。
专家会指出，\Q 可以被任何近似域替代，
\indexdef{field!approximate}%
即 \Q 的一个子环，其中存在任意精确的近似逆
\index{inverse!approximate}%
。
一个例子是2进有理数环，
\index{rational numbers!dyadic}%
即形如 $n/2^k$ 的有理数。
如果我们在计算机上实现构造性数学，近似域会更合适，但我们将这种精细之处留给那些在意~$\pi$ 的数字的人。

我们在 \cref{sec:set-quotients} 中将整数 \Z 构造为 $\N\times \N$ 的一个商集，并观察到这个商集由一个幂等元生成。
在 \cref{sec:free-algebras} 中我们看到 \Z 是 \unit 上的自由群；类似地，可以证明它也是 \emptyt 上的自由交换环\index{ring}。
有理数域 \Q 的构造也遵循同样的思路，即作为商集
%
\[ \Q \defeq (\Z \times \N)/{\approx} \]
%
其中
\[ (u,a) \approx (v,b) \defeq (u (b + 1) = v (a + 1)). \]
%
换言之，对子 $(u, a)$ 代表有理数 $u / (1 + a)$。这里不可能出现除以零的情况，因为我们巧妙地给分母~$a$ 加了 1。
这里同样有典范的代表元选择方式，即最简分数。因此，我们可以应用 \cref{lem:quotient-when-canonical-representatives} 来得到一个集合 \Q，它同样具有可判定相等性。
\index{decidable!equality}%

我们不专门写出 \Q 上的算术运算，相信读者即使面对分母加了 1 的情形，也知道如何计算分数。
我们仅记录如下结论：有理数 \Q 的有序域构造是完全不成问题的，它具有可判定的相等性与可判定的序关系。
它也可刻画为初始有序域。
\index{initial!ordered field}%

\symlabel{positive-rationals}
\indexdef{rational numbers!positive}%
\indexdef{positive!rational numbers}%
最后，我们将用 $\Qp \defeq \setof{ q : \Q | q > 0 }$ 表示正有理数的类型。

\section{Dedekind 实数}
\label{sec:dedekind-reals}

\index{real numbers!Dedekind|(}%
我们首先回顾 Dedekind构造的基本思想。
我们采用双侧Dedekind 分割，而非常用的单侧版本，因为对称性使构造更为优雅，且在构造性数学与经典数学中均适用。
\index{mathematics!constructive}%
一个\emph{Dedekind 分割}\index{cut!Dedekind} 由 \Q 的一对子集 $(L, U)$ 组成，分别称为\emph{下}分割与\emph{上}分割，它们满足：
%

\begin{enumerate}
\item \emph{有实例的（inhabited）：} 存在 $q \in L$ 和 $r \in U$，
\item \emph{舍入的（rounded）：} $q \in L \Leftrightarrow \exis {r \in \Q} q < r \land r \in L$
  且 $r \in U \Leftrightarrow \exis {q \in \Q} q \in U \land q < r$，
  \index{rounded!Dedekind cut}
\item \emph{不相交的（disjoint）：} $\lnot (q \in L \land q \in U)$，以及
\item \emph{定位的（located）：} $q < r \Rightarrow q \in L \lor r \in U$。
  \index{locatedness}%
\end{enumerate}

%
从两个方向理解``舍入性''条件：自左向右读，它告诉我们分割是\emph{开的（open）}，
\index{open!cut}%
而从右向左读，它说明 $L$ 和 $U$ 分别是\emph{下集（lower set）}与\emph{上集（upper set）}。
定位性条件指出，$L$ 与 $U$ 之间不存在大的间隙。
因为分割总是开的，它们从不包含``中间的那个点''——即使该点为有理数时亦然。
一个典型的Dedekind 分割如下图所示：
%

\begin{center}
  \begin{tikzpicture}[x=\textwidth]
    \draw[<-),line width=0.75pt] (0,0) -- (0.297,0) node[anchor=south east]{$L\ $};
    \draw[(->,line width=0.75pt] (0.300, 0) node[anchor=south west]{$\ U$} -- (0.9, 0) ;
  \end{tikzpicture}
\end{center}

%
我们或许会朴素地将非形式化的定义翻译到类型论中，即认为一个分割就是一对映射 $L, U : \Q \to \prop$。但在 \cref{subsec:prop-subsets} 中我们已经看到，$\prop$ 是 $\prop_{\UU_i}$ 的一种带有\index{typical ambiguity}类型歧义的记法，其中 $\UU_i$ 是某个宇宙。
一旦我们使用某个特定的 $\UU_i$ 来定义分割，实数类型便位于下一个宇宙 $\UU_{i+1}$ 中，实数的一个性质位于更高两层的宇宙 $\UU_{i+2}$ 中，实数子集的一个性质位于 $\UU_{i+3}$ 中，依此类推。
原则上，我们完全可以跟踪这些\index{universe level}宇宙层级，尤其是在证明助手的帮助下；但在此处这样做只会带来我们宁可避免的繁琐事务。
因此，我们将做出一个简化假设：就我们的所有目的而言，一个单一的命题类型 $\Omega$ 就已足够。

事实上，Dedekind 实数的构造对逻辑层面的处理相当稳健。我们可以通过几种方式来理解这里使用单一类型 $\Omega$ 的做法：
%

\begin{enumerate}

\item 我们可以将 $\Omega$ 等同于那个有歧义的 $\prop$，并追踪定义和构造中出现的所有宇宙。

\item 我们可以假设命题大小重整公理，
  \index{propositional!resizing}%
  如 \cref{subsec:prop-subsets} 所示，
  它本质上将各层的 $\prop_{\UU_i}$ “坍缩”到最低的
  宇宙层级\index{universe level}，我们称其为 $\Omega$。

\item 一位对类型论宇宙或计算的复杂细节不感兴趣的经典数学家，可以直接假设命题的排中律~\eqref{eq:lem}，从而令 $\Omega \jdeq \bool$。
  \index{excluded middle}
  这不仅根除了关于 $\prop$ 的层级\index{universe level}问题，还将我们所做的一切转化为标准的经典\index{mathematics!classical}实数构造。

\item 在谱系的另一端，我们或许会问：让这些构造得以成立的最低要求是什么。一个谓词成为Dedekind 分割的条件，可以只用合取、析取以及对可数集~\Q 的存在量词\index{quantifier!existential}来表达。因此，我们可以取 $\Omega$ 为初始 \emph{$\sigma$-框架}，
  \index{initial!sigma-frame@$\sigma$-frame}%
  \index{sigma-frame@$\sigma$-frame!initial|defstyle}%
  即一个带有可数并\index{join!in a lattice}的格\index{lattice}，其中二元交运算对可数并满足分配律。（初始 $\sigma$-框架不可能是双点格 $\bool$，因为除非假设排中律，否则 $\bool$ 对可数并不封闭。）这将导致把~$\Omega$ 构造为一个高阶归纳-归纳类型，但像 \cref{sec:cauchy-reals} 那样的一个实验已经足够了。
\end{enumerate}

在上述所有情况下，$\Omega$ 都是一个集合。
%
事不宜迟，我们现在就将这个非形式的定义翻译为类型论的形式。在本章中，我们使用的逻辑记号遵循 \cref{defn:logical-notation}。

\begin{defn} \label{defn:dedekind-reals}
  一个 \define{Dedekind 分割（Dedekind cut）}
  \indexsee{Dedekind!cut}{cut, Dedekind}%
  \indexdef{cut!Dedekind}%
  是一对谓词 $L : \Q \to \Omega$ 与 $U : \Q \to \Omega$ 构成的 $(L, U)$，满足以下条件：
  %
  \begin{enumerate}
  \item \label{defn:dedekind-reals-inhabited}
    \emph{有实例的（即有界的）：} $\exis{q : \Q} L(q)$ 且 $\exis{r : \Q} U(r)$，
  \item \emph{舍入的（rounded）：} 对所有 $q, r : \Q$，
    \index{rounded!Dedekind cut}
    %
    \begin{align*}
      L(q) &\Leftrightarrow \exis{r : \Q} (q < r) \land L(r)
      \qquad\text{and}\\
      U(r) &\Leftrightarrow \exis{q : \Q} (q < r) \land U(q),
    \end{align*}
  \item \emph{不相交的（disjoint）：} $\lnot (L(q) \land U(q))$ 对所有 $q : \Q$，
  \item \emph{定位的（located）：} $(q < r) \Rightarrow L(q) \lor U(r)$ 对所有 $q, r : \Q$。
  \index{locatedness}%
  \end{enumerate}
  %
  我们用 $\dcut(L, U)$ 表示这些条件的合取。\define{Dedekind 实数（Dedekind reals）}的类型定义为
  \indexsee{Dedekind!real numbers}{real numbers, De\-de\-kind}%
  \indexdef{real numbers!Dedekind}%
  %
  \begin{equation*}
    \RD \defeq \setof{ (L, U) : (\Q \to \Omega) \times (\Q \to \Omega) | \dcut(L,U)}.
  \end{equation*}
\end{defn}

显然 $\dcut(L, U)$ 是一个命题，并且由于 $\Q \to \Omega$ 是一个集合，Dedekind 实数也构成一个集合。关于Dedekind 分割的变体（这些变体会导出扩充实数、下实数、上实数以及区间域），请参见
\cref{ex:RD-extended-reals,ex:RD-lower-cuts,ex:RD-interval-arithmetic}。

存在一个嵌入 $\Q \to \RD$，它将每个有理数 $q : \Q$ 与分割 $(L_q, U_q)$ 关联起来，其中
%
\begin{equation*}
  L_q(r) \defeq (r < q)
  \qquad\text{and}\qquad
  U_q(r) \defeq (q < r).
\end{equation*}
%
对于与有理数关联的分割 $(L_q, U_q)$，我们简记为 $q$。

\subsection{Dedekind 实数的代数结构}
\label{sec:algebr-struct-dedek}

Dedekind 实数的代数结构及序结构的构造，在直觉主义逻辑中照常进行。我们不再赘述细节，而是指出经典\index{mathematics!classical}构造与直觉主义构造之间的差异。用 $L_x$ 和 $U_x$ 表示实数 $x : \RD$ 的下分割与上分割，我们将加法定义为%
%
\indexdef{addition!of Dedekind reals}%
\begin{align*}
  L_{x + y}(q) &\defeq \exis{r, s : \Q} L_x(r) \land L_y(s) \land q = r + s, \\
  U_{x + y}(q) &\defeq \exis{r, s : \Q} U_x(r) \land U_y(s) \land q = r + s,
\end{align*}
%
加法逆元则定义为
%
\begin{align*}
  L_{-x}(q) &\defeq \exis{r : \Q} U_x(r) \land q = - r, \\
  U_{-x}(q) &\defeq \exis{r : \Q} L_x(r) \land q = - r.
\end{align*}
%
在这些运算下，$(\RD, 0, {+}, {-})$ 构成一个交换\index{group!abelian}群。乘法的定义则略显繁琐：
%
\indexdef{multiplication!of Dedekind reals}%
\begin{align*}
  L_{x \cdot y}(q) &\defeq
  \begin{aligned}[t]
    \exis{a, b, c, d : \Q} & L_x(a) \land U_x(b) \land L_y(c) \land U_y(d) \land {}\\
                           & \qquad q < \min (a \cdot c, a \cdot d, b \cdot c, b \cdot d),
  \end{aligned} \\
  U_{x \cdot y}(q) &\defeq
  \begin{aligned}[t]
    \exis{a, b, c, d : \Q} & L_x(a) \land U_x(b) \land L_y(c) \land U_y(d) \land {}\\
                           & \qquad \max (a \cdot c, a \cdot d, b \cdot c, b \cdot d) < q.
  \end{aligned}
\end{align*}
%
\index{interval!arithmetic}%
这些公式与区间算术中的区间乘法有关。在区间算术中，端点均为有理数的区间 $[a,b]$ 和 $[c,d]$ 相乘得到区间
%
\begin{equation*}
  [a,b] \cdot [c,d] =
  [\min(a c, a d, b c, b d), \max(a c, a d, b c, b d)].
\end{equation*}
%
例如，下分割的公式可以理解为：当存在分别包含 $x$ 和 $y$ 的区间 $[a,b]$ 和 $[c,d]$，使得 $q$ 位于 $[a,b] \cdot [c,d]$ 的左侧时，就有 $q < x \cdot y$。一般而言，将一个满足 $L_x(a)$ 和 $U_x(b)$ 的区间 $[a,b]$ 视为~$x$ 的一种近似是很有用的，参见
\cref{ex:RD-interval-arithmetic}。

我们现在已经得到了一个含幺交换环\index{ring}
\index{unit!of a ring}%
$(\RD, 0, 1, {+}, {-}, {\cdot})$。为了处理乘法逆元，我们必须先引入序。定义 $\leq$ 和 $<$ 为
%
\begin{align*}
  (x \leq y) &\ \defeq \ \fall{q : \Q} L_x(q) \Rightarrow L_y(q), \\
  (x < y)    &\ \defeq \ \exis{q : \Q} U_x(q) \land L_y(q).
\end{align*}

\begin{lem} \label{dedekind-in-cut-as-le}
  对于所有 $x : \RD$ 和 $q : \Q$，$L_x(q) \Leftrightarrow (q < x)$ 且 $U_x(q)
  \Leftrightarrow (x < q)$。
\end{lem}

\begin{proof}
  如果 $L_x(q)$，则由舍入性，仅存在 $r > q$ 使得 $L_x(r)$ 成立；而 $U_q(r)$ 成立，由此可得 $q < x$。
  反之，如果 $q < x$，则存在 $r : \Q$ 使得 $U_q(r)$ 和 $L_x(r)$ 同时成立；又因为 $L_x$ 是下集，所以 $L_x(q)$ 成立。
  另一半证明是对称的。
\end{proof}

\index{partial order}%
\index{transitivity!of . for reals@of $<$ for reals}
\index{transitivity!of . for reals@of $\leq$ for reals}
\index{relation!irreflexive}
\index{irreflexivity!of . for reals@of $<$ for reals}
关系 $\leq$ 是偏序，而 $<$ 传递且不自反。线性序
\index{order!linear}%
\index{linear order}%
%
\begin{equation*}
  (x < y) \lor (y \leq x)
\end{equation*}
%
在假设排中律时成立，但若不假设排中律，则得到弱线性序
%
\index{order!weakly linear}
\index{weakly linear order}
\begin{equation} \label{eq:RD-linear-order}
  (x < y) \Rightarrow (x < z) \lor (z < y).
\end{equation}
%
乍一看，或许不清楚~\eqref{eq:RD-linear-order} 与线性序有什么关系。但如果我们取 $x \jdeq u - \epsilon$ 和 $y \jdeq u + \epsilon$，其中 $\epsilon > 0$，则得到
%
\begin{equation*}
  (u - \epsilon < z) \lor (z < u + \epsilon).
\end{equation*}
%
这就是``在微小数值误差意义下''的线性序。换句话说，既然我们实际上无法期望以无限精度进行计算，那么 $<$ 只能在我们已计算出的有限精度内判定，也就不足为奇了。

为了证明~\eqref{eq:RD-linear-order} 成立，假设 $x < y$。则仅存在 $q : \Q$ 使得 $U_x(q)$ 且 $L_y(q)$。
由舍入性，仅存在 $r, s : \Q$ 使得 $r < q < s$，$U_x(r)$ 且 $L_y(s)$。
于是，由定位性可得 $L_z(r)$ 或 $U_z(s)$。
在第一种情形得到 $x < z$，在第二种情形得到 $z < y$。

在经典数学中，所有非零的数都存在乘法逆元。然而，在没有排中律的情况下，需要一个更强的条件。称两个实数 $x, y : \RD$ 是\define{分离的（apart）}
\indexdef{apartness}%
记作 $x \apart y$，当$(x < y) \lor (y < x)$：
%
\symlabel{apart}
\begin{equation*}
  (x \apart y) \defeq (x < y) \lor (y < x).
\end{equation*}
%
如果 $x \apart y$，则 $\lnot (x = y)$。
其逆命题在假设排中律时成立，但在构造性数学中不可证。
\index{mathematics!constructive}%
事实上，如果 $\lnot (x = y)$ 蕴含 $x \apart y$，则可推出一小部分排中律；参见\cref{ex:reals-apart-neq-MP}。

\begin{thm} \label{RD-inverse-apart-0}
  一个实数可逆，当且仅当，它与 $0$ 分离。
\end{thm}

\begin{rmk}
  我们注意到，一个实数可逆当且仅当它仅可逆。实际上，这在任何环\index{ring}中都成立，因为环是一个集合，而乘法逆元若存在则唯一。相关讨论参见\cref{cor:UC}之后的内容。
\end{rmk}

\begin{proof}
  假设 $x \cdot y = 1$。则仅存在 $a, b, c, d : \Q$ 使得 $a < x < b$，$c < y < d$ 且$0 < \min (a c, a d, b c, b d)$。由 $0 < a c$ 和 $0 < b c$ 可知，$a$、$b$ 和 $c$ 要么全为正，要么全为负。
  因此，要么 $0 < a < x$，要么 $x < b < 0$，于是 $x \apart 0$。

  反过来，如果 $x \apart 0$，则要么 $x > 0$，要么 $x < 0$。
  若 $x > 0$，我们如下定义 $x^{-1}$：
  %
  \begin{align*}
    L_{x^{-1}}(q) &\defeq
    (q > 0) \Rightarrow \exis{r : \Q} U_x(r) \land (q r < 1),
    \\
    U_{x^{-1}}(q) &\defeq
    (q > 0) \land \exis{r : \Q} L_x(r) \land (q r > 1).
  \end{align*}
  %
  若 $x < 0$，则如下定义：
  %
  \begin{align*}
    L_{x^{-1}}(q) &\defeq
    (q < 0) \land \exis{r : \Q} U_x(r) \land (q r > 1),
    \\
    U_{x^{-1}}(q) &\defeq
    (q < 0) \Rightarrow \exis{r : \Q} L_x(r) \land (q r < 1). \qedhere
  \end{align*}
\end{proof}

\index{ordered field!archimedean}%
\index{dense}%
\indexsee{order-dense}{dense}%
Archimedean 原理可以有多种陈述方式。我们发现，以``$\Q$ 在 $\RD$ 中稠密''这种形式来陈述，最具启发性。

\begin{thm}[$\RD$ 的 Archimedean 原理] \label{RD-archimedean}
  %
  对于所有 $x, y : \RD$，如果 $x < y$，则仅存在 $q : \Q$ 使得 $x < q < y$。
\end{thm}

\begin{proof}
  由 $<$ 的定义直接得出。
\end{proof}

在讨论Dedekind 实数的完备性之前，让我们先精确地说明它们所具有的代数结构。在下面的定义中，我们的目标不是追求公理化的最小化，而是希望获得实用的结构与性质。

\begin{defn} \label{ordered-field} 一个\define{有序域}
  \indexdef{ordered field}%
  \indexsee{field!ordered}{ordered field}%
  是一个集合 $F$，连同常量 $0$、$1$，运算 $+$、$-$、$\cdot$、$\min$、$\max$，以及纯粹关系 $\leq$、$<$、$\apart$，使得下列条件成立：
  %
  \begin{enumerate}
  \item $(F, 0, 1, {+}, {-}, {\cdot})$ 构成一个带单位的交换环；
    \index{unit!of a ring}%
    \index{ring}%
  \item $x : F$ 可逆当且仅当 $x \apart 0$；
  \item $(F, {\leq}, {\min}, {\max})$ 构成一个格；
  \item 严格序 $<$ 是传递的、不自反的，
    \index{relation!irreflexive}
    \index{irreflexivity!of . in a field@of $<$ in a field}%
    且弱线性的（$x < y \Rightarrow x < z \lor z < y$）；\index{transitivity!of . in a field@of $<$ in a field}
    \index{order!weakly linear}
    \index{weakly linear order}
    \index{strict!order}%
    \index{order!strict}%
  \item 分离关系 $\apart$ 是不自反的、对称的且协传递的（$x \apart y \Rightarrow x \apart z \lor y \apart z$）；
    \index{relation!irreflexive}
    \index{irreflexivity!of apartness}%
    \indexdef{relation!cotransitive}%
    \index{cotransitivity of apartness}%
  \item 对所有 $x, y, z : F$：
    %
    \begin{align*}
      x \leq y &\Leftrightarrow \lnot (y < x), &
      x < y \leq z &\Rightarrow x < z, \\
      x \apart y &\Leftrightarrow (x < y) \lor (y < x), &
      x \leq y < z &\Rightarrow x < z, \\
      x \leq y &\Leftrightarrow x + z \leq y + z, &
      x \leq y \land 0 \leq z &\Rightarrow x z \leq y z, \\
      x < y &\Leftrightarrow x + z < y + z, &
      0 < z \Rightarrow (x < y &\Leftrightarrow x z < y z), \\
      0 < x + y &\Rightarrow 0 < x \lor 0 < y, &
      0 &< 1.
    \end{align*}
  \end{enumerate}
  %
  每个这样的域都有一个标准嵌入 $\Q \to F$。一个有序域是\define{Archimedean 的}
  \indexdef{ordered field!archimedean}%
  \indexsee{archimedean property}{ordered field, archi\-mede\-an}%
  ，如果对所有 $x, y : F$，当 $x < y$ 时纯粹存在 $q :
  \Q$ 使得 $x < q < y$。
\end{defn}

\begin{thm} \label{RD-archimedean-ordered-field}
  Dedekind 实数构成一个有序的 Archimedean 域。
\end{thm}

\begin{proof}
  我们省略证明，希望迄今为止所展示的内容足以使这一定理显得可信。
\end{proof}

\subsection{Dedekind 实数是 Cauchy 完备的}
\label{sec:RD-cauchy-complete}

回忆一下，$x : \N \to \Q$ 称为一个\emph{柯西序列}\indexdef{Cauchy!sequence}，如果它满足
%
\begin{equation} \label{eq:cauchy-sequence}
  \prd{\epsilon : \Qp} \sm{n : \N} \prd{m, k \geq n} |x_m - x_k| < \epsilon.
\end{equation}
%
请注意，我们\emph{没有}截断内部的存在量词，因为我们实际上想要计算收敛速率——一个没有误差估计的逼近几乎不提供任何有用信息。由 \cref{thm:ttac}，\eqref{eq:cauchy-sequence} 给出一个函数 $M
: \Qp \to \N$，称为\emph{收敛模（modulus of convergence）}\indexdef{modulus!of convergence}，使得 $m, k \geq M(\epsilon)$
蕴含 $|x_m - x_k| < \epsilon$。由此可得，对所有 $\delta, \epsilon : \Qp$ 均有 $|x_{M(\delta/2)} - x_{M(\epsilon/2)}|<
\delta + \epsilon$。事实上，映射 $(\epsilon \mapsto
x_{M(\epsilon/2)}) : \Qp \to \Q$ 承载了与原始 Cauchy 条件~\eqref{eq:cauchy-sequence} 相同的关于极限的信息。后续我们将使用这些逼近函数而非柯西序列。

\begin{defn} \label{defn:cauchy-approximation}
  一个\define{Cauchy 逼近（Cauchy approximation）}
  \indexdef{Cauchy!approximation}%
  是一个映射 $x : \Qp \to \RD$，它满足
  %
  \begin{equation}
    \label{eq:cauchy-approx}
    \fall{\delta, \epsilon :\Qp} |x_\delta - x_\epsilon| < \delta + \epsilon.
  \end{equation}
  %
  一个 Cauchy 逼近 $x : \Qp \to \RD$ 的\define{极限（limit）}
  \index{limit!of a Cauchy approximation}%
  是一个数 $\ell : \RD$，使得
  %
  \begin{equation*}
    \fall{\epsilon, \theta : \Qp} |x_\epsilon - \ell| < \epsilon + \theta.
  \end{equation*}
\end{defn}

\begin{thm} \label{RD-cauchy-complete}
  $\RD$ 中的每个 Cauchy 逼近都有一个极限。
\end{thm}

\begin{proof}
  注意，我们这里证明的是极限的存在性，而不仅仅是纯粹存在性。
  给定一个 Cauchy 逼近 $x : \Qp \to \RD$，定义
  %
  \begin{align*}
    L_y(q) &\defeq \exis{\epsilon, \theta : \Qp} L_{x_\epsilon}(q + \epsilon + \theta),\\
    U_y(q) &\defeq \exis{\epsilon, \theta : \Qp} U_{x_\epsilon}(q - \epsilon - \theta).
  \end{align*}
  %
  显然，$L_y$ 和 $U_y$ 都是有实例且舍入的。互不相交性可由 Cauchy 逼近条件得出。为证明定位性，考虑任意
  满足 $q < r$ 的 $q, r : \Q$。存在 $\epsilon : \Qp$ 使得
  $4 \epsilon < r - q$。由于 $q + 2 \epsilon < r - 2 \epsilon$，纯粹地
  $L_{x_\epsilon}(q + 2 \epsilon)$ 或者 $U_{x_\epsilon}(r - 2 \epsilon)$。在第一种情况下，
  我们有 $L_y(q)$，而在第二种情况下有 $U_y(r)$。

  为了证明 $y$ 是 $x$ 的极限，考虑任意 $\epsilon, \theta : \Qp$。因为
  $\Q$ 在 $\RD$ 中稠密，所以纯粹存在 $q, r : \Q$ 使得
  %
  \begin{narrowmultline*}
    x_\epsilon - \epsilon - \theta < q < x_\epsilon - \epsilon - \theta/2
    < x_\epsilon < \\
    x_\epsilon + \epsilon + \theta/2 < r < x_\epsilon + \epsilon + \theta,
  \end{narrowmultline*}
  %
  因此 $q < y < r$。由此可得 $|y - x_\epsilon| < \epsilon + \theta$。
\end{proof}

为了完备性，我们也记录下经典的表述。

\begin{cor}
  假设 $x : \N \to \RD$ 满足 Cauchy 条件~\eqref{eq:cauchy-sequence}。则
  存在 $y : \RD$ 使得
  %
  \begin{equation*}
    \prd{\epsilon : \Qp} \sm{n : \N} \prd{m \geq n} |x_m - y| < \epsilon.
  \end{equation*}
\end{cor}

\begin{proof}
  由 \cref{thm:ttac}，存在 $M : \Qp \to \N$ 使得 $\bar{x}(\epsilon) \defeq
  x_{M(\epsilon/2)}$ 是一个 Cauchy 逼近。令 $y$ 为其极限（该极限的存在性由
  \cref{RD-cauchy-complete} 保证）。给定任意的 $\epsilon : \Qp$，令 $n \defeq M(\epsilon/4)$
  并观察到，对于任意 $m \geq n$，
  %
  \begin{narrowmultline*}
    |x_m - y| \leq |x_m - x_n| + |x_n - y| =
    |x_m - x_n| + |\bar{x}(\epsilon/2) - y| < \narrowbreak
    \epsilon/4 + \epsilon/2 + \epsilon/4 = \epsilon.\qedhere
  \end{narrowmultline*}
\end{proof}

\subsection{Dedekind 实数是Dedekind 完备的}
\label{sec:RD-dedekind-complete}

我们通过取 $\Q$ 上Dedekind 分割的类型得到了 $\RD$。但我们本也可以从任一 Archimedean 有序域 $F$ 出发，构造 $F$ 上的\index{cut!Dedekind}Dedekind 分割。这些分割同样构成一个 Archimedean 有序域 $\bar{F}$，即\define{$F$ 的Dedekind 完备化}，%
\index{completion!Dedekind}%
\indexsee{Dedekind!completion}{completion, Dedekind}
其中 $F$ 作为子域被包含。如果将这一构造应用于 $\RD$ 本身，我们会得到更多的实数吗？答案是否定的。事实上，我们将证明一个更强的结论：$\RD$ 是终极的。

一个有序域~$F$ 称为\define{对 $\Omega$ 可容许的}，
\indexsee{admissible!ordered field}{ordered field, admissible}%
\indexdef{ordered field!admissible}%
如果 $F$ 上的严格序 $<$ 是一个映射 ${<} : F \to F \to \Omega$。

\begin{thm} \label{RD-final-field}
  每个对 $\Omega$ 可容许的 Archimedean 有序域都是 $\RD$ 的子域。
\end{thm}

\begin{proof}
  设 $F$ 为一个 Archimedean 有序域。对每个 $x : F$，由下式定义 $L_x, U_x : \Q \to
  \Omega$：
  %
  \begin{equation*}
    L_x(q) \defeq (q < x)
    \qquad\text{and}\qquad
    U_x(q) \defeq (x < q).
  \end{equation*}
  %
  （此处用到了 $F$ 对 $\Omega$ 可容许的假设。）
  则 $(L_x, U_x)$ 是一个\index{cut!Dedekind}Dedekind 分割。事实上，由于 $F$ 是 Archimedean 的且 $<$ 具有传递性，该分割是有实例的且为舍入的；由于 $<$ 不自反，它是不交的；由于 $<$ 是弱线性序，它是定位的。令 $e : F \to \RD$ 为映射 $e(x) \defeq (L_x,
  U_x)$。

  我们断言 $e$ 是一个保持且反映序的域嵌入。首先注意到，对于有理数 $q$，有 $e(q) = q$。其次，对所有 $x, y : F$，有如下等价：
  %
  \begin{narrowmultline*}
    x < y \Leftrightarrow
    (\exis{q : \Q} x < q < y) \Leftrightarrow \narrowbreak
    (\exis{q : \Q} U_x(q) \land L_y(q)) \Leftrightarrow
    e(x) < e(y),
  \end{narrowmultline*}
  %
  因此 $e$ 确实保持并反映序关系。$e(x + y) = e(x) + e(y)$ 成立是因为，对所有 $q : \Q$：
  %
  \begin{equation*}
    q < x + y \Leftrightarrow
    \exis{r, s : \Q} r < x \land s < y \land q = r + s.
  \end{equation*}
  %
  从右到左的蕴含是显然的。至于另一个方向，若 $q < x + y$，则仅仅存在 $r : \Q$ 使得 $q - y < r < x$；取 $s \defeq q - r$，即得所需的 $r$ 和 $s$。我们将 $e$ 保持乘法的验证留作练习。
\end{proof}

为了证明 $\RD$ 上的Dedekind 分割不会给出任何新东西，我们只需再有一个引理。

\begin{lem} \label{lem:cuts-preserve-admissibility}
  如果 $F$ 对 $\Omega$ 可容许，那么它的Dedekind 完备化也对 $\Omega$ 可容许。
  \index{completion!Dedekind}%
\end{lem}

\begin{proof}
  设 $\bar{F}$ 为 $F$ 的Dedekind 完备化。$\bar{F}$ 上的严格序定义为
  %
  \begin{equation*}
    ((L,U) < (L',U')) \defeq \exis{q : \Q} U(q) \land L'(q).
  \end{equation*}
  %
  由于 $U(q)$ 和 $L'(q)$ 都是 $\Omega$ 的元素，只要 $\Omega$ 对合取与可数存在量化封闭——而我们从一开始就做了这一假设——引理即成立。
\end{proof}

\begin{cor} \label{RD-dedekind-complete}
  %
  \indexdef{complete!ordered field, Dedekind}%
  \indexdef{Dedekind!completeness}%
  Dedekind 实数是Dedekind 完备的：对每个取实数值的Dedekind 分割 $(L, U)$，都存在唯一的 $x : \RD$，使得对所有 $y : \RD$，有 $L(y) = (y < x)$ 且 $U(y) = (x < y)$。
\end{cor}

\begin{proof}
  由 \cref{lem:cuts-preserve-admissibility}，$\RD$ 的Dedekind 完备化 $\barRD$ 对 $\Omega$ 可容许，故由 \cref{RD-final-field} 可得嵌入 $\barRD \to \RD$，以及嵌入 $\RD \to \barRD$。但这两个嵌入必为同构，因为它们的复合是固定稠密子域~$\Q$ 的保序域\index{homomorphism!field}同态，从而就是恒等映射。由 $\barRD \to \RD$ 为同构，推论立即得证。
\end{proof}

\index{real numbers!Dedekind|)}%

\section{Cauchy 实数}
\label{sec:cauchy-reals}

\index{real numbers!Cauchy|(}%
\index{completion!Cauchy|(}%
\indexsee{Cauchy!completion}{completion, Cauchy}%
Cauchy 实数按其意图是 \Q 在柯西序列极限下的完备化。\index{Cauchy!sequence}
在 Cauchy 实数的经典构造中，我们考虑 \Q 中所有柯西序列的集合 $\mathcal{C}$，然后形成适当的商 $\mathcal{C}/{\approx}$。
然后，为证明 $\mathcal{C}/{\approx}$ 是 Cauchy 完备的，我们考虑一个柯西序列 $x : \N \to \mathcal{C}/{\approx}$，将其提升为一个序列的序列 $\bar{x} : \N \to \mathcal{C}$，并用 $\bar{x}$ 构造 $x$ 的极限。然而，将~$x$ 提升为 $\bar{x}$ 需要用到可数选择公理（即~\eqref{eq:ac} 在 $X=\N$ 时的特例）或排中律，而这是我们可能希望避免的。
\indexdef{axiom!of choice!countable}%
每个以取商为最后一步的实数构造都有此缺陷。
在构造性数学中，摆脱这一困境有三种常见途径：
\index{mathematics!constructive}%
%
\index{bargaining}%

\begin{enumerate}
\item 假装实数是一个 setoid $(\mathcal{C}, {\approx})$，即柯西序列的类型 $\mathcal{C}$ 附带一个通过行政命令加上去的重合\index{coincidence, of Cauchy approximations}关系。如此一来，一个实数序列就\emph{只不过}是一个代表这些实数的柯西序列的序列。
\item 向诱惑低头，接受可数选择公理。毕竟，该公理在大多数基于计算观点的构造性数学模型（如可实现性模型）中都是成立的。
\item 宣布 Cauchy 实数不足取，转而构造Dedekind 实数。在某些语境下，例如在构造性数学的层论模型中，这样的判决完全成立。然而，正如我们在 \cref{sec:dedekind-reals} 所见，构造性的Dedekind 实数也有其自身的问题。
\end{enumerate}

然而，利用高阶归纳类型，我们还有第四种解决方案。我们相信它比以上任何一种都更可取，甚至对经典数学家来说也很有趣。
这个想法是：Cauchy 实数应当是由~\Q 生成的\emph{自由完备度量空间}\index{free!complete metric space}。
一般来说，构造任何种类的自由结构，都需要将该结构的操作反复应用于生成元。
例如，集合 $X$ 上自由群的元素，并不仅仅是 $X$ 中元素的二元乘积和逆，还包括通过迭代乘积和求逆构造而得到的``词''。
因此，我们自然可以预期 Cauchy 完备化也是如此，这里相关的\emph{操作}就是\emph{取一个柯西序列的极限}。
（在这种情况下，迭代必须超限地进行，因为即使在无穷多步之后，也仍然会有新的柯西序列需要取极限。）

前文提到的论证表明，如果排中律或可数选择公理成立，那么 Cauchy 完备化会变得非常特殊：在构造一个空间的完备化时，只需在应用该操作\emph{一步}之后即可停止。
这可以类比于自由幺半群和自由群能够用（约化）词来显式描述的事实。
然而，我们在 \cref{sec:free-algebras} 看到，高阶归纳类型允许我们\emph{直接}构造自由结构，无论是否还有可用的显式描述。
在本节中我们将说明，这对于 Cauchy 实数同样成立（类似的技术可用于构造任何度量空间的 Cauchy 完备化；参见 \cref{ex:metric-completion}）。
具体来说，高阶归纳类型允许我们\emph{同时}添加柯西序列的极限并对重合关系作商，从而避免了将实数序列提升为代表元序列的问题。
\index{completion!Cauchy|)}%

\subsection{Cauchy 实数的构造}
\label{sec:constr-cauchy-reals}

将 Cauchy 实数 $\RC$ 构造为高阶归纳类型，比 \cref{sec:free-algebras} 中考虑的自由代数结构的构造更为微妙。
我们打算引入一个``取极限''构造子，它以实数的柯西序列作为输入；但``实数的柯西序列''这一概念本身，又依赖于某种度量实数间``距离''的方法。
当然，一般来说，两个实数之间的距离本身也是一个实数，这可能带来有问题的循环。

然而，对于``实数的柯西序列''这一概念，我们实际需要的并非一般的``距离''概念，而是一种能够表达``两个实数之间的距离\index{distance}小于 $\epsilon$''的方法，对任意 $\epsilon:\Qp$ 成立。
这可以用一个二元关系族来表示，记作 $\mathord{\close\epsilon} : \RC\to\RC\to \prop$。
关系 $x \close\epsilon y$ 的预期含义是 $|x - y| < \epsilon$；但由于我们尚未定义减法、绝对值或不等式等概念（毕竟我们才刚刚在定义 $\RC$），我们就不得不在定义 $\RC$ 本身的同时，一并定义这些关系 $\close\epsilon$。
又因为 $\close\epsilon$ 是一个以两份 $\RC$ 为指标的类型族，所以我们无法用普通的相互（高阶）归纳定义来完成；取而代之，我们必须使用一种\emph{高阶归纳-归纳定义}。
\index{inductive-inductive type!higher}

回顾 \cref{sec:generalizations} 可知，普通意义上的归纳-归纳定义允许我们通过同时归纳的方式，定义一个类型以及一个以该类型为指标的类型族。
当然，它的``高阶''版本则允许该类型和该类型族都既包含点构造子，也包含道路构造子。
本书不会尝试建立高阶归纳-归纳定义的一般性理论，但我们希望接下来对 $\RC$ 和 $\close\epsilon$ 的描述，能够使这个想法清晰明了。

\begin{rmk}
  我们也可以考虑\emph{高阶归纳-递归定义（higher inductive-recursive definition）}，其中 $\close\epsilon$ 是利用 $\RC$ 的\emph{递归}原理来定义的，而这一定义又与 $\RC$ 本身的\emph{归纳}定义同时进行。
  但我们选择了归纳-归纳的路线，原因有二。
  首先，高阶归纳-递归定义似乎更难在同伦语义中给出证成。
  其次，也是更重要的一点，归纳-归纳的定义能给出更强的归纳原理，而发展 Cauchy 实数的基本理论就已经需要用到它。
\end{rmk}

最后，如同我们在讨论 Dedekind 实数的 Cauchy 完备性时（见 \cref{sec:RD-cauchy-complete}）所做的那样，我们将使用\emph{Cauchy 逼近（Cauchy approximations）}（\cref{defn:cauchy-approximation}）而不是 Cauchy 序列。
当然，现在的 Cauchy 逼近将是由 Cauchy 实数构成的，而不是 Dedekind 实数或有理数。

\begin{defn}\label{defn:cauchy-reals}
  令 $\RC$ 以及关系 $\closesym:\RC \times \RC \times \Qp \to \type$ 为以下的高阶归纳-归纳类型族。
  \define{Cauchy 实数（Cauchy reals）}\indexdef{real numbers!Cauchy}%\indexsee{Cauchy!real numbers}{real numbers, Cau\-chy}的类型 $\RC$ 由以下构造子生成：
  \begin{itemize}
  \item \emph{有理点（rational points）}：
    对任意 $q : \Q$，存在一个实数 $\rcrat(q)$。
    \index{rational numbers!as Cauchy real numbers}%
  \item \emph{极限点（limit points）}：
    对任意 $x : \Qp \to \RC$，若满足
    %
    \begin{equation}
      \label{eq:RC-cauchy}
      \fall{\delta, \epsilon : \Qp} x_\delta \close{\delta + \epsilon} x_\epsilon
    \end{equation}
    %
    则存在一点 $\rclim(x) : \RC$。我们称 $x$ 为一个\define{Cauchy 逼近（Cauchy approximation）}\indexdef{Cauchy!approximation}%
\index{limit!of a Cauchy approximation}。%
  \item \emph{道路（paths）}：
    对 $u, v : \RC$，若满足
    %
    \begin{equation}
      \label{eq:RC-path}
      \fall{\epsilon : \Qp} u \close\epsilon v
    \end{equation}
    %
    则存在一条道路 $\rceq(u, v) : \id[\RC]{u}{v}$。
  \end{itemize}
  同时，类型族 $\closesym:\RC\to\RC\to\Qp \to\type$ 由以下构造子生成。
  此处 $q$ 与 $r$ 表示有理数；$\delta$、$\epsilon$ 与 $\eta$ 表示正有理数；$u$ 与 $v$ 表示 Cauchy 实数；$x$ 与 $y$ 表示 Cauchy 逼近：
  \begin{itemize}
  \item 对任意 $q,r,\epsilon$，若 $-\epsilon < q - r < \epsilon$，则 $\rcrat(q) \close\epsilon \rcrat(r)$，
  \item 对任意 $q,y,\epsilon,\delta$，若 $\rcrat(q) \close{\epsilon - \delta} y_\delta$，则 $\rcrat(q) \close{\epsilon} \rclim(y)$，
  \item 对任意 $x,r,\epsilon,\delta$，若 $x_\delta \close{\epsilon - \delta} \rcrat(r)$，则 $\rclim(x) \close\epsilon \rcrat(r)$，
  \item 对任意 $x,y,\epsilon,\delta,\eta$，若 $x_\delta \close{\epsilon - \delta - \eta} y_\eta$，则 $\rclim(x) \close\epsilon \rclim(y)$，
  \item 对任意 $u,v,\epsilon$，若 $\xi,\zeta : u \close{\epsilon} v$，则 $\xi=\zeta$（命题截断）。
  \end{itemize}
\end{defn}

\mentalpause

$\RC$ 的第一个构造子说明任何有理数都可以看作一个实数。
第二个构造子说明，从任意 Cauchy 逼近出发，我们可以得到一个新的实数，称为它的“极限”。
第三个构造子则表达了这样一个想法：如果两个 Cauchy 逼近相一致，那么它们的极限相等。

$\closesym$ 的前四个构造子分别规定：两个有理数何时接近、一个有理数与一个极限何时接近，以及两个极限何时接近。
对于两个有理数的情况，这正好是通常意义下有理数的 $\epsilon$-接近；而其他情况则可以通过以下观察导出：每个逼近项 $x_\delta$ 都应当在其极限 $\rclim(x)$ 的 $\delta$ 范围之内。

我们提醒自己注意证明相关性：通过 $\rclim$ 得到的实数，其表示不仅包含一个 Cauchy 逼近 $x$，还包含 \eqref{eq:RC-cauchy} 的一个证明 $p$，所以严格来说我们应该写成 $\rclim(x,p)$ 而不是 $\rclim(x)$。
类似的观察也适用于 $\rceq$ 和 \eqref{eq:RC-path}，但我们将仅写作 $\rceq : u = v$ 而非 $\rceq(u, v, p) : u = v$。
这些记法滥用之所以可以接受，是因为我们省略了那些命题以及可以轻易推断出的信息。
同样地，$\mathord{\close\epsilon}$ 的最后一个构造子也使得我们可以不给其他四个构造子命名。

我们立刻就能用许多实数来填充 $\RC$。假设 $x : \N \to
\Q$ 是一个传统的有理数 Cauchy 序列\index{Cauchy!sequence}，并令 $M : \Qp \to \N$ 为其收敛模。
那么 $\rcrat \circ x \circ M : \Qp \to \RC$ 就是一个 Cauchy 逼近（利用 $\closesym$ 的第一个构造子来生成必要的见证）。
因此，$\rclim(\rcrat \circ x \circ M)$ 就是一个实数。
像 $\sqrt{2}$、$\pi$、$e$…… 这样的著名实数，都是这类有理数 Cauchy 序列的极限。

\subsection{Cauchy 实数上的归纳与递归}
\label{sec:induct-recurs-cauchy}

当然，要想对 $\RC$ 做些有用的事情，我们需要给出它的归纳原理。
正如同时归纳定义两个或更多对象时的情况一样，$\RC$ 和 $\closesym$ 的基本归纳原理要求对两者同时进行归纳。
因此，我们预期该原理是说：假设有两个类型族分别居于 $\RC$ 和 $\closesym$ 之上，并给定对应于每个构造子的数据，那么这两个类型族就都存在截面。
然而，由于 $\closesym$ 以 $\RC$ 的两个拷贝为指标，这些类型族的精确依赖关系就有些微妙了。
归纳原理将适用于任意一对类型族：
\begin{align*}
A&:\RC\to\type\\
B&:\prd{x,y:\RC} A(x) \to A(y) \to \prd{\epsilon:\Qp} (x\close\epsilon y) \to \type.
\end{align*}
$A$ 的类型是明显的，但 $B$ 的类型就需要稍加思索了。
既然 $B$ 必须依赖于 $\closesym$，而 $\closesym$ 本身又依赖于两个 $\RC$ 的拷贝和一个 $\Qp$ 的拷贝，那么很明显，$B$ 也必须依赖于变量 $x,y:\RC$、$\epsilon:\Qp$ 以及 $(x\close\epsilon y)$ 的一个元素。
稍不明显的是，$B$ 还必须依赖于 $A(x)$ 和 $A(y)$。

如果我们考虑非依值情形（即递归原理），这一点可能会更清楚；在那种情形下，$A$ 是一个简单的类型（而不是类型族）。
此时，我们会希望 $B$ 不依赖于 $x,y:\RC$ 或 $x\close\epsilon y$。
但递归原理（连同其相关的唯一性原理）意在表明，带有关联关系 $\close\epsilon$ 的 $\RC$ 是某个范畴中的``始对象''。
因此，在这种情况下，$A$ 和 $B$ 的依赖结构应当反映出 $\RC$ 和 $\close\epsilon$ 的依赖结构：也就是说，我们应有 $B:A\to A\to \Qp \to \type$。
将这一观察与依值情形下 $B$ 还必须依赖于 $x,y:\RC$ 和 $x\close\epsilon y$ 的事实结合起来，便不可避免地导出了上面给出的 $B$ 的类型。

\symlabel{RC-recursion}
将 $B$ 看作一个以 $\epsilon$ 为索引的、介于类型 $A(x)$ 与 $A(y)$ 之间的关系族，是很有帮助的。
据此，我们可以将 $B(x,y,a,b,\epsilon,\xi)$ 写作 $(x,a) \bsim_\epsilon^\xi (y,b)$。
由于 $\xi:x\close\epsilon y$ 在存在时是唯一的，我们通常将其从记号中省略，而写作 $(x,a) \bsim_\epsilon (y,b)$；只要我们记住这个关系仅在 $x\close\epsilon y$ 成立时有定义，这样做就不会出问题。
我们有时还会进一步简化，写作 $a\bsim_\epsilon b$，此时 $x$ 和 $y$ 可以从 $a$ 和 $b$ 的类型推断出来；但为了清晰起见，有时仍有必要将它们包含在记号中。

\index{induction principle!for Cauchy reals}%
现在，给定类型族 $A:\RC\to\type$ 以及如上所述的关系族 $\bsim$，归纳原理的前提由以下数据组成，每一条对应 $\RC$ 或 $\closesym$ 的一个构造子：

\begin{itemize}
\item 对于任意 $q : \Q$，有一个元素 $f_q:A(\rcrat(q))$。
\item 对于任意 Cauchy 逼近 $x$，以及任意满足 $a:\prd{\epsilon:\Qp} A(x_\epsilon)$ 的 \begin{equation}
    \fall{\delta, \epsilon : \Qp}
    (x_\delta,a_\delta) \bsim_{\delta+\epsilon} (x_\epsilon,a_\epsilon),
    \label{eq:depCauchyappx}
  \end{equation}，有一个元素 $f_{x,a}:A(\rclim(x))$。
  我们将这样的 $a$ 称为一个\define{依值 Cauchy 逼近（dependent Cauchy approximation）}
  \indexdef{Cauchy!approximation!dependent}%
  \indexsee{approximation, Cauchy}{Cauchy approximation}%
  \indexdef{dependent!Cauchy approximation}%
  在 $x$ 上。
\item 对于满足 $h:\fall{\epsilon : \Qp} u \close\epsilon v$ 的 $u, v : \RC$，以及所有满足 $\fall{\epsilon:\Qp} (u,a) \bsim_\epsilon (v,b)$ 的 $a:A(u)$ 与 $b:A(v)$，有一条依值道路 $\dpath{A}{\rceq(u,v)}{a}{b}$。
\item 对于 $q,r:\Q$ 和 $\epsilon:\Qp$，如果 $-\epsilon < q - r < \epsilon$，则有
  \narrowequation{(\rcrat(q),f_q) \bsim_\epsilon (\rcrat(r),f_r).}
\item 对于 $q:\Q$、$\delta,\epsilon:\Qp$、Cauchy 逼近 $y$，以及 $y$ 上的依值 Cauchy 逼近 $b$，如果 $\rcrat(q) \close{\epsilon - \delta} y_\delta$，那么
  \[(\rcrat(q),f_q) \bsim_{\epsilon-\delta} (y_\delta,b_\delta)
  \;\Rightarrow\;
  (\rcrat(q),f_q) \bsim_\epsilon (\rclim(y),f_{y,b}).\]
\item 类似地，对于 $r:\Q$、$\delta,\epsilon:\Qp$、Cauchy 逼近 $x$，以及 $x$ 上的依值 Cauchy 逼近 $a$，如果 $x_\delta \close{\epsilon - \delta} \rcrat(r)$，那么
  \[(x_\delta,a_\delta) \bsim_{\epsilon-\delta} (\rcrat(r),f_r)
  \;\Rightarrow\;
  (\rclim(x),f_{x,a}) \bsim_\epsilon (\rcrat(q),f_r).
  \]
\item 对于 $\epsilon,\delta,\eta:\Qp$、Cauchy 逼近 $x,y$，以及分别在 $x$ 和 $y$ 上的依值 Cauchy 逼近 $a$ 与 $b$，如果我们有 $x_\delta \close{\epsilon - \delta - \eta} y_\eta$，那么
  \[ (x_\delta,a_\delta) \bsim_{\epsilon - \delta - \eta} (y_\eta,b_\eta)
  \;\Rightarrow\;
  (\rclim(x),f_{x,a}) \bsim_\epsilon (\rclim(y),f_{y,b}).\]
\item 对于 $\epsilon:\Qp$、$x,y:\RC$、$\xi,\zeta:x\close{\epsilon} y$，以及 $a:A(x)$ 和 $b:A(y)$，$(x,a) \bsim_\epsilon^\xi (y,b)$ 与 $(x,a) \bsim_\epsilon^\zeta (y,b)$ 中的任意两个元素沿 $\xi=\zeta$ 依值相等。
  注意，按照惯例，这等价于要求 $\bsim$ 取值于命题。
\end{itemize}

在这些前提下，我们推导出函数
\begin{align*}
  f&:\prd{x:\RC} A(x)\\
  g&:\prd{x,y:\RC}{\epsilon:\Qp}{\xi:x\close{\epsilon} y}
  (x,f(x)) \bsim_\epsilon^\xi (y,f(y))
\end{align*}
它们按照预期计算：
\begin{align}
  f(\rcrat(q)) &\defeq f_q, \label{eq:rcsimind1}\\
  f(\rclim(x)) &\defeq f_{x,(f,g)[x]}. \label{eq:rcsimind2}
\end{align}
这里 $(f,g)[x]$ 表示将 $f$ 和 $g$ 应用于 Cauchy 逼近 $x$ 以得到 $x$ 上的依值 Cauchy 逼近。
也就是说，我们定义 $(f,g)[x]_\epsilon \defeq f(x_\epsilon) : A(x_\epsilon)$，然后对于任意 $\epsilon,\delta:\Qp$，我们有 $g(x_\epsilon,x_\delta,\epsilon+\delta,\xi)$ 来见证 $(f,g)[x]$ 是一个依值 Cauchy 逼近，其中 $\xi: x_\epsilon \close{\epsilon+\delta} x_\delta$ 来自 $x$ 是 Cauchy 逼近的假设。

我们之后再也不会用到这个记号了，所以不必费心去记它。
通常我们采用模式匹配的约定，其中 $f$ 通过形如~\eqref{eq:rcsimind1} 与~\eqref{eq:rcsimind2} 的等式来定义，而~\eqref{eq:rcsimind2} 的右侧可能出现符号 $f(x_\epsilon)$，以及它们构成依值 Cauchy 逼近这一假设。

不过，这条归纳原理的表述确实还是相当繁琐。
为了便于理解，我们注意到它作为特例包含了两个独立的归纳原理，分别关于~$\RC$ 和~$\closesym$。
首先，假设只给定一个类型族 $A:\RC\to\type$，并将 $\bsim$ 定义为取常值 \unit。
这样一来，许多所需的数据就变得平凡，余下仅需提供：

\begin{itemize}
\item 对任意 $q : \Q$，有一个元素 $f_q:A(\rcrat(q))$，
\item 对任意 Cauchy 逼近 $x$，以及任意 $a:\prd{\epsilon:\Qp} A(x_\epsilon)$，有一个元素 $f_{x,a}:A(\rclim(x))$，
\item 对 $u, v : \RC$ 和 $h:\fall{\epsilon : \Qp} u \close\epsilon v$，以及 $a:A(u)$ 与 $b:A(v)$，有 $\dpath{A}{\rceq(u,v)}{a}{b}$。
\end{itemize}

给定这些数据，归纳原理就给出一个函数 $f:\prd{x:\RC} A(x)$，使得
\begin{align*}
  f(\rcrat(q)) &\defeq f_q,\\
  f(\rclim(x)) &\defeq f_{x,f(x)}.
\end{align*}
我们称此原理为 \define{$\RC$-归纳}；它本质上说的是，若将 $\close\epsilon$ 视为给定的，则 $\RC$ 就是由其构造子归纳生成的。

请注意，若 $A$ 是命题，则 $\RC$-归纳中的第三个假设自动成立（稍后我们会看到，这二者实际上是等价的）。
因此，要证明关于实数的命题，只需对有理数以及 Cauchy 逼近的极限加以证明即可。
下面是一个例子。

\begin{lem} \label{lem:close-reflexive}
  对任意 $u:\RC$ 与 $\epsilon:\Qp$，有 $u\close\epsilon u$。
\end{lem}

\begin{proof}
  定义 $A(u) \defeq \fall{\epsilon:\Qp} (u\close\epsilon u)$。
  由于这是一个命题（根据 $\closesym$ 的最后一个构造子），由 $\RC$-归纳，只需对 $u$ 为 $\rcrat(q)$ 和 $u$ 为 $\rclim(x)$ 两种情形加以证明。
  在第一种情形，显然对任意 $\epsilon$ 都有 $|q-q|<\epsilon$，于是根据 $\closesym$ 的第一个构造子，有 $\rcrat(q) \close\epsilon \rcrat(q)$。
  %
  在第二种情形，我们可以归纳地假设对所有 $\delta,\epsilon:\Qp$ 都有 $x_\delta \close\epsilon x_\delta$。
  那么特别地，我们有 $x_{\epsilon/3} \close{\epsilon/3} x_{\epsilon/3}$，从而根据 $\closesym$ 的第四个构造子，得到 $\rclim(x) \close{\epsilon} \rclim(x)$。
\end{proof}

由 \cref{lem:close-reflexive} 可知，只有当族 $A:\RC\to\type$ 是命题时，直接应用 $\RC$-归纳才有可能成功。
为了看清这一点，固定 $u:\RC$。
取 $v$ 为 $u$，则 $\RC$-归纳的第三个假设告诉我们，对任意 $a : A(u)$，有 $\dpath{A}{\rceq(u,u)}{a}{a}$。
此外，若再给定一个点 $b : A(u)$，我们还可得到 $\dpath{A}{\rceq(u,u)}{a}{b}$。
根据依值道路类型的定义，我们得出结论：这两条道路合起来意味着 $a = b$，即 $A(u)$ 中所有的点都相等。

\begin{thm}\label{thm:Cauchy-reals-are-a-set}
  $\RC$ 是一个集合。
\end{thm}

\begin{proof}
  我们刚刚证明了，纯粹关系
  \narrowequation{P(u,v) \defeq \fall{\epsilon:\Qp} (u\close\epsilon v)}
  是自反的。
  由于它蕴含相等关系（根据 $\RC$ 的道路构造子），结论由 \cref{thm:h-set-refrel-in-paths-sets} 可得。
\end{proof}

我们还可以证明，尽管 $\RC$ 可能不是\emph{有理数} Cauchy 序列集合的商集，但它仍然是\emph{实数} Cauchy 序列集合的商集。
（当然，这不能作为 $\RC$ 的有效\emph{定义}，但它是一个有用的性质。）
我们将 Cauchy 逼近的类型定义为
%
\symlabel{cauchy-approximations}%
\index{Cauchy!approximation!type of}%
\begin{equation*}
  \CAP \defeq
  \setof{ x : \Qp \to \RC |
    \fall{\epsilon, \delta : \Qp} x_\delta \close{\delta + \epsilon} x_\epsilon
  }.
\end{equation*}
$\RC$ 的第二个构造子给出了一个函数 $\rclim:\CAP\to\RC$。

\begin{lem} \label{RC-lim-onto}
  每个实数都仅仅是极限点：$\fall{u : \RC} \exis{x : \CAP} u = \rclim(x)$。
  换言之，$\rclim:\CAP\to\RC$ 是满的。
\end{lem}

\begin{proof}
  根据 $\RC$-归纳，我们可以对 $u$ 分情况讨论。
  显然，如果 $u$ 是一个极限 $\rclim(x)$，结论是平凡的。
  因此假设 $u$ 是一个有理点 $\rcrat(q)$；我们断言 $u$ 等于 $\rclim(\lam{\epsilon} \rcrat(q))$。
  根据 $\RC$ 的道路构造子，只需证明对所有 $\epsilon:\Qp$ 有 $\rcrat(q) \close\epsilon \rclim(\lam{\epsilon} \rcrat(q))$。
  而根据 $\closesym$ 的第二个构造子，为此只需找到 $\delta:\Qp$ 使得 $\rcrat(q)\close{\epsilon-\delta} \rcrat(q)$。
  但根据 $\closesym$ 的第一个构造子，我们可以取任意满足 $\delta<\epsilon$ 的 $\delta:\Qp$。
\end{proof}

%

\begin{lem} \label{RC-lim-factor}
  如果 $A$ 是一个集合，且函数 $f : \CAP \to A$ 保持 Cauchy 逼近的\index{coincidence!of Cauchy approximations}重合，意指
  %
  \begin{equation*}
    \fall{x, y : \CAP} \rclim(x) = \rclim(y) \Rightarrow f(x) = f(y),
  \end{equation*}
  %
  则 $f$ 沿 $\rclim : \CAP \to \RC$ 唯一分解。
\end{lem}

\begin{proof}
  由于 $\rclim$ 是满的，根据 \cref{lem:images_are_coequalizers}，$\RC$ 是 $\CAP$ 关于 $\rclim$ 的\index{kernel!pair}核对的商集。
  而这正是本引理的断言。
\end{proof}

现在来看归纳原理的第二个特殊情况。假设我们将 $A$ 取为取值恒为 $\unit$ 的常值类型族。
在这种情况下，$\bsim$ 就只是一个以 $\epsilon$ 为指标、定义在 $\epsilon$-接近的实数对上的关系族，因此我们可以记 $u\bsim_\epsilon v$ 来代替 $(u,\ttt)\bsim_\epsilon (v,\ttt)$。
此时所需的数据简化为如下几项，其中 $q, r$ 表示有理数，$\epsilon, \delta, \eta$ 表示正有理数，$x, y$ 表示 Cauchy 逼近：

\begin{itemize}
\item 若 $-\epsilon < q - r < \epsilon$，则
  $\rcrat(q) \bsim_\epsilon \rcrat(r)$，
\item 若 $\rcrat(q) \close{\epsilon - \delta} y_\delta$ 且
  $\rcrat(q)\bsim_{\epsilon-\delta} y_\delta$，
  则 $\rcrat(q) \bsim_\epsilon \rclim(y)$，
\item 若 $x_\delta \close{\epsilon - \delta} \rcrat(r)$ 且
  $x_\delta \bsim_{\epsilon-\delta} \rcrat(r)$，
  则 $\rclim(y) \bsim_\epsilon \rcrat(q)$，
\item 若 $x_\delta \close{\epsilon - \delta - \eta} y_\eta$ 且
  $x_\delta\bsim_{\epsilon - \delta - \eta} y_\eta$，
  则 $\rclim(x) \bsim_\epsilon \rclim(y)$。
\end{itemize}

最终得到的结论是 $\fall{u,v:\RC}{\epsilon:\Qp} (u\close\epsilon v) \to (u \bsim_\epsilon v)$。
我们称这一原理为 \define{$\closesym$-归纳（$\closesym$-induction）}；它本质上说的是，如果我们把 $\RC$ 视为给定的，那么 $\close\epsilon$ 作为类型族，是由\emph{它自己的}构造子归纳生成的。
例如，我们可以用它来证明 $\closesym$ 是对称的。

\begin{lem}\label{thm:RCsim-symmetric}
  对于任意 $u,v:\RC$ 和 $\epsilon:\Qp$，有 $(u\close\epsilon v) = (v\close\epsilon u)$。
\end{lem}

\begin{proof}
  由于两边都是命题，利用对称性只需证明其中一个蕴含关系即可。
  因此，设 $(u\bsim_\epsilon v) \defeq (v\close\epsilon u)$。
  根据 $\closesym$-归纳，我们可以将问题归约到下述情形：$u\close\epsilon v$ 是由 $\closesym$ 的四个有趣的构造子之一推导出来的。
  在第一种情形，即 $u$ 和 $v$ 均为有理数时，结果是显然的（我们可以再次应用第一个构造子）。
  在另外三种情形中，归纳假设（结合 $\Q$ 中加法的交换性）恰好给出了 $\closesym$ 的另一个构造子所需的输入（第二和第三个构造子互换角色，而第四个则保持不变）。
\end{proof}

因此，一般的归纳原理——我们可称之为 \define{$(\RC,\closesym)$-归纳（$(\RC,\closesym)$-induction）}——是 $\RC$-归纳与 $\closesym$-归纳的一种联合形式。
例如，考虑它的非依值版本，我们称之为 \define{$(\RC,\closesym)$-递归（$(\RC,\closesym)$-recursion）}，这也是我们最常用到的形式。
\index{recursion principle!for Cauchy reals}%
普通的 $\RC$-递归告诉我们，要定义一个函数 $f : \RC \to A$，只需满足：

\begin{enumerate}
\item 对每个 $q : \Q$，构造 $f(\rcrat(q)) : A$，
\item 对每个 Cauchy 逼近 $x : \Qp \to \RC$，构造 $f(x) : A$，
  这里假定对所有的 $\epsilon : \Qp$，$f(x_\epsilon)$ 已被定义，
\item 对所有满足 $\fall{\epsilon:\Qp} u\close\epsilon v$ 的 $u, v : \RC$，证明 $f(u) = f(v)$。\label{item:rcrec3}
\end{enumerate}

然而，在不知道 $f$ 如何作用在 $\epsilon$-接近的 Cauchy 实数上的情况下，通常很难证明~\ref{item:rcrec3}。
增强的 $(\RC,\closesym)$-递归原理弥补了这一缺陷，它允许我们指定一个\emph{任意的} ``$f$ 在 $\epsilon$-接近的 Cauchy 实数上的作用方式''，然后我们可以通过与 $f$ 的定义同步进行的归纳来证明这一点确实成立。
这个``作用方式''就是关系族 $\bsim$。
由于 $A$ 独立于 $\RC$，为了简化，我们可以假设 $\bsim$ 只依赖于 $A$ 和 $\Qp$，因此用 $a\bsim_\epsilon b$ 代替 $(u,a) \bsim_\epsilon (v,b)$ 不会产生歧义。
在这种情况下，使用 $(\RC,\closesym)$-递归定义一个函数 $f:\RC\to A$ 需要处理以下几种情形（现在我们使用模式匹配约定的写法来书写）。

\begin{itemize}
\item 对每个 $q : \Q$，构造 $f(\rcrat(q)) : A$。
\item 对每个 Cauchy 逼近 $x : \Qp \to \RC$，构造 $f(\rclim(x)) : A$，这里归纳地假设对所有 $\epsilon : \Qp$，$f(x_\epsilon)$ 已经定义，并且构成一个``关于 $\bsim$ 的 Cauchy 逼近''，即 $\fall{\epsilon,\delta:\Qp} (f(x_\epsilon) \bsim_{\epsilon+\delta} f(x_\delta))$。
\item 证明关系 $\bsim$ 是\emph{分离的}，即对任意 $a,b:A$，
  \indexdef{relation!separated family of}%
  \indexdef{separated family of relations}%
\narrowequation{(\fall{\epsilon:\Qp} a\bsim_\epsilon b) \Rightarrow (a=b).}
\item 证明对 $q,r:\Q$，若 $-\epsilon< q-r <\epsilon$，则 $f(\rcrat(q))\bsim_\epsilon f(\rcrat(r))$。
\item 对任意 $q:\Q$ 与任意 Cauchy 逼近 $y$，证明
\narrowequation{f(\rcrat(q)) \bsim_\epsilon f(\rclim(y)),} 这里归纳地假设 $\rcrat(q)\close{\epsilon-\delta} y_\delta$ 和 $f(\rcrat(q)) \bsim_{\epsilon-\delta} f(y_\delta)$ 对某个 $\delta:\Qp$ 成立，并且 $\eta \mapsto f(x_\eta)$ 是关于 $\bsim$ 的 Cauchy 逼近。
\item 对任意 Cauchy 逼近 $x$ 与任意 $r:\Q$，证明
\narrowequation{f(\rclim(x)) \bsim_\epsilon f(\rcrat(r)),}
这里归纳地假设 $x_\delta \close{\epsilon-\delta} \rcrat(r)$ 和 $f(x_\delta) \bsim_{\epsilon-\delta} f(\rcrat(r))$ 对某个 $\delta:\Qp$ 成立，并且 $\eta\mapsto f(x_\eta)$ 是关于 $\bsim$ 的 Cauchy 逼近。
\item 对任意 Cauchy 逼近 $x,y$，证明
\narrowequation{f(\rclim(x)) \bsim_\epsilon f(\rclim(y)),}
这里归纳地假设 $x_\delta \close{\epsilon-\delta-\eta} y_\eta$ 和 $f(x_\delta) \bsim_{\epsilon-\delta-\eta} f(y_\eta)$ 对某个 $\delta,\eta:\Qp$ 成立，并且 $\theta\mapsto f(x_\theta)$ 与 $\theta\mapsto f(y_\theta)$ 都是关于 $\bsim$ 的 Cauchy 逼近。
\end{itemize}


注意，在后四个证明中，我们可以自由使用前两条数据中给出的 $f(\rcrat(q))$ 和 $f(\rclim(x))$ 的具体定义。
然而，分离性的证明必须适用于 $A$ 的\emph{任意}两个元素，而与 $f$ 无关：它是关系族 $\bsim$ 的一种``可接受性''条件。
因此，我们通常在定义 $\bsim$ 之后立即验证它，然后再去定义 $f(\rcrat(q))$ 和 $f(\rclim(x))$。

在上述假设下，$(\RC,\closesym)$-递归给出一个函数 $f:\RC\to A$，使得 $f(\rcrat(q))$ 和 $f(\rclim(x))$ 与前两条中给出的定义判断相等。
此外，我们还可以得出
\begin{equation}
  \fall{u,v:\RC}{\epsilon:\Qp} (u\close\epsilon v) \to (f(u) \bsim_\epsilon f(v)).\label{eq:RC-sim-recursion-extra}
\end{equation}

作为一个典型例子，$(\RC,\closesym)$-递归允许我们将定义在 $\Q$ 上的函数延拓到整个 $\RC$，只要这些函数是充分连续的。
\index{function!continuous}%

\begin{defn}\label{defn:lipschitz}
  函数 $f:\Q\to\RC$ 称为 \define{Lipschitz 的}
  \indexdef{function!Lipschitz}%
  \indexdef{Lipschitz!function}%
  \indexdef{Lipschitz!constant}%
  \indexdef{constant!Lipschitz}%
  ，如果存在 $L:\Qp$（称为 \define{Lipschitz 常数}）使得
  \[ |q - r|<\epsilon \Rightarrow (f(q) \close{L\epsilon} f(r)) \]
  对所有 $\epsilon:\Qp$ 和 $q,r:\Q$ 成立。
  %
  类似地，$g:\RC\to\RC$ 是 \define{Lipschitz 的}，如果存在 $L:\Qp$ 使得
  \[ (u\close\epsilon v) \Rightarrow (g(u) \close{L\epsilon} g(v)) \]
  对所有 $\epsilon:\Qp$ 和 $u,v:\RC$ 成立。
\end{defn}

特别地，注意根据 $\closesym$ 的第一个构造子，如果 $f:\Q\to\Q$ 在显然的意义上是 Lipschitz 的，那么复合 $\Q\xrightarrow{f} \Q \to \RC$ 也是 Lipschitz 的。

\begin{lem}\label{RC-extend-Q-Lipschitz}
  设 $f : \Q \to \RC$ 是 Lipschitz 的，常数为 $L : \Qp$。
  则存在一个 Lipschitz 映射 $\bar{f} : \RC \to \RC$，同样以 $L$ 为常数，使得 $\bar{f}(\rcrat(q)) \jdeq f(q)$ 对所有 $q:\Q$ 成立。
\end{lem}

\begin{proof}
  % Uniqueness follows directly from \cref{RC-continuous-eq}.
  我们通过 $(\RC,\closesym)$-递归来定义 $\bar{f}$，陪域为 $A\defeq \RC$。
  我们将关系 $\mathord{\bsim}: \RC \to \RC \to \Qp \to \prop$ 定义为
  \begin{align*}
    (u \bsim_\epsilon v) &\defeq (u \close{L\epsilon} v).
  \end{align*}
  对于 $q : \Q$，我们定义
  %
  \begin{equation*}
    \bar{f}(\rcrat(q)) \defeq \rcrat(f(q)).
  \end{equation*}
  %
  对于一个 Cauchy 逼近 $x : \Qp \to \RC$，我们定义
  %
  \begin{equation*}
    \bar{f}(\rclim(x)) \defeq \rclim(\lamu{\epsilon : \Qp} \bar{f}(x_{\epsilon/L})).
  \end{equation*}
  %
  为使此定义有意义，我们必须验证 $y \defeq \lamu{\epsilon : \Qp} \bar{f}(x_{\epsilon/L})$ 确实是一个 Cauchy 逼近。
  而这步的归纳假设是：对任意 $\delta,\epsilon:\Qp$ 都有 $\bar{f}(x_\delta) \bsim_{\delta+\epsilon} \bar{f}(x_\epsilon)$，即 $\bar{f}(x_\delta) \close{L\delta+L\epsilon} \bar{f}(x_\epsilon)$。
  由此我们得到
  \[y_\delta \jdeq f(x_{\delta/L}) \close{\delta + \epsilon} f(x_{\epsilon/L})   \jdeq y_\epsilon. \]

  为证明分离性，我们只须观察到：$\fall{\epsilon:\Qp} a\bsim_\epsilon b$ 意味着 $\fall{\epsilon:\Qp} a\close{L\epsilon} b$，这蕴含着 $\fall{\epsilon:\Qp}a\close\epsilon b$，从而 $a=b$。

  为了完成 $(\RC,\closesym)$-递归，剩下还需验证 $\bsim$ 所满足的四个条件。
  这基本相当于证明 $\bar f$ 对于 $\closesym$ 的所有四个构造子都是 Lipschitz 的。
  \begin{enumerate}
  \item 当 $u$ 是 $\rcrat(q)$ 而 $v$ 是 $\rcrat(r)$，且满足 $-\epsilon < |q-r| <\epsilon$ 时，由 $f$ 是 Lipschitz 的假设可得 $f(q) \close{L\epsilon} f(r)$，从而根据定义有 $\bar{f}(\rcrat(q)) \bsim_\epsilon \bar{f}(\rcrat(r))$。
  \item 当 $u$ 是 $\rclim(x)$ 而 $v$ 是 $\rcrat(q)$，且满足 $x_\eta \close{\epsilon - \eta} \rcrat(q)$ 时，归纳假设是 $\bar{f}(x_\eta) \close{L \epsilon - L \eta} \rcrat(f(q))$；这结合 $\closesym$ 的第三个构造子，即证明了
      \narrowequation{\bar{f}(\rclim(x)) \close{L \epsilon} \bar{f}(\rcrat(q)).}
  \item 当 $u$ 为有理数而 $v$ 为极限的对称情形，证明本质上相同。
  \item 当 $u$ 是 $\rclim(x)$ 而 $v$ 是 $\rclim(y)$，并且存在 $\delta, \eta : \Qp$ 使得 $x_\delta \close{\epsilon - \delta - \eta} y_\eta$ 时，归纳假设是 $\bar{f}(x_\delta) \close{L \epsilon - L \delta - L \eta} \bar{f}(y_\eta)$；这结合 $\closesym$ 的第四个构造子，即证明了 $\bar{f}(\rclim(x)) \close{L
        \epsilon} \bar{f}(\rclim(y))$。
  \end{enumerate}
  至此完成了 $(\RC,\closesym)$-递归，从而完成了 $\bar f$ 的构造。
  所期望的等式 $\bar f(\rcrat(q))\jdeq f(q)$ 正是 $(\RC,\closesym)$-递归的第一个计算规则，而附加条件~\eqref{eq:RC-sim-recursion-extra} 则正好说明 $\bar f$ 是常量为 $L$ 的 Lipschitz 函数。
\end{proof}

在对 $\closesym$ 有更好的刻画之前，我们目前能做的也就只有这些了。
我们在 $\closesym$ 的构造子中规定了两种不同形式的 Cauchy 实数在何种条件下 $\epsilon$-接近。
但是，我们如何知道，在最终得到的归纳-归纳类型族中，这些就是该事实的\emph{唯一}见证呢？
我们已经看到，归纳类型族（例如相等类型，参见 \cref{sec:identity-systems}）和高阶归纳类型往往包含“比最初放入的还要多”的东西，因此这并非一个无关紧要的问题。

为了更精确地刻画 $\closesym$，我们将在 $\RC$ 上\emph{递归地}定义一个关系族 $\approx_\epsilon$，使其能在构造子上进行计算，然后证明这个关系族等价于 $\close\epsilon$。

\begin{thm}\label{defn:RC-approx}
  存在一族纯粹关系 $\mathord\approx:\RC\to\RC\to\Qp\to\prop$，使得
  \begin{align}
    (\rcrat(q) \approx_\epsilon \rcrat(r))  &\defeq
    (-\epsilon < q - r < \epsilon)\label{eq:RCappx1}\\
    (\rcrat(q) \approx_\epsilon \rclim(y)) &\defeq
    \exis{\delta : \Qp} \rcrat(q) \approx_{\epsilon - \delta} y_\delta\label{eq:RCappx2}\\
    (\rclim(x) \approx_\epsilon \rcrat(r)) &\defeq
    \exis{\delta : \Qp} x_\delta \approx_{\epsilon - \delta} \rcrat(r)\label{eq:RCappx3}\\
    (\rclim(x) \approx_\epsilon \rclim(y)) &\defeq
    \exis{\delta, \eta : \Qp} x_\delta \approx_{\epsilon - \delta - \eta} y_\eta.\label{eq:RCappx4}
  \end{align}
  此外，我们还有
  \begin{gather}
    (u \approx_\epsilon v) \Leftrightarrow \exis{\theta:\Qp} (u \approx_{\epsilon-\theta} v) \label{RC-sim-rounded}\\
    (u \approx_\epsilon v) \to (v\close\delta w) \to (u\approx_{\epsilon+\delta} w)\label{eq:RC-sim-rtri}\\
    (u \close\epsilon v) \to (v\approx_\delta w) \to (u\approx_{\epsilon+\delta} w)\label{eq:RC-sim-ltri}.
  \end{gather}
\end{thm}

附加条件~\eqref{RC-sim-rounded}--\eqref{eq:RC-sim-ltri} 是使归纳定义得以成立所必需的。
条件~\eqref{RC-sim-rounded} 称为\define{舍入的（rounded）}。
\indexsee{relation!rounded}{rounded relation}%
\indexdef{rounded!relation}%
从右向左读，便得出 $\approx$ 的\define{单调性（monotonicity）}，
\index{monotonicity}%
\index{relation!monotonic}%
%
\begin{equation*}
  (\delta < \epsilon) \land (u \approx_\delta v) \Rightarrow (u \approx_\epsilon v)
\end{equation*}
%
而从左向右读，则给出 $\approx$ 的\define{开性（openness）}，
\index{open!relation}%
\index{relation!open}%
%
\begin{equation*}
  (u \approx_\epsilon v) \Rightarrow \exis{\delta : \Qp} (\delta < \epsilon) \land (u \approx_\delta v).
\end{equation*}
%
条件~\eqref{eq:RC-sim-rtri} 与~\eqref{eq:RC-sim-ltri} 是三角不等式的两种表现形式，它们表明 $\approx$ 是 $\closesym$ 两侧的``模''。

\begin{proof}
  我们将通过双重 $(\RC,\closesym)$-递归来定义 $\mathord\approx:\RC\to\RC\to\Qp\to\prop$。
  首先，我们将目标域为 $\RC\to\Qp\to\prop$ 的某个子集进行 $(\RC,\closesym)$-递归，该子集由满足舍入性及某一种适当形式的三角形不等式的谓词族构成。
  若将这些谓词视作某个二元关系的一半，我们便将其记作 $(u,\epsilon) \mapsto (\hapx_\epsilon u)$，并用符号 $\hapname$ 指代整个关系。
  现在我们可以将 $A$ 精确地写为
  \begin{multline*}
    A \defeq\; \Bigg\{ \hapname :\RC\to\Qp\to\prop \;\bigg|\; \\
    \Big(\fall{u:\RC}{\epsilon:\Qp}
    \big((\hapx_\epsilon u) \Leftrightarrow \exis{\theta:\Qp} (\hapx_{\epsilon-\theta} u)\big)\Big)  \\
    \land \Big(\fall{u,v:\RC}{\eta,\epsilon:\Qp} (u\close\epsilon v) \to\\
    \big((\hapx_\eta u) \to (\hapx_{\eta+\epsilon} v) \big) \land \big((\hapx_\eta v) \to (\hapx_{\eta+\epsilon} u) \big)\Big)\Bigg\}
  \end{multline*}
  与子集的惯常记法一致，我们对 $A$ 的一个居民与其第一个分量 $\hapname$ 使用相同的记号。
  作为 $(\RC,\closesym)$-递归所需的关系族，我们考虑以下关系，它将确保另一种形式的三角形不等式：
  \begin{narrowmultline*}
    (\hapname \bsim_\epsilon \hapbname ) \defeq \narrowbreak
    \fall{u:\RC}{\eta:\Qp} ((\hapx_\eta u) \to (\hapxb_{\epsilon+\eta} u))
    \land \narrowbreak
    ((\hapxb_\eta u) \to (\hapx_{\epsilon+\eta} u)).
  \end{narrowmultline*}
  我们观察到这些关系是分离的。
  因为假设
  \narrowequation{\fall{\epsilon:\Qp} (\hapname \bsim_\epsilon \hapbname),}
  要证明 $\hapname = \hapbname$，只需对所有 $u:\RC$ 证明 $(\hapx_\epsilon u) \Leftrightarrow (\hapxb_\epsilon u)$。
  但由舍入性，$\hapx_\epsilon u$ 蕴含存在某个 $\theta$ 使得 $\hapx_{\epsilon-\theta} u$，这结合 $\hapname \bsim_\epsilon \hapbname$ 便蕴含 $\hapxb_\epsilon u$；反之亦然。

  现在递归原理所需的前两个数据如下。
  \begin{itemize}
  \item 对于任意 $q:\Q$，我们必须给出 $A$ 的一个元素，记作 $(\rcrat(q)\approx_{(\blank)} \blank)$。
  \item 对于任意 Cauchy 逼近 $x$，若假设已定义了一个函数 $\Qp \to A$（记作 $\epsilon \mapsto (x_\epsilon \approx_{(\blank)} \blank)$）满足性质
    % \[ \fall{u,v:\RC}{\delta,\epsilon,\eta:\Qp} (x_\delta \approx_\eta u) \to (u\close{\delta+\epsilon} v) \to (x_\epsilon \approx_{\eta+\delta+\epsilon} v) \]
    \begin{equation}
      \fall{u:\RC}{\delta,\epsilon,\eta:\Qp} (x_\delta \approx_\eta u) \to (x_\epsilon \approx_{\eta+\delta+\epsilon} u),\label{eq:appxrec2}
    \end{equation}
    我们就必须给出 $A$ 的一个元素，记作 $(\rclim(x)\approx_{(\blank)} \blank)$。
  \end{itemize}
  在这两种情形下，我们都通过嵌套的 $(\RC,\closesym)$-递归来给出所需的定义，此递归的目标域是 $\Qp\to\prop$ 中由舍入的命题族构成的子集。
  若将这些命题视作某个二元关系的“零半边”，我们便将其记作 $\epsilon \mapsto (\tap{\epsilon})$，并用符号 $\tapname$ 指代整个族。
  现在我们可以将这些内层递归的目标域精确地写为
  \begin{narrowmultline*}
    C \defeq
    \bigg\{ \tapname :\Qp\to\prop \;\;\Big|\;\; \narrowbreak
    \fall{\epsilon:\Qp} \Big((\tap\epsilon) \Leftrightarrow \exis{\theta:\Qp} (\tap{\epsilon-\theta})\Big)\bigg\}
  \end{narrowmultline*}
  我们取所需的关系族为三角形不等式的剩余部分：
  \begin{narrowmultline*}
    (\tapname \bbsim_\epsilon \tapbname) \defeq
    \fall{\eta:\Qp} ((\tap\eta) \to (\tapb{\epsilon+\eta})) \land
    \narrowbreak
    ((\tapb\eta) \to (\tap{\epsilon+\eta})).
  \end{narrowmultline*}
  利用 $C$ 中所有元素的舍入性，通过与 $\bsim$ 相同的论证可知这些关系是分离的。

  注意，若这样一个内层递归成功，它将产生一个谓词族 $\hapname : \RC\to\Qp\to \prop$，这些谓词是舍入的（因为它们在 $\Qp\to\prop$ 中的像落在 $C$ 中）并且满足
  \[ \fall{u,v:\RC}{\epsilon:\Qp} (u\close\epsilon v) \to \big((\hapx_{(\blank)} u) \bbsim_\epsilon (\hapx_{(\blank)} u)\big). \]
  展开 $\bbsim$ 的定义，这恰是 $\hapname$ 属于 $A$ 所需的第三个条件；因此这正是我们需要的。

  此时，我们可以给出定义~\eqref{eq:RCappx1}--\eqref{eq:RCappx4}，作为两个内层递归各自对应于有理点和极限的前两个子句。
  在每种情形下，我们都必须验证该关系是舍入的，从而位于 $C$ 中。
  在有理数-有理数情形~\eqref{eq:RCappx1} 中，这是显然的；而在其他情形中，则由归纳假设得出。（在~\eqref{eq:RCappx2} 中，相关的归纳假设是 $(\rcrat(q) \approx_{(\blank)} y_\delta) : C$，而在~\eqref{eq:RCappx3} 和~\eqref{eq:RCappx4} 中则是 $(x_\delta \approx_{(\blank)} \blank) : A$。）

  子递归剩余的数据在于证明 \eqref{eq:RCappx1}--\eqref{eq:RCappx4} 对 $\closesym$ 的构造子满足右侧的三角形不等式。
  共有八种情形——每个子递归四种——对应于~\eqref{eq:RC-sim-rtri} 中 $u$、$v$ 和 $w$ 选为有理点或极限的八种可能组合。
  首先我们考虑 $u$ 是 $\rcrat(q)$ 的情形。
  \begin{enumerate}
  \item 假设 $\rcrat(q)\approx_\phi \rcrat(r)$ 且 $-\epsilon<|r-s|<\epsilon$，我们必须证明 $\rcrat(q)\approx_{\phi+\epsilon} \rcrat(s)$。
    但由 $\approx$ 的定义，这归约为有理数的三角形不等式。
  \item 我们假设 $\phi,\epsilon,\delta:\Qp$ 满足 $\rcrat(q)\approx_\phi \rcrat(r)$ 与 $\rcrat(r) \close{\epsilon-\delta} y_\delta$，并归纳地假设
    \begin{equation}
      \fall{\psi:\Qp}(\rcrat(q) \approx_{\psi} \rcrat(r)) \to (\rcrat(q) \approx_{\psi+\epsilon-\delta} y_\delta).\label{eq:RCappx-rtri-rrl1}
    \end{equation}
    我们还假设 $\psi,\delta\mapsto (\rcrat(q) \approx_{\psi} y_\delta)$ 是一个关于 $\bbsim$ 的 Cauchy 逼近，即
    \begin{equation}
      \fall{\psi,\xi,\zeta:\Qp} (\rcrat(q) \approx_{\psi} y_\xi) \to (\rcrat(q) \approx_{\psi+\xi+\zeta} y_\zeta),\label{eq:RCappx-rtri-rrl2}
    \end{equation}
    不过在此情形中我们并不需要这个假设。
    事实上，令 $\psi\defeq \phi$，\eqref{eq:RCappx-rtri-rrl1} 立即给出 $\rcrat(q) \approx_{\phi+\epsilon-\delta} y_\delta$，从而由 $\approx$ 的定义即得 $\rcrat(q) \approx_{\phi+\epsilon} \rclim(y)$。
  \item 我们假设 $\phi,\epsilon,\delta:\Qp$ 满足 $\rcrat(q)\approx_\phi \rclim(y)$ 与 $y_\delta \close{\epsilon-\delta} \rcrat(r)$，并归纳地假设
    \begin{gather}
      \fall{\psi:\Qp}(\rcrat(q) \approx_{\psi} y_\delta) \to (\rcrat(q) \approx_{\psi+\epsilon-\delta} \rcrat(r)).\label{eq:RCappx-rtri-rlr1}\\
      \fall{\psi,\xi,\zeta:\Qp} (\rcrat(q) \approx_{\psi} y_\xi) \to (\rcrat(q) \approx_{\psi+\xi+\zeta} y_\zeta).\label{eq:RCappx-rtri-rlr2}
    \end{gather}
    根据定义，$\rcrat(q)\approx_\phi \rclim(y)$ 意味着存在 $\theta:\Qp$ 使得 $\rcrat(q) \approx_{\phi-\theta} y_\theta$。
    因此，由假设~\eqref{eq:RCappx-rtri-rlr2}，我们也有 $\rcrat(q) \approx_{\phi+\delta} y_\delta$，然后由~\eqref{eq:RCappx-rtri-rlr1} 可知 $\rcrat(q) \approx_{\phi+\epsilon} \rcrat(r)$，正如所求。
  \item 我们假设 $\phi,\epsilon,\delta,\eta:\Qp$ 满足 $\rcrat(q)\approx_\phi \rclim(y)$ 与 $y_\delta \close{\epsilon-\delta-\eta} z_\eta$，并归纳地假设
    \begin{gather}
      \fall{\psi:\Qp}(\rcrat(q) \approx_{\psi} y_\delta) \to (\rcrat(q) \approx_{\psi+\epsilon-\delta-\eta} z_\eta), \label{eq:RCappx-rtri-rll1}\\
      \fall{\psi,\xi,\zeta:\Qp} (\rcrat(q) \approx_{\psi} y_\xi) \to (\rcrat(q) \approx_{\psi+\xi+\zeta} y_\zeta), \label{eq:RCappx-rtri-rll2}\\
      \fall{\psi,\xi,\zeta:\Qp} (\rcrat(q) \approx_{\psi} z_\xi) \to (\rcrat(q) \approx_{\psi+\xi+\zeta} z_\zeta). \label{eq:RCappx-rtri-rll3}
    \end{gather}
    同样，$\rcrat(q)\approx_\phi \rclim(y)$ 意味着存在 $\xi:\Qp$ 使得 $\rcrat(q) \approx_{\phi-\xi} y_\xi$，而~\eqref{eq:RCappx-rtri-rll2} 蕴含 $\rcrat(q) \approx_{\phi+\delta} y_\delta$，~\eqref{eq:RCappx-rtri-rll1} 蕴含 $\rcrat(q) \approx_{\phi+\epsilon-\eta} z_\eta$。
    但由 $\approx$ 的定义，这蕴含了即得所求的 $\rcrat(q) \approx_{\phi+\epsilon} \rclim(z)$。
  \end{enumerate}
  现在我们考虑 $u$ 是 $\rclim(x)$ 的情形，其中 $x$ 是一个 Cauchy 逼近。
  在此情形下，定义 $(\rclim(x) \approx_{(\blank)} {\blank}) : A$ 时的外围归纳假设为 ${(x_\delta \approx_{(\blank)} {\blank})}: A$，因此这些关系除了舍入性之外还满足右侧的三角形不等式。
  \begin{enumerate}\setcounter{enumi}{4}
  \item 假设 $\rclim(x)\approx_\phi \rcrat(r)$ 且 $-\epsilon<|r-s|<\epsilon$，我们必须证明 $\rclim(x)\approx_{\phi+\epsilon} \rcrat(s)$。
    根据 $\approx$ 的定义，前者意味着 $x_\delta \approx_{\phi-\delta} \rcrat(r)$，因此上述三角形不等式蕴含 $x_\delta \approx_{\epsilon+\phi-\delta} \rcrat(s)$，从而即得所求的 $\rclim(x)\approx_{\phi+\epsilon} \rcrat(s)$。
  \item 我们假设 $\phi,\epsilon,\delta:\Qp$ 满足 $\rclim(x)\approx_\phi \rcrat(r)$ 与 $\rcrat(r) \close{\epsilon-\delta} y_\delta$，以及两个多余的归纳假设。
    %
    根据定义，存在 $\eta:\Qp$ 使得 $x_\eta \approx_{\phi-\eta} \rcrat(r)$，因此归纳的三角形不等式给出 $x_\eta \approx_{\phi+\epsilon-\eta-\delta} y_\delta$。
    由 $\approx$ 的定义立得 $\rclim(x) \approx_{\phi+\epsilon} \rclim(y)$。
  \item 我们假设 $\phi,\epsilon,\delta:\Qp$ 满足 $\rclim(x)\approx_\phi \rclim(y)$ 与 $y_\delta \close{\epsilon-\delta} \rcrat(r)$，以及两个多余的归纳假设。
    根据定义，存在 $\xi,\theta:\Qp$ 使得 $x_\xi \approx_{\phi-\xi-\theta} y_\theta$。
    由于 $y$ 是一个 Cauchy 逼近，我们有 $y_\theta \close{\theta+\delta} y_\delta$，因此由归纳的三角形不等式可得 $x_\xi \approx_{\phi+\delta-\xi} y_\delta$，进而 $x_\xi \close{\phi+\epsilon-\xi} \rcrat(r)$。
    再由 $\approx$ 的定义即得 $\rclim(x) \approx_{\phi+\epsilon}\rcrat(r)$，正如所求。
  \item 最后，我们假设 $\phi,\epsilon,\delta,\eta:\Qp$ 满足 $\rclim(x)\approx_\phi \rclim(y)$ 与 $y_\delta \close{\epsilon-\delta-\eta} z_\eta$。
    那么如前所述，存在 $\xi,\theta:\Qp$ 使得 $x_\xi \approx_{\phi-\xi-\theta} y_\theta$，如前一样两次应用三角形不等式足矣。
  \end{enumerate}

  这就完成了两个内层递归，从而定义了关系族 $(\rcrat(q)\approx_{(\blank)}\blank)$ 和 $(\rclim(x)\approx_{(\blank)}\blank)$。
  因为它们都是 $A$ 的元素，所以是舍入的，并且对 $\closesym$ 满足右侧的三角形不等式。
% , and satisfy~\eqref{eq:appxrec2}.
  余下的工作是验证与 $\bsim$ 相关的条件，即这些关系对 $\closesym$ 的构造子满足\emph{左侧}的三角形不等式。
  四种情形对应于~\eqref{eq:RC-sim-ltri} 中 $u$ 和 $v$ 取有理点或极限的四种选择，且由于它们均为命题，我们可以应用 $\RC$-归纳法并假设 $w$ 也是有理点或极限。
  这又产生八种情形，其证明与刚给出的那些本质相同；因此我们就不再拿它们来折磨读者了。
\end{proof}

现在我们可以证明：

\begin{thm}\label{thm:RC-sim-characterization}
  对任意 $u,v:\RC$ 和 $\epsilon:\Qp$，有 $(u\close\epsilon v) = (u\approx_\epsilon v)$。
\end{thm}

\begin{proof}
  由于两边都是命题，只需证明双向蕴含。
  对于从左到右的方向，我们对 $C(u,v,\epsilon)\defeq (u\approx_\epsilon v)$ 使用 $\closesym$-归纳。
  因此，只需考虑 $\closesym$ 的四个构造子。
  在每种情形中，$u$ 和 $v$ 都被特化为有理点或极限，从而 $\approx$ 的定义可直接计算，归纳假设也总是适用。

  对于从右到左的方向，我们使用 $\RC$-归纳，假设 $u$ 和 $v$ 是有理点或极限，从而 $\approx$ 可求值。
  此时 $\approx$ 的定义以及归纳假设恰好提供了 $\closesym$ 相应构造子所需的数据。
\end{proof}

\index{encode-decode method}%
略加引申，我们可以将 $\approx$ 称为 $\closesym$ 的``代码''纤维化，而上述证明的两个方向分别对应 \encode 和 \decode。
根据 $\approx$ 的定义，由 \cref{thm:RC-sim-characterization} 我们得到等价
\begin{align*}
  (\rcrat(q) \close\epsilon \rcrat(r))  &=
  (-\epsilon < q - r < \epsilon)\\
  (\rcrat(q) \close\epsilon \rclim(y)) &=
  \exis{\delta : \Qp} \rcrat(q) \close{\epsilon - \delta} y_\delta\\
  (\rclim(x) \close\epsilon \rcrat(r)) &=
  \exis{\delta : \Qp} x_\delta \close{\epsilon - \delta} \rcrat(r)\\
  (\rclim(x) \close\epsilon \rclim(y)) &=
  \exis{\delta, \eta : \Qp} x_\delta \close{\epsilon - \delta - \eta} y_\eta.
\end{align*}
我们的证明还给出以下附加信息。

\begin{cor}
  \index{triangle!inequality for R@inequality for $\RC$}%
  \indexsee{inequality!triangle}{triangle inequality}%
  $\closesym$ 是\index{rounded!relation}舍入的，并且满足三角不等式：
    \begin{gather}
      \eqvspaced{
        (u \close\epsilon v)
      }{
        \exis{\theta : \Qp} u \close{\epsilon - \theta} v
      }\\
      (u\close\epsilon v) \to (v\close\delta w) \to (u\close{\epsilon+\delta} w). \label{item:RC-sim-triangle}
    \end{gather}
\end{cor}


% \begin{proof}
%   The construction of $\approx$ showed simultaneously that it is rounded, and satisfies ``triangle inequalities'' such as
%   \[ (u\approx_\epsilon v) \to (v\close\delta w) \to (u\approx_{\epsilon+\delta} w). \]
%   Thus, both properties follow from \cref{thm:RC-sim-characterization}.
% \end{proof}

有了三角不等式，我们就可以证明 Cauchy 逼近的``极限''确实表现得像极限。

\begin{lem}\label{thm:RC-sim-lim}
  对任意 $u:\RC$，任意 Cauchy 逼近 $y$，以及任意 $\epsilon,\delta:\Qp$，若 $u\close\epsilon y_\delta$，则 $u\close{\epsilon+\delta} \rclim(y)$。
\end{lem}

\begin{proof}
  我们对 $u$ 使用 $\RC$-归纳。
  若 $u$ 是 $\rcrat(q)$，这正是 $\closesym$ 的第二个构造子。
  现在假设 $u$ 是 $\rclim(x)$，并且每个 $x_\eta$ 都满足如下性质：对任意 $y,\epsilon,\delta$，若 $x_\eta\close\epsilon y_\delta$，则 $x_\eta \close{\epsilon+\delta} \rclim(y)$。
  特别地，在此假设中取 $y\defeq x$ 且 $\delta\defeq\eta$，可得：对任意 $\eta,\theta:\Qp$，有 $x_\eta \close{\eta+\theta} \rclim(x)$。

  现在任取 $y,\epsilon,\delta$ 并假设 $\rclim(x) \close\epsilon y_\delta$。
  由舍入性，存在 $\theta$ 使得 $\rclim(x) \close{\epsilon-\theta} y_\delta$。
  那么根据上面的观察，对任意 $\eta$ 有 $x_\eta \close{\eta+\theta/2} \rclim(x)$，进而由三角不等式得到 $x_\eta \close{\epsilon+\eta-\theta/2} y_\delta$。
  因此，$\closesym$ 的第四个构造子给出 $\rclim(x) \close{\epsilon+2\eta+\delta-\theta/2} \rclim(y)$。
  于是，取 $\eta \defeq \theta/4$，即得所求。
\end{proof}

\begin{lem}\label{thm:RC-sim-lim-term}
  对任意 Cauchy 逼近 $y$ 及任意 $\delta,\eta:\Qp$，有 $y_\delta \close{\delta+\eta} \rclim(y)$。
\end{lem}

\begin{proof}
  在前一个引理中取 $u\defeq y_\delta$ 和 $\epsilon\defeq \eta$ 即可。
\end{proof}

\begin{rmk}
  我们可能会期望 $y_\delta \close{\delta} \rclim(y)$ 成立，但这一性质在某些例子中并不成立。
  例如，考虑由 $x_\epsilon \defeq \epsilon$ 定义的 $x$。
  其极限显然是 $0$，但 $|\epsilon - 0 |<\epsilon$ 并不成立，只有 $\le$ 成立。
\end{rmk}

作为应用，\cref{thm:RC-sim-lim-term} 使我们能够证明 Lipschitz 函数从 \cref{RC-extend-Q-Lipschitz} 出发的扩展是唯一的。

\begin{lem}\label{RC-continuous-eq}
  \index{function!continuous}%
  设 $f,g:\RC\to\RC$ 连续，即满足
  \[ \fall{u:\RC}{\epsilon:\Qp}\exis{\delta:\Qp}\fall{v:\RC} (u\close\delta v) \to (f(u) \close\epsilon f(v)) \]
  对 $g$ 亦同。
  若对所有 $q:\Q$ 有 $f(\rcrat(q))=g(\rcrat(q))$，则 $f=g$。
\end{lem}

\begin{proof}
  我们通过 $\RC$-归纳证明对所有 $u$ 有 $f(u)=g(u)$。
  有理数的情形就是假设本身。
  于是，假设对所有 $\delta$ 有 $f(x_\delta)=g(x_\delta)$。
  我们将证明对所有 $\epsilon$ 有 $f(\rclim(x))\close\epsilon g(\rclim(x))$，从而可应用 $\RC$ 的道路构造子。

  由于 $f$ 和 $g$ 连续，故存在 $\theta,\eta$ 使得对所有 $v$ 有
  \begin{align*}
    (\rclim(x)\close\theta v) &\to (f(\rclim(x)) \close{\epsilon/2} f(v))\\
    (\rclim(x)\close\eta v) &\to (g(\rclim(x)) \close{\epsilon/2} g(v)).
  \end{align*}
  取 $\delta < \min(\theta,\eta)$，由 \cref{thm:RC-sim-lim-term} 我们既有 $\rclim(x)\close\theta y_\delta$ 也有 $\rclim(x)\close\eta y_\delta$。
  因此
  \[ f(\rclim(x)) \close{\epsilon/2} f(y_\delta) = g(y_\delta) \close{\epsilon/2} g(\rclim(x))\]
  于是由三角不等式可得 $f(\rclim(x))\close\epsilon g(\rclim(x))$。
\end{proof}

\subsection{Cauchy 实数的代数结构}
\label{sec:algebr-struct-cauchy}

我们首先定义加法结构 $(\RC, 0, {+}, {-})$。显然，加法单位元 $0$ 就是 $\rcrat(0)$，而加法逆元 ${-} : \RC \to \RC$ 则是利用 \cref{RC-extend-Q-Lipschitz} 将 $\Q$ 上的加法逆元 ${-} : \Q \to \Q$ 扩展而得，Lipschitz 常数为 $1$。加法本身的定义则需要多花一些功夫。

\begin{lem} \label{RC-binary-nonexpanding-extension}
  假设 $f : \Q \times \Q \to \Q$ 满足，对所有 $q, r, s : \Q$，
  %
  \begin{equation*}
    |f(q, s) - f(r, s)| \leq |q - r|
    \qquad\text{and}\qquad
    |f(q, r) - f(q, s)| \leq |r - s|.
  \end{equation*}
  %
  那么存在函数 $\bar{f} : \RC \times \RC \to \RC$，使得对所有 $q, r : \Q$ 有 $\bar{f}(\rcrat(q), \rcrat(r)) = f(q,r)$。此外，对所有 $u, v, w : \RC$ 和 $q : \Qp$，
  %
  \begin{equation*}
    u \close\epsilon v \Rightarrow \bar{f}(u,w) \close\epsilon \bar{f}(v,w)
    \quad\text{and}\quad
    v \close\epsilon w \Rightarrow \bar{f}(u,v) \close\epsilon \bar{f}(u,w).
  \end{equation*}
\end{lem}

\begin{proof}
  我们使用 $(\RC, {\closesym})$-递归来构造 $\bar{f}$ 的柯里化形式，即一个映射
  $\RC \to A$，其中 $A$ 是取值于实数的非扩张\index{function!non-expanding}\index{non-expanding function}函数构成的空间：
  %
  \begin{equation*}
    A \defeq
    \setof{ h : \RC \to \RC |
      \fall{\epsilon : \Qp} \fall{u, v : \RC}
      u \close\epsilon v \Rightarrow h(u) \close\epsilon h(v)
    }.
  \end{equation*}
  %
  我们还需要 $A$ 上的一个合适的 $\bsim_\epsilon$，其定义为
  %
  \begin{equation*}
    (h \bsim_\epsilon k) \defeq \fall{u : \RC} h(u) \close\epsilon k(u).
  \end{equation*}
  %
  显然，若 $\fall{\epsilon : \Qp} h \bsim_\epsilon k$ 则对所有 $u : \RC$ 都有 $h(u) = k(u)$，因此 $\bsim$ 是分离的。

  对于基础情况，我们将 $\bar{f}(\rcrat(q)) : A$（其中 $q : \Q$）定义为 Lipschitz 映射 $\lam{r} f(q,r)$ 从 $\Q \to \Q$ 到 $\RC \to \RC$ 的延拓，这由 \cref{RC-extend-Q-Lipschitz} 以 Lipschitz 常数~$1$ 构造而成。
  接着，对于一个 Cauchy 逼近 $x$，我们将 $\bar{f}(\rclim(x)) : \RC \to \RC$ 定义为
  %
  \begin{equation*}
    \bar{f}(\rclim(x))(v) \defeq \rclim (\lam{\epsilon} \bar{f}(x_\epsilon)(v)).
  \end{equation*}
  %
  为使此定义有效，$\lam{\epsilon} \bar{f}(x_\epsilon)(v)$ 应当是一个 Cauchy 逼近。为此，考虑任意 $\delta, \epsilon : \Q$。
  则由假设有 $\bar{f}(x_\delta) \bsim_{\delta + \epsilon} \bar{f}(x_\epsilon)$，于是
  $\bar{f}(x_\delta)(v) \close{\delta + \epsilon} \bar{f}(x_\epsilon)(v)$。
  此外，$\bar{f}(\rclim(x))$ 是非扩张的，因为由归纳假设，$\bar{f}(x_\epsilon)$ 是非扩张的。
  事实上，若 $u \close\epsilon v$，则对所有 $\epsilon : \Q$，
  %
  \begin{equation*}
    \bar{f}(x_{\epsilon/3})(u) \close{\epsilon/3} \bar{f}(x_{\epsilon/3})(v),
  \end{equation*}
  %
  因此，由 $\closesym$ 的第四个构造子可得 $\bar{f}(\rclim(x))(u) \close\epsilon \bar{f}(\rclim(x))(v)$。

  我们还需要验证另外四个条件，这里仅举例说明其中一个。
  假设 $\epsilon : \Qp$，且对某个 $\delta : \Qp$ 有 $\rcrat(q) \close{\epsilon - \delta}
  y_\delta$ 与 $\bar{f}(\rcrat(q)) \bsim_{\epsilon - \delta} \bar{f}(y_\delta)$。
  为证明 $\bar{f}(\rcrat(q)) \bsim_\epsilon \bar{f}(\rclim(y))$，考虑任意 $v : \RC$ 并注意到
  %
  \begin{equation*}
    \bar{f}(\rcrat(q))(v) \close{\epsilon - \delta} \bar{f}(y_\delta)(v).
  \end{equation*}
  %
  因此，由 $\closesym$ 的第二个构造子，我们得到
  \narrowequation{\bar{f}(\rcrat(q))(v) \close\epsilon \bar{f}(\rclim(y))(v)}
  如所欲证。
\end{proof}

我们可以将 \cref{RC-binary-nonexpanding-extension} 应用于任何在每个变量上分别非扩张的二元有理函数。
加法就是这样一个函数，因此我们得到了 ${+} : \RC \times \RC \to \RC$。
\indexdef{addition!of Cauchy reals}%
此外，只要要求延拓在每个变量上非扩张，延拓就是唯一的；并且正如单变量情形一样，有理数上的等式可延拓为实数上的等式。
由于非扩张映射的复合仍是非扩张的，我们可以断定加法满足通常的性质，如交换律与结合律。
\index{associativity!of addition!of Cauchy reals}%
因此，$(\RC, 0, {+}, {-})$ 是一个交换群。

我们也可以将 \cref{RC-binary-nonexpanding-extension} 应用于函数 $\min : \Q \times
\Q \to \Q$ 和 $\max : \Q \times \Q \to \Q$，从而使 $\RC$ 成为一个格。
$\RC$ 上的偏序 $\leq$ 通过 $\max$ 定义为
%
\symlabel{leq-RC}
\index{order!non-strict}%
\index{non-strict order}%
\begin{equation*}
  (u \leq v) \defeq (\max(u, v) = v).
\end{equation*}
%
关系 $\leq$ 是偏序，因为它在 $\Q$ 上就是偏序，而偏序公理可以用关于 $\min$ 和 $\max$ 的等式来表达，所以这些性质传递到了 $\RC$。

\index{absolute value}%
另一个通过同样方法可延拓到 $\RC$ 的函数是绝对值 $|{\blank}|$。
同样，它具有预期的性质，因为这些性质从 $\Q$ 传递到了 $\RC$。

\symlabel{lt-RC}
由 $\leq$ 定义严格序 $<$ 如下：
\index{strict!order}%
\index{order!strict}%
%
\begin{equation*}
  (u < v) \defeq \exis{q, r : \Q} (u \leq \rcrat(q)) \land (q < r) \land (\rcrat(r) \leq v).
\end{equation*}
%
也就是说，当 merely 存在一对有理数 $q < r$ 使得 $x \leq
\rcrat(q)$ 且 $\rcrat(r) \leq v$ 时，$u < v$ 成立。
不难验证 $<$ 不自反且传递，并满足有序域的其他预期性质。
Archimedean 原理直接由 $<$ 的定义得出。

\index{ordered field!archimedean}%


\begin{thm}[$\RC$ 的Archimedean 原理] \label{RC-archimedean}
  %
  对于所有满足 $u < v$ 的 $u, v : \RC$，mere 存在 $q : \Q$ 使得 $u < q < v$。
\end{thm}

\begin{proof}
  由 $u < v$，我们 merely 得到 $r, s : \Q$ 使得 $u \leq r < s \leq v$，此时可取 $q
  \defeq (r + s) / 2$。
\end{proof}

现在我们在 $\RC$ 上已具备足够的结构，可以用标准概念来表达 $u \close\epsilon v$。

\begin{lem}\label{thm:RC-le-grow}
  如果 $q:\Q$ 和 $u:\RC$ 满足 $u\le \rcrat(q)$，那么对于任意 $v:\RC$ 和 $\epsilon:\Qp$，若 $u\close\epsilon v$ 则 $v\le \rcrat(q+\epsilon)$。
\end{lem}

\begin{proof}
  注意到函数 $\max(\rcrat(q),\blank):\RC\to\RC$ 是 Lipschitz 的，其常数为 $1$。
  先考虑 $u=\rcrat(r)$ 为有理数的情形。
  为此，我们对 $v$ 施行归纳。
  如果 $v$ 是有理数，那么结论是显然的。
  如果 $v$ 是 $\rclim(y)$，我们归纳地假设：对任意 $\epsilon,\delta$，如果 $\rcrat(r)\close\epsilon y_\delta$，则 $y_\delta \le \rcrat(q+\epsilon)$，即 $\max(\rcrat(q+\epsilon),y_\delta)=\rcrat(q+\epsilon)$。

  现在假设给定 $\epsilon$ 且 $\rcrat(r)\close\epsilon \rclim(y)$，那么存在 $\theta$ 使得 $\rcrat(r)\close{\epsilon-\theta} \rclim(y)$ 成立，从而当 $\delta<\theta$ 时就有 $\rcrat(r)\close\epsilon y_\delta$。
  于是由归纳假设，对这样的 $\delta$ 有 $\max(\rcrat(q+\epsilon),y_\delta)=\rcrat(q+\epsilon)$。
  但根据定义，
  \[\max(\rcrat(q+\epsilon),\rclim(y)) \jdeq \rclim(\lam{\delta} \max(\rcrat(q+\epsilon),y_\delta)).\]
  由于一个最终常值的 Cauchy 逼近的极限就是该常值，我们有
  \[\max(\rcrat(q+\epsilon),\rclim(y)) = \rcrat(q+\epsilon),\]从而得到 $\rclim(y)\le \rcrat(q+\epsilon)$。

  现在考虑一般的 $u:\RC$。
  因为 $u\le \rcrat(q)$ 意味着 $\max(\rcrat(q),u)=\rcrat(q)$，结合假设 $u\close\epsilon v$ 以及 $\max(\rcrat(q),-)$ 的 Lipschitz 性质，可推出 $\max(\rcrat(q),v) \close\epsilon \rcrat(q)$。
  这样，既然 $\rcrat(q)\le \rcrat(q)$，由第一种情形得 $\max(\rcrat(q),v) \le \rcrat(q+\epsilon)$，再由 $\le$ 的传递性即得 $v\le \rcrat(q+\epsilon)$。
\end{proof}

\begin{lem}\label{thm:RC-lt-open}
  设 $q:\Q$ 和 $u:\RC$ 满足 $u<\rcrat(q)$。则：
  \begin{enumerate}
  \item 对任意 $v:\RC$ 和 $\epsilon:\Qp$，若 $u\close\epsilon v$，则 $v< \rcrat(q+\epsilon)$。\label{item:RCltopen1}
  \item 存在 $\epsilon:\Qp$，使得对任意 $v:\RC$，只要 $u\close\epsilon v$，就有 $v<\rcrat(q)$。\label{item:RCltopen2}
  \end{enumerate}
\end{lem}

\begin{proof}
  根据定义，$u<\rcrat(q)$ 意味着存在 $r:\Q$，满足 $r<q$ 且 $u\le \rcrat(r)$。
  于是由 \cref{thm:RC-le-grow}，对任意 $\epsilon$，如果 $u\close\epsilon v$，则 $v\le \rcrat(r+\epsilon)$。
  结论~\ref{item:RCltopen1} 直接可得，因为 $r+\epsilon<q+\epsilon$；至于~\ref{item:RCltopen2}，只需取任意 $\epsilon <q-r$ 即可。
\end{proof}

现在我们可以证明，辅助关系 $\closesym$ 正是我们所预期的。

\begin{thm} \label{RC-sim-eqv-le}
  \index{distance}%
  $\eqv{(u \close\epsilon v)}{(|u - v| < \rcrat(\epsilon))}$
  对所有 $u, v : \RC$ 和 $\epsilon : \Qp$ 成立。
\end{thm}

\begin{proof}
  减法与绝对值的 Lipschitz 性质表明：如果 $u\close\epsilon v$，那么 $|u-v| \close\epsilon |u-u| = 0$。
  因此，对于从左到右的方向，只需证明：如果 $u\close\epsilon 0$，那么 $|u|<\rcrat(\epsilon)$。
  我们对 $u$ 进行 $\RC$-归纳。

  如果 $u$ 是有理数，那么结论直接成立，因为绝对值与序延拓了 $\Qp$ 上的标准定义。
  如果 $u$ 是 $\rclim(x)$，那么由舍入性，存在 $\theta:\Qp$ 使得 $\rclim(x)\close{\epsilon-\theta} 0$。
  于是，由三角不等式，我们有 $x_{\theta/3} \close{\epsilon-2\theta/3} 0$，从而归纳假设给出 $|x_{\theta/3}|<\rcrat(\epsilon-2\theta/3)$。
  但是 $x_{\theta/3} \close{2\theta/3} \rclim(x)$，因此由 Lipschitz 性质得到 $|x_{\theta/3}| \close{2\theta/3} |\rclim(x)|$，这样一来，\cref{thm:RC-lt-open}\ref{item:RCltopen1} 蕴涵 $|\rclim(x)|<\rcrat(\epsilon)$。

  对于另一个方向，我们对 $u$ 和 $v$ 进行 $\RC$-归纳。
  如果两者都是有理数，那么这就是 $\closesym$ 的第一个构造子。

  如果 $u$ 是 $\rcrat(q)$ 而 $v$ 是 $\rclim(y)$，我们归纳地假设：对于任意 $\epsilon,\delta$，如果 $|\rcrat(q)-y_\delta|<\rcrat(\epsilon)$ 那么 $\rcrat(q) \close{\epsilon} y_\delta$。
  取定一个满足 $|\rcrat(q) - \rclim(y)|<\rcrat(\epsilon)$ 的 $\epsilon$。
  由于 $\Q$ 在 $\RC$ 中序稠密，存在 $\theta<\epsilon$ 使得 $|\rcrat(q) - \rclim(y)|<\rcrat(\theta)$。
  现在，对任意 $\delta,\eta$ 我们有 $\rclim(y)\close{2\delta} y_\delta$，因此由 Lipschitz 性质，
  \[ |\rcrat(q) - \rclim(y)| \close{\delta+\eta} |\rcrat(q) - y_\delta|. \]
  于是，由 \cref{thm:RC-lt-open}\ref{item:RCltopen1}，我们得到 $|\rcrat(q) - y_\delta| < \rcrat(\theta+2\delta)$。
  那么由归纳假设，$\rcrat(q) \close{\theta+2\delta} y_\delta$，从而由三角不等式得到 $\rcrat(q)\close{\theta+4\delta} \rclim(y)$。
  因此，只需取 $\delta \defeq (\epsilon-\theta)/4$ 即可。

  剩下的两种情况完全类似。
\end{proof}

\indexdef{multiplication!of Cauchy reals}%
接下来，我们希望为 $\RC$ 配备乘法结构。
对每个 $q : \Q$，映射 $r \mapsto q \cdot r$ 是 Lipschitz 的，其常数为\footnote{我们定义 Lipschitz 常数为\emph{正的}有理数。} $|q| + 1$，因此我们可以将其延拓为实数上乘以 $q$ 的乘法。
所以 $\RC$ 是 $\Q$ 上的向量空间\index{vector!space}。
一般来说，我们可以将实数的乘法定义为
%
\begin{equation}
  u \cdot v \defeq
  {\textstyle \frac{1}{2}} \cdot ((u + v)^2 - u^2 - v^2),\label{mult-from-square}
\end{equation}
%
因此我们只需要平方函数\index{squaring function} $u \mapsto u^2$ 作为映射 $\RC \to \RC$。
平方函数不是 Lipschitz 映射，但在每个有界域上它是 Lipschitz 的，据此可以将其拼接起来。
定义开区间和闭区间为
%
\indexdef{interval!open and closed}%
\indexdef{open!interval}%
\indexdef{closed!interval}%
\begin{equation*}
  [u,v] \defeq \setof{ x : \RC | u \leq x \leq v }
  \qquad\text{and}\qquad
  (u,v) \defeq \setof{ x : \RC | u < x < v }.
\end{equation*}
%
虽然严格来说，$[u,v]$ 或 $(u,v)$ 的元素是一个 Cauchy 实数连同其一个证明，但由于后者仅居于一个命题，故无关紧要。
因此，正如处理子集类型时的常见做法，每当 $x:\RC$ 满足 $u\leq x \leq v$ 时，我们通常就简记为 $x:[u,v]$，其余类似。

\begin{thm} \label{RC-squaring}
  %
  存在唯一的函数 ${(\blank)}^2 : \RC \to \RC$，它延拓有理数的平方映射 $q \mapsto
  q^2$，并且满足
  %
  \begin{equation*}
    \fall{n : \N}
    \fall{u, v : [-n, n]}
    |u^2 - v^2| \leq 2 \cdot n \cdot |u - v|.
  \end{equation*}
\end{thm}

\begin{proof}
  我们首先观察到，对于每个 $u : \RC$，总存在 $n : \N$ 使得 $-n \leq u \leq n$，参见 \cref{ex:traditional-archimedean}，因此映射
  %
  \begin{equation*}
    e : \Parens{\sm{n : \N} [-n, n]} \to \RC
    \qquad\text{defined by}\qquad
    e(n, x) \defeq x
  \end{equation*}
  %
  是满射。接着，对于每个 $n : \N$，平方映射
  %
  \begin{equation*}
    s_n : \setof{ q : \Q | -n \leq q \leq n } \to \Q
    \qquad\text{defined by}\qquad
    s_n(q) \defeq q^2
  \end{equation*}
  %
  具有 Lipschitz 常数 $2 n$，于是我们可以利用 \cref{RC-extend-Q-Lipschitz} 将其延拓为一个具有 Lipschitz 常数 $2 n$ 的映射 $\bar{s}_n : [-n, n] \to \RC$，具体细节参见 \cref{RC-Lipschitz-on-interval}。这些映射 $\bar{s}_n$ 是相容的：如果对某个 $m, n : \N$ 有 $m < n$，那么 $s_n$ 在 $[-m, m]$ 上的限制必须与 $s_m$ 一致，因为二者都是 Lipschitz 的，因此在~\cref{RC-continuous-eq} 的意义下是连续的。因此，根据 \cref{lem:images_are_coequalizers}，映射
  %
  \begin{equation*}
    \Parens{\sm{n : \N} [-n, n]} \to \RC,
    \qquad\text{given by}\qquad
    (n, x) \mapsto s_n(x)
  \end{equation*}
  %
  可以唯一地通过 $\RC$ 分解，从而给出我们想要的函数。
\end{proof}

至此，我们已经有了实数的环结构和 Archimedean 序。为了确立 $\RC$ 是一个 Archimedean 有序域，还需要证明逆元的存在性。

\begin{thm}
  \index{apartness}%
  一个 Cauchy 实数可逆当且仅当它与零分离。
\end{thm}

\begin{proof}
  首先，假设 $u : \RC$ 有一个逆元 $v : \RC$。根据 Archimedean 原理，存在 $q : \Q$ 使得 $|v| < q$。那么 $1 = |u v| < |u| \cdot v < |u| \cdot q$，因此 $|u| > 1/q$，也就是说 $u \apart 0$。

  反之，与 \cref{RC-squaring} 中平方函数的构造类似，我们通过粘合函数来构造逆映射
  %
  \begin{equation*}
    ({\blank})^{-1} : \setof{ u : \RC | u \apart 0 } \to \RC
  \end{equation*}
  %
  。我们仅概述主要步骤。对于每个 $q : \Q$，令
  %
  \begin{equation*}
    [q, \infty) \defeq \setof{u : \RC | q \leq u}
    \qquad\text{and}\qquad
    (-\infty, q] \defeq \setof{u : \RC | u \leq -q}.
  \end{equation*}
  %
  那么，当 $q$ 取遍 $\Qp$ 时，类型 $(-\infty, q]$ 和 $[q, \infty)$ 共同覆盖了 $\setof{u : \RC | u \apart 0}$。在每个这样的 $[q, \infty)$ 和 $(-\infty, q]$ 上，应用 \cref{RC-extend-Q-Lipschitz} 即可得到具有 Lipschitz 常数 $1/q^2$ 的逆函数。最后，\cref{lem:images_are_coequalizers} 保证逆函数可唯一地经由 $\setof{ u : \RC | u
    \apart 0 }$ 分解。
\end{proof}

我们用一个定理来总结 $\RC$ 的代数结构。

\begin{thm} \label{RC-archimedean-ordered-field}
  Cauchy 实数构成一个 Archimedean 有序域。
\end{thm}

\subsection{Cauchy 实数是 Cauchy 完备的}
\label{sec:cauchy-reals-cauchy-complete}

我们通过在有理数 $\Q$ 上对 Cauchy 逼近取极限来构造 $\RC$，因此 $\RC$ 理应是 Cauchy 完备的。得益于 \cref{RC-sim-eqv-le}，在 $\RC$ 构造中定义的 Cauchy 逼近 $x : \Qp \to \RC$，与 \cref{defn:cauchy-approximation} 中定义的（适应于 $\RC$ 的）Cauchy 逼近之间没有区别。

这样一来，给定一个 Cauchy 逼近 $x : \Qp \to \RC$，自然期望 $\rclim(x)$ 就是它的极限，这里极限的概念是按 \cref{defn:cauchy-approximation} 定义的。而根据 \cref{RC-sim-eqv-le} 和 \cref{thm:RC-sim-lim-term}，情况正是如此。我们已经证明了：

\begin{thm}
  $\RC$ 中的每个 Cauchy 逼近都有一个极限。
\end{thm}

若某 Archimedean 有序域中每个 Cauchy 逼近都有极限，则称之为
\define{Cauchy 完备的（Cauchy complete）}。
\indexdef{Cauchy!completeness}%
\indexdef{complete!ordered field, Cauchy}%
\index{ordered field}%
Cauchy 实数是这样的最小域。

\begin{thm} \label{RC-initial-Cauchy-complete}
  Cauchy 实数可以嵌入到每个 Cauchy 完备的 Archimedean 有序域中。
\end{thm}

\begin{proof}
  \index{limit!of a Cauchy approximation}%
  设 $F$ 是一个 Cauchy 完备的 Archimedean 有序域。由于极限是唯一的，
  存在一个算子 $\lim$，将 $F$ 中的 Cauchy 逼近映为其极限。我们
  通过 $(\RC, {\closesym})$-递归定义嵌入 $e : \RC \to F$ 如下
  %
  \begin{equation*}
    e(\rcrat(q)) \defeq q
    \qquad\text{and}\qquad
    e(\rclim(x)) \defeq \lim (e \circ x).
  \end{equation*}
  %
  $F$ 上的一个合适的 $\bsim$ 为
  %
  \begin{equation*}
    (a \bsim_\epsilon b) \defeq |a - b| < \epsilon.
  \end{equation*}
  %
  因为 $F$ 是 Archimedean 的，所以这是一个分离关系。$(\RC, {\closesym})$-递归的其余条款
  很容易验证。此外还需要验证 $e$ 是一个固定有理数的有序域嵌入。
\end{proof}

\index{real numbers!Cauchy|)}%

\section{Cauchy 实数与 Dedekind 实数的比较}
\label{sec:comp-cauchy-dedek}

\index{real numbers!Dedekind|(}%
\index{real numbers!Cauchy|(}%
\index{depression|(}

下面我们也讨论一下 Cauchy 实数与 Dedekind 实数之间的关系。由
\cref{RC-archimedean-ordered-field}，$\RC$ 是一个 Archimedean 有序域。它对于 $\Omega$ 也是
\define{可容许的}\index{ordered field!admissible}，这很容易验证。（若 $\Omega$ 为初始
$\sigma$-帧
\index{initial!sigma-frame@$\sigma$-frame}%
\index{sigma-frame@$\sigma$-frame!initial}%
则只需一个简单的归纳，而在其他情况下则是显然的。）
因此，由 \cref{RD-final-field}，存在一个有序域嵌入
%
\begin{equation*}
  \RC \to \RD
\end{equation*}
%
该嵌入固定有理数。
（我们也可以从 \cref{RC-initial-Cauchy-complete,RD-cauchy-complete} 得到这个嵌入。）
一般来说，如果没有进一步的假设，我们不期望 $\RC$ 与 $\RD$ 重合。

\begin{lem} \label{lem:untruncated-linearity-reals-coincide}
  %
  若对每个 $x : \RD$，都仅仅存在
  %
  \begin{equation}
    \label{eq:untruncated-linearity}
    c : \prd{q, r : \Q} (q < r) \to (q < x) + (x < r)
  \end{equation}
  %
  则 Cauchy 实数与Dedekind 实数重合。
\end{lem}

\begin{proof}
  注意，\eqref{eq:untruncated-linearity} 中的类型是 \eqref{eq:RD-linear-order} 的一个未截断变体，后者断言 $<$ 是弱线性序。
  我们已知 $\RC$ 可嵌入 $\RD$，因此只需证明每个Dedekind 实数仅仅是有理数柯西序列\index{Cauchy!sequence}的极限。

  考虑任意的 $x : \RD$。由假设，仅仅存在满足引理陈述中条件的 $c$；又由分割\index{cut!Dedekind}的非空性，仅仅存在 $a, b : \Q$ 使得 $a < x < b$。
  我们用递归构造序列\index{sequence} $f : \N \to \setof{ \pairr{q, r} \in \Q \times \Q | q < r }$：
  %
  \begin{enumerate}
  \item 令 $f(0) \defeq \pairr{a, b}$。
  \item 假设 $f(n)$ 已定义为 $\pairr{q_n, r_n}$ 且 $q_n < r_n$。
    定义 $s \defeq (2 q_n + r_n)/3$ 与 $t \defeq (q_n + 2 r_n)/3$。于是 $c(s,t)$ 可判定 $s < x$ 与 $x < t$ 二者孰是。
    若判定为 $s < x$，则令 $f(n+1) \defeq
    \pairr{s, r_n}$；否则令 $f(n+1) \defeq \pairr{q_n, t}$。
  \end{enumerate}
  %
  将序列~$f$ 的第 $n$ 项记作 $\pairr{q_n, r_n}$。不难看出，对所有 $n : \N$，有 $q_n < x < r_n$ 及 $|q_n - r_n| \leq (2/3)^n \cdot |q_0 - r_0|$。
  因此 $q_0, q_1, \ldots$ 和 $r_0, r_1, \ldots$ 都是收敛到Dedekind 分割~$x$ 的柯西序列。
  我们已证明，对每个 $x : \RD$，都仅仅存在收敛到 $x$ 的柯西序列。
\end{proof}

该引理表明，可数选择公理或排中律均足以保证 $\RC$ 与 $\RD$ 重合。

\begin{cor} \label{when-reals-coincide}
  \index{axiom!of choice!countable}%
  \index{excluded middle}%
  若排中律或可数选择公理成立，则 $\RC$ 与 $\RD$ 等价。
\end{cor}

\begin{proof}
  若排中律成立，则 $(x < y) \to (x < z) + (z < y)$ 可证：要么 $x < z$，要么 $\lnot (x < z)$。
  前一种情况即已完成，而后一种情况下，由 $z \leq x < y$ 得 $z < y$。
  因此，我们得到~\eqref{eq:untruncated-linearity}，从而可以应用 \cref{lem:untruncated-linearity-reals-coincide}。

  假设可数选择公理成立。集合 $S = \setof{ \pairr{q, r} \in \Q \times \Q | q <
    r }$ 与 $\N$ 等价，因此可以对``$x$ 可定位''这一陈述
  %
  \begin{equation*}
    \fall{\pairr{q, r} : S} (q < x) \lor (x < r).
  \end{equation*}
  %
  应用可数选择公理。
  注意，$(q < x) \lor (x < r)$ 可表达为一个存在性陈述 $\exis{b :
    \bool} (b = \bfalse \to q < x) \land (b = \btrue \to x < r)$。该选择函数的柯里化形式
  恰好就是~\eqref{eq:untruncated-linearity}，从而
  \cref{lem:untruncated-linearity-reals-coincide} 再次适用。
\end{proof}

\index{real numbers!Dedekind|)}%
\index{real numbers!Cauchy|)}%
\index{real numbers!agree}%

\index{depression|)}

\section{区间的紧性}
\label{sec:compactness-interval}

\index{mathematics!classical|(}%
\index{mathematics!constructive|(}%

我们此前已指出，对实数的构造与经典逻辑完全兼容。
因此，通过假设排中律~\eqref{eq:lem} 和选择公理~\eqref{eq:ac}，我们可以发展经典分析\index{classical!analysis}\index{analysis!classical}，其结果实质上等同于照搬任何一本标准分析教科书。

\index{analysis!constructive}%
\index{constructive!analysis}%
尽管如此，任何关注计算的人——例如数值分析师——都应当对在计算上有意义的环境中发展分析学感到好奇。
在构造性框架下进行分析甚至是可能的，这一点已由~\cite{Bishop1967} 所证明。
作为经典分析与构造性分析之异同的一个示例，我们将仅简要讨论一个主题——闭区间 $[0,1]$ 的紧性以及围绕这一概念的若干定理。

在构造性数学中，经典等价的概念往往会产生分叉，这是一个常见现象，而紧性也不例外。
最常用的三种紧性概念是：
%
\indexdef{compactness}%


\begin{enumerate}
\item \define{度量紧的（metrically compact）}：“Cauchy 完备且完全有界”，
  \indexdef{metrically compact}%
  \indexdef{compactness!metric}%
\item \define{Bolzano--Weierstra\ss{} 紧的（Bolzano--Weierstra\ss{} compact）}：“每个序列都有收敛子列”，
  \index{compactness!Bolzano--Weierstrass@Bolzano--Weierstra\ss{}}%
  \indexsee{Bolzano--Weierstrass@Bolzano--Weierstra\ss{}}{compactness}%
  \index{sequence}%
\item \define{Heine--Borel 紧的（Heine--Borel compact）}：“每个开覆盖都有有限子覆盖”。
  \index{compactness!Heine--Borel}%
  \indexsee{Heine--Borel}{compactness}%
\end{enumerate}


%
它们在经典数学中都是等价的。
让我们看看它们在同伦类型论中表现如何。我们可以使用Dedekind 实数或 Cauchy 实数，因此直接将实数记作~$\R$。我们先回顾几个基本定义。

\indexsee{space!metric}{metric space}
\index{metric space|(}%

\begin{defn} \label{defn:metric-space}
  一个\define{度量空间（metric space）}
  \indexdef{metric space}%
  $(M, d)$ 由一个集合 $M$ 连同一个映射 $d : M \times M \to \R$ 组成，满足对所有的 $x, y, z : M$，
  %
  \begin{align*}
    d(x,y) &\geq 0, &
    d(x,y) &= d(y,x), \\
    d(x,y) &= 0 \Leftrightarrow x = y, &
    d(x,z) &\leq d(x,y) + d(y,z).
  \end{align*}
  %
\end{defn}

\begin{defn} \label{defn:complete-metric-space}
  $M$ 中的一个\define{Cauchy 逼近（Cauchy approximation）}
  \index{Cauchy!approximation}%
  是一个序列 $x : \Qp \to M$，满足
  %
  \begin{equation*}
    \fall{\delta, \epsilon} d(x_\delta, x_\epsilon) < \delta + \epsilon.
  \end{equation*}
  %
  \index{limit!of a Cauchy approximation}%
  一个 Cauchy 逼近 $x : \Qp \to M$ 的\define{极限（limit）}是一个点 $\ell : M$，满足
  %
  \begin{equation*}
    \fall{\epsilon, \theta : \Qp} d(x_\epsilon, \ell) < \epsilon + \theta.
  \end{equation*}
  %
  \indexdef{metric space!complete}%
  \indexdef{complete!metric space}%
  若一个度量空间中每个 Cauchy 逼近都有极限，则称其为\define{完备度量空间（complete metric space）}。
\end{defn}

\begin{defn} \label{defn:total-bounded-metric-space}
  对于一个正有理数 $\epsilon$，度量空间 $(M, d)$ 中的一个\define{$\epsilon$-网（$\epsilon$-net）}
  \indexdef{epsilon-net@$\epsilon$-net}%
  是
  %
  \begin{equation*}
    \sm{n : \N}{x_1, \ldots, x_n : M}
    \fall{y : M} \exis{k \leq n} d(x_k, y) < \epsilon.
  \end{equation*}
  %
  的一个元素。换言之，这是一组有限点列 $x_1, \ldots, x_n$，使得 $M$ 中的每一点仅位于某个 $x_k$ 的 $\epsilon$ 邻域内。

  一个度量空间 $(M, d)$ 是\define{完全有界的（totally bounded）}
  \indexdef{totally bounded metric space}%
  \indexdef{metric space!totally bounded}%
  ，若它对每个 $\epsilon$ 都有 $\epsilon$-网：
  %
  \begin{equation*}
    \prd{\epsilon : \Qp}
    \sm{n : \N}{x_1, \ldots, x_n : M}
    \fall{y : M} \exis{k \leq n} d(x_k, y) < \epsilon.
  \end{equation*}
\end{defn}

\begin{rmk}
  在完全有界性的定义中，我们使用了一种不够严格的记法 $\sm{n : \N}{x_1, \ldots, x_n : M}$。严格来说，我们应该写成 $\sm{x : \lst{M}}$，其中 $\lst{M}$ 是从 \cref{sec:bool-nat} 引入的有限列表\index{type!of lists}的归纳类型。不过，那样会使其余部分的表述稍显繁琐。
\end{rmk}

请注意，在完全有界性的定义中，我们要求的是 $\epsilon$-网的纯粹存在性，而非仅存在性。这样我们就能得到一个函数，为每个 $\epsilon : \Qp$ 指派一个具体的 $\epsilon$-网。这样的函数或许可以称为``完全有界模''。一般而言，在将经典的度量概念移植到同伦类型论时，我们应当节制地使用命题截断，通常是为了避免要求一个从 $\R$ 到 $\Q$ 或 $\N$ 的非常值映射。例如，以下是``一致连续''的``正确''定义。

\begin{defn} \label{defn:uniformly-continuous}
  度量空间上的映射 $f : M \to \R$ 称为\define{一致连续的（uniformly continuous）}
  \indexdef{function!uniformly continuous}%
  \indexdef{uniformly continuous function}%
  ，若满足
  %
  \begin{equation*}
    \prd{\epsilon : \Qp}
    \sm{\delta : \Qp}
    \fall{x, y : M}
    d(x,y) < \delta \Rightarrow |f(x) - f(y)| < \epsilon.
  \end{equation*}
  %
  特别地，一致连续映射有一个\define{一致连续模（modulus of uniform continuity）}\indexdef{modulus!of uniform continuity}，这是一个为每个 $\epsilon$ 指派对应 $\delta$ 的函数。
\end{defn}

现在我们来证明 $[0,1]$ 在第一种意义下是紧的。

\begin{thm} \label{analysis-interval-ctb}
  \index{compactness!metric}%
  \index{interval!open and closed}%
  闭区间 $[0,1]$ 是完备且完全有界的。
\end{thm}

\begin{proof}
  给定 $\epsilon : \Qp$，存在 $k : \N$ 使得 $2/k < \epsilon$，因此我们可以取 $\epsilon$-网 $x_i = i/k$，其中 $i = 0, \ldots, k$。这是一个 $\epsilon$-网，因为对于每个 $y : [0,1]$，仅存在 $i$ 使得 $0 \leq i \leq k$ 且 $(i -
  1)/k < y < (i+1)/k$，从而 $|y - x_i| < 2/k < \epsilon$。

  为证 $[0,1]$ 的完备性，考虑一个 Cauchy 逼近 $x : \Qp \to [0,1]$，并令 $\ell$ 为其在 $\R$ 中的极限。由于 $\max$ 和 $\min$ 是 Lipschitz 映射，由 $r(x) \defeq \max(0, \min(1, x))$ 定义的收缩 $r : \R \to [0,1]$ 与 Cauchy 逼近的极限交换，因此有
  %
  \begin{equation*}
    r(\ell) =
    r (\lim x) =
    \lim (r \circ x) =
    \lim x =
    \ell,
  \end{equation*}
  %
  这意味着 $0 \leq \ell \leq 1$，如所需。
\end{proof}

于是，我们至少在同伦类型论中获得了一个好的紧性概念。遗憾的是，它局限于度量空间，因为完全有界是一个度量概念。我们稍后会考虑另外两个概念，但首先我们证明，在完全有界空间上的一致连续映射具有\define{上确界（supremum）}，
\indexsee{least upper bound}{supremum}%
即一个小于或等于所有其他上界的上界。

\begin{thm} \label{ctb-uniformly-continuous-sup}
  %
  \indexdef{supremum!of uniformly continuous function}%
  在完全有界度量空间 $(M, d)$ 上的一致连续映射 $f : M \to \R$ 有一个上确界 $m : \R$。对于每个 $\epsilon : \Qp$，都存在 $u : M$ 使得 $|m - f(u)| < \epsilon$。
\end{thm}

\begin{proof}
  令 $h : \Qp \to \Qp$ 为 $f$ 的一致连续模。
  我们如下定义一个逼近 $x : \Qp \to \R$：对于任意的 $\epsilon : \Q$，由 $M$ 的完全有界性可以得到一个 $h(\epsilon)$-网 $y_0, \ldots, y_n$。定义
  %
  \begin{equation*}
    x_\epsilon \defeq \max (f(y_0), \ldots, f(y_n)).
  \end{equation*}
  %
  我们断言 $x$ 是一个 Cauchy 逼近。考虑任意的 $\epsilon, \eta : \Q$，此时
  %
  \begin{equation*}
    x_\epsilon \jdeq \max (f(y_0), \ldots, f(y_n))
    \quad\text{and}\quad
    x_\eta \jdeq \max (f(z_0), \ldots, f(z_m))
  \end{equation*}
  %
  其中的 $y_0, \ldots, y_n$ 是某个 $h(\epsilon)$-网，$z_0, \ldots, z_m$ 是某个 $h(\eta)$-网。
  每一个 $z_i$ 仅仅与某个 $y_j$ 的距离在 $h(\epsilon)$ 以内，因此 $|f(z_i) - f(y_j)| <
  \epsilon$，由此可得
  %
  \begin{equation*}
    f(z_i) < \epsilon + f(y_j) \leq \epsilon + x_\epsilon,
  \end{equation*}
  %
  故 $x_\eta < \epsilon + x_\epsilon$。对称地，我们得到 $x_\eta < \eta +
  x_\eta$，因此 $|x_\eta - x_\epsilon| < \eta + \epsilon$。

  我们断言 $m \defeq \lim x$ 是 $f$ 的上确界。要证明对所有 $x : M$ 有 $f(x) \leq m$，只需证明 $\lnot (m < f(x))$。于是，反设 $m <
  f(x)$。存在 $\epsilon : \Qp$ 使得 $m + \epsilon < f(x)$。但此时，仅由参与定义 $x_\epsilon$ 的某个 $y_i$ 便得 $|f(x) - f(y_i)| <
  \epsilon$，因此 $m < f(x) - \epsilon < f(y_i) \leq m$，矛盾。

  最后，我们通过证明 $m$ 满足定理的第二部分来完成证明，因为
  这样 $m$ 就自动成为最小上界。给定任意的 $\epsilon : \Qp$，一方面
  $|m - f(x_{\epsilon/2})| < 3 \epsilon/4$，另一方面，仅由参与定义 $x_{\epsilon/2}$ 的某个 $y_i$ 便得 $|f(x_{\epsilon/2}) - f(y_i)| <
  \epsilon/4$，于是取 $u \defeq y_i$，由三角不等式即得 $|m - f(u)| < \epsilon$。
\end{proof}

现在，如果在 \cref{ctb-uniformly-continuous-sup} 中我们还知道 $M$ 是完备的，我们或许可以寄望于将一致连续的假设弱化为连续，并将结论加强为存在一个使上确界得以达到的点。这些改进的通常证明依赖于如下事实：在一个完备且完全有界的空间中
%


\begin{enumerate}
\item 连续蕴涵一致连续，并且
\item 每个序列都有收敛的子序列。
\end{enumerate}


%
第一条陈述可由 Heine--Borel 紧性容易地推出，第二条则正是 Bolzano--Weierstra\ss{} 紧性。
\index{compactness!Bolzano--Weierstrass@Bolzano--Weierstra\ss{}}%
遗憾的是，这两条都或多或少存在问题。我们先证明 Bolzano--Weierstra\ss{} 紧性蕴含一种被称为\define{有限全知原理}的排中律特例：
\indexsee{axiom!limited principle of omniscience}{limited principle of omniscience}%
\indexdef{limited principle of omniscience}%
即对于每个 $\alpha : \N \to \bool$，
%
\begin{equation} \label{eq:lpo}
  \Parens{\sm{n : \N} \alpha(n) = \btrue} +
  \Parens{\prd{n : \N} \alpha(n) = \bfalse}.
\end{equation}
%
从计算的角度来说，我们并不期望这条原理成立，因为它要求我们判断一个函数是否有无穷多个取值为 $\bfalse$ 的值。

\begin{thm} \label{analysis-bw-lpo}
  %
  区间 $[0,1]$ 的 Bolzano--Weierstra\ss{} 紧性蕴涵有限全知原理。
  \index{compactness!Bolzano--Weierstrass@Bolzano--Weierstra\ss{}}%
\end{thm}

\begin{proof}
  给定任意 $\alpha : \N \to \bool$，定义序列\index{sequence} $x : \N \to [0,1]$ 如下：
  %
  \begin{equation*}
    x_n \defeq
    \begin{cases}
      0 & \text{if $\alpha(k) = \bfalse$ for all $k < n$,}\\
      1 & \text{if $\alpha(k) = \btrue$ for some $k < n$}.
    \end{cases}
  \end{equation*}
  %
  如果 Bolzano--Weierstra\ss{} 性质成立，则存在严格递增的 $f : \N \to
  \N$ 使得 $x \circ f$ 是柯西序列\index{Cauchy!sequence}。对于充分大的 $n:
  \N$，第 $n$ 项 $x_{f(n)}$ 与其极限的距离在 $1/6$ 以内。要么 $x_{f(n)} < 2/3$，
  要么 $x_{f(n)} > 1/3$。若 $x_{f(n)} < 2/3$，则 $x_n$ 收敛于 $0$，于是 $\prd{n : \N}
  \alpha(n) = \bfalse$。若 $x_{f(n)} > 1/3$，则 $x_{f(n)} = 1$，因此 $\sm{n : \N}
  \alpha(n) = \btrue$。
\end{proof}

也许 Bolzano--Weierstra\ss{} 紧性的丧失并不会让我们太过惋惜，但要舍弃 Heine--Borel 紧性则似乎要困难得多——经典数学和 Brouwer（布劳威尔）的直觉主义都接受了它，这一事实便足以印证。由于我们不想过多涉入一般拓扑，我们将采用基本开集来展开讨论。对于 $\R$ 而言，这些就是以有理数为端点的开区间。一个由某类型 $I$ 索引的此类区间族，就是一个映射
%
\begin{equation*}
  \mathcal{F} : I \to \setof{(q, r) : \Q \times \Q | q < r},
\end{equation*}
%
其想法是：满足 $q < r$ 的有理数对 $(q, r)$ 确定了类型 $\setof{ x : \R | q < x < r}$。为稍作方便起见，我们也允许退化区间，因此我们将\define{基本区间族}
\indexdef{family!of basic intervals}%
\indexdef{interval!family of basic}%
取为一个映射
%
\begin{equation*}
  \mathcal{F} : I \to \Q \times \Q.
\end{equation*}
%
更确切地说，一个族是一个依值对 $(I, \mathcal{F})$，而不只是
$\mathcal{F}$。一个\define{有限基本区间族}是由 $\setof{ m :
  \N | m < n}$（对于某个 $n : \N$）索引的族。我们通常用有限列表 $[(q_0, r_0), \ldots,
(q_{n-1}, r_{n-1})]$ 来表示它。最后，$(I, \mathcal{F})$ 的一个\define{有限子族}\indexdef{subfamily, finite, of intervals} 由一个指标列表 $[i_1, \ldots, i_n]$ 给出，该列表随后确定了有限族
$[\mathcal{F}(i_1), \ldots, \mathcal{F}(i_n)]$。

只要我们清楚有序对 $(q, r)$ 与它所对应的区间 $\setof{ x : \R | q < x < r}$ 之间的区别，便可以放心地用同一个记号 $(q, r)$ 来表示二者。
区间之间的\define{交}与\define{包含}关系可以用它们的端点来表示：
%
\symlabel{interval-intersection}
\symlabel{interval-subset}
\begin{align*}
  (q, r) \cap (s, t) &\ \defeq\  (\max(q, s), \min(r, t)),\\
  (q, r) \subseteq (s, t) &\ \defeq\ (q < r \Rightarrow s \leq q < r \leq t).
\end{align*}
%
我们说 $\intfam{i}{I}{(q_i, r_i)}$ \define{（逐点）覆盖} $[a,b]$，若
\indexdef{interval!pointwise cover}%
\indexdef{cover!pointwise}%
\indexdef{pointwise!cover}%
%
\begin{equation} \label{eq:cover-pointwise-truncated}
  \fall{x : [a,b]} \exis{i : I} q_i < x < r_i.
\end{equation}
%
$[0,1]$ 的 \define{Heine–Borel 紧性}
\indexdef{compactness!Heine--Borel}%
是指：对于 $[0,1]$ 的每一个覆盖族，都仅仅存在一个仍能覆盖 $[0,1]$ 的有限子族。

\index{depression}

\begin{thm} \label{classical-Heine-Borel}
  \index{excluded middle}%
  若排中律成立，则 $[0,1]$ 是 Heine–Borel 紧的。
\end{thm}

\begin{proof}
  为推出矛盾，假定有一个族 $\intfam{i}{I}{(a_i,
    b_i)}$ 覆盖了 $[0,1]$，但其任意有限子族都不能覆盖 $[0,1]$。
  我们构造一系列闭区间 $[q_n, r_n]$，它们是嵌套的，长度趋于 $0$，并且其中任何一个都不能被 $\intfam{i}{I}{(a_i, b_i)}$ 的有限子族覆盖。

  我们设 $[q_0, r_0] \defeq [0,1]$。
  假定已经构造好 $[q_n, r_n]$，令 $s
  \defeq (2 q_n + r_n)/3$ 与 $t \defeq (q_n + 2 r_n)/3$。
  区间 $[q_n, t]$ 和 $[s, r_n]$ 都被 $\intfam{i}{I}{(a_i, b_i)}$ 覆盖，但它们不可能同时拥有有限子覆盖，否则 $[q_n, r_n]$ 也会有有限子覆盖。
  要么 $[q_n, t]$ 有有限子覆盖，要么没有。
  若有，则设 $[q_{n+1}, r_{n+1}] \defeq [s, r_n]$；否则设 $[q_{n+1},
  r_{n+1}] \defeq [q_n, t]$。

  序列 $q_0, q_1, \ldots$ 和 $r_0, r_1, \ldots$ 都是柯西序列，它们收敛到同一个点 $x : [0,1]$，该点包含于每一个 $[q_n, r_n]$ 中。
  仅仅存在 $i : I$ 使得 $a_i < x < b_i$。
  由于区间 $[q_n, r_n]$ 的长度趋于零，存在 $n : \N$ 使得 $a_i < q_n \leq x
  \leq r_n < b_i$，但这就意味着 $[q_n, r_n]$ 被单个区间 $(a_i, b_i)$ 所覆盖，同时却又没有有限子覆盖。矛盾。
\end{proof}

没有了排中律，或是缺少一点布劳威尔直觉主义，我们似乎就束手无策了。
然而，$[0,1]$ 的 Heine–Borel 紧性\emph{是能够}在构造性框架下恢复的，而且其方式仍然与经典数学兼容！
为此，我们需要重新审视``覆盖''这一概念。
\eqref{eq:cover-pointwise-truncated} 的问题在于，截断的存在量词允许空间以任意杂乱的方式被覆盖，因此从计算的角度来说，我们便毫无机会``仅仅''提取出一个有限的子覆盖。
去掉截断，我们得到
%
\begin{equation} \label{eq:cover-pointwise}
  \prd{x : [0,1]} \sm{i : I} q_i < x < r_i,
\end{equation}
%
这或许能有所帮助，但它对覆盖的要求实在太高了。
按照这个定义，我们甚至无法证明 $(0,3)$ 和 $(2,5)$ 覆盖了 $[1,4]$，因为那相当于给出一个非常值的映射 $[1,4] \to \bool$，参见
\cref{ex:reals-non-constant-into-Z}。
这里，我们可以从``无点拓扑''
\index{pointfree topology}%
\index{topology!pointfree}%
（即 locale 理论）
\index{locale}%
中汲取经验：覆盖的概念应当基于开集来表达，而不涉及点。
这种对空间的``整体性''观点将使我们能够分析覆盖的概念，进而恢复 Heine–Borel 紧性。
locale 理论使用幂集，
\index{power set}%
我们可以通过假设命题大小重整来获得它；
\index{propositional!resizing}%
不过，我们可以转而借鉴 locale 理论的直谓近亲的思想，
\index{mathematics!predicative}%
它被称为``形式拓扑''。
\index{formal!topology}%

\index{acceptance|(}

假设我们有一个族 $\pairr{I, \mathcal{F}}$ 和一个区间 $(a, b)$。
我们如何在不涉及点的情况下表达``$(a,b)$ 被该族覆盖''这一事实呢？
这里有一个：如果 $(a, b)$ 等于某个 $\mathcal{F}(i)$，那么它被该族覆盖。
还有一个：如果 $(a,b)$ 被另一个族 $(J, \mathcal{G})$ 覆盖，并且每个 $\mathcal{G}(j)$ 又都被 $\pairr{I, \mathcal{F}}$ 覆盖，那么 $(a,b)$ 就被 $\pairr{I, \mathcal{F}}$ 覆盖。
请注意，我们正在列举可以用来\emph{推导}出 $\pairr{I, \mathcal{F}}$ 覆盖 $(a,b)$ 的\emph{规则}。
我们需要找到一组足够好的规则，并将它们转化为一个归纳定义。

\begin{defn} \label{defn:inductive-cover}
  %
  所谓\define{归纳覆盖 $\cover$}
  \indexdef{inductive!cover}%
  \indexdef{cover!inductive}%
  是一个纯粹关系
  %
  \begin{equation*}
    {\cover} : (\Q \times \Q) \to \Parens{\sm{I : \type} (I \to \Q \times \Q)} \to \prop
  \end{equation*}
  %
  它由以下规则归纳地定义，其中 $q, r, s, t$ 是有理数，而
  $\pairr{I, \mathcal{F}}$ 与 $\pairr{J, \mathcal{G}}$ 是基本区间族：
  %
  \begin{enumerate}

  \item \emph{自反性：}
    \index{reflexivity!of inductive cover}%
    $\mathcal{F}(i) \cover \pairr{I, \mathcal{F}}$ 对所有 $i : I$ 成立，

  \item \emph{传递性：}
    \index{transitivity!of inductive cover}%
    若 $(q, r) \cover \pairr{J, \mathcal{G}}$ 且 $\fall{j : J} \mathcal{G}(j) \cover \pairr{I,\mathcal{F}}$，
    则 $(q, r) \cover \pairr{I, \mathcal{F}}$，

  \item \emph{单调性：}
    \index{monotonicity!of inductive cover}%
    若 $(q, r) \subseteq (s, t)$ 且 $(s,t) \cover \pairr{I, \mathcal{F}}$ 则 $(q, r) \cover
    \pairr{I, \mathcal{F}}$，

  \item \emph{局部化：}
    \index{localization of inductive cover}%
    若 $(q, r) \cover (I, \mathcal{F})$ 则 $(q, r) \cap (s, t) \cover
    \intfam{i}{I}{(\mathcal{F}(i) \cap (s, t))}$。

  \item \label{defn:inductive-cover-interval-1}
    若 $q < s < t < r$，则 $(q, r) \cover [(q, t), (s, r)]$，

  \item \label{defn:inductive-cover-interval-2}
    $(q, r) \cover \intfam{u}{\setof{ (s,t) : \Q \times \Q | q < s < t < r}}{u}$。
  \end{enumerate}
\end{defn}

此定义应理解为一个高阶归纳类型，其中所列规则为点构造子，且该类型是 $(-1)$-截断的。
前四条规则具有一般性，直观上应不难理解。
后两条规则则是针对实直线特有的：一条是说，若两个区间有重叠，则它们可以覆盖一个区间；另一条是说，一个区间可以从内部被覆盖。
顺便指出，若 $r \leq q$，则根据最后一条规则，区间 $(q, r)$ 被空族覆盖。

归纳覆盖具有 Heine--Borel 性质，其证明需要下面一个引理。

\begin{lem} \label{reals-formal-topology-locally-compact}
  假设 $q < s < t < r$ 且 $(q, r) \cover \pairr{I, \mathcal{F}}$。那么纯粹存在 $\pairr{I, \mathcal{F}}$ 的一个有限子族，它归纳地覆盖 $(s, t)$。
\end{lem}

\begin{proof}
  我们对 $(q, r) \cover \pairr{I, \mathcal{F}}$ 进行归纳来证明此命题。共有六种情况：
  %
  \begin{enumerate}

  \item 自反性：若 $(q, r) = \mathcal{F}(i)$，则由单调性知，$(s, t)$ 被有限子族 $[\mathcal{F}(i)]$ 覆盖。

  \item 传递性：
    假设 $(q, r) \cover \pairr{J, \mathcal{G}}$ 且 $\fall{j : J} \mathcal{G}(j) \cover
    \pairr{I, \mathcal{F}}$。根据归纳假设，纯粹存在 $[\mathcal{G}(j_1), \ldots, \mathcal{G}(j_n)]$ 覆盖 $(s, t)$。
    再次应用归纳假设，每一个 $\mathcal{G}(j_k)$ 都被 $\pairr{I, \mathcal{F}}$ 的一个有限子族覆盖，我们可以将这些有限子族收集起来，构成一个覆盖 $(s, t)$ 的有限子族。

  \item 单调性：
    若 $(q, r) \subseteq (u, v)$ 且 $(u, v) \cover \pairr{I, \mathcal{F}}$，则可以对 $(u, v) \cover \pairr{I, \mathcal{F}}$ 应用归纳假设，因为此时有 $u <
    s < t < v$。

  \item 局部化：
    假设 $(q', r') \cover \pairr{I, \mathcal{F}}$ 且 $(q, r) = (q', r') \cap (a, b)$。
    由于 $q' < s < t < r'$，根据归纳假设，存在 $(s, t)$ 的一个有限子覆盖 $[\mathcal{F}(i_1), \ldots, \mathcal{F}(i_n)]$。同时我们也知道 $a < s <
    t < b$，因此 $(s, t) = (s, t) \cap (a, b)$ 被
    $[\mathcal{F}(i_1) \cap (a,b), \ldots, \mathcal{F}(i_n) \cap (a,b)]$ 覆盖，而后者正是
    $\intfam{i}{I}{(\mathcal{F}(i) \cap (a, b))}$ 的一个有限子族。

  \item 若对某个 $q < u < v < r$ 有 $(q, r) \cover [(q, v), (u, r)]$，则由单调性可得
    $(s, t) \cover [(q, v), (u, r)]$。

  \item 最后，由自反性即得 $(s, t) \cover \intfam{z}{\setof{ (u,v):\Q \times \Q | q < u < v < r}}{z}$。\qedhere
  \end{enumerate}
\end{proof}

若纯粹存在 $\epsilon : \Qp$ 使得 $(a - \epsilon, b +
\epsilon) \cover \pairr{I, \mathcal{F}}$，则称 \define{$\pairr{I, \mathcal{F}}$ 归纳地覆盖 $[a, b]$}。

\begin{cor} \label{interval-Heine-Borel}
  \index{compactness!Heine-Borel}%
  \index{interval!open and closed}%
  对于归纳覆盖而言，闭区间是 Heine--Borel 紧的。
\end{cor}

\begin{proof}
  假设 $[a, b]$ 被 $\pairr{I, \mathcal{F}}$ 归纳覆盖，即纯粹存在 $\epsilon : \Qp$ 使得 $(a - \epsilon, b + \epsilon) \cover \pairr{I, \mathcal{F}}$。
  由 \cref{reals-formal-topology-locally-compact}，存在
  $(a - \epsilon/2, b + \epsilon/2)$ 的一个有限子覆盖，因而它也是 $[a, b]$ 的有限子覆盖。
\end{proof}

从形式拓扑学\index{topology!formal}的经验可知，归纳覆盖的这些规则足以开展构造性的无点拓扑学。
不过，我们也能自行提供证据表明，它们确实是一个合理的概念。

\begin{thm} \label{inductive-cover-classical}
  \mbox{}
  %
  \begin{enumerate}
  \item 归纳覆盖同时也是逐点覆盖。
  \item 若假设排中律，则逐点覆盖同时也是归纳覆盖。
  \end{enumerate}
\end{thm}

\begin{proof}
  \mbox{}
  %
  \begin{enumerate}

  \item
    考虑一个基本区间族 $\pairr{I, \mathcal{F}}$，其中我们记 $(q_i,
    r_i) \defeq \mathcal{F}(i)$，$(a,b)$ 是一个由 $\pairr{I,
      \mathcal{F}}$ 归纳覆盖的区间，$x$ 满足 $a < x < b$。
    %
    我们通过对 $(a,b) \cover \pairr{I, \mathcal{F}}$ 作归纳来证明，仅仅存在指标 $i : I$ 使得 $q_i < x < r_i$。大部分情形都相当显然，因此我们只展示其中两种。
    若 $(a,b) \cover \pairr{I, \mathcal{F}}$ 由自反性得到，则仅仅存在某个 $i : I$ 使得 $(a,b) = (q_i, r_i)$，从而有 $q_i < x < r_i$。
    若 $(a,b)
    \cover \pairr{I, \mathcal{F}}$ 经 $\intfam{j}{J}{(s_j, t_j)}$ 由传递性得到，则由归纳假设，仅仅存在 $j : J$ 使得 $s_j < x < t_j$；
    继而由于 $(s_j, t_j) \cover \pairr{I, \mathcal{F}}$，再次利用归纳假设可知，仅仅存在某个 $i : I$ 使得 $q_i < x < r_i$。其他情形也同样直截了当。

  \item 假设 $\intfam{i}{I}{(q_i, r_i)}$ 点态覆盖 $(a, b)$。根据
    \cref{defn:inductive-cover} 的 \cref{defn:inductive-cover-interval-2}，只需证明每当
    $a < c < d < b$ 时，$\intfam{i}{I}{(q_i, r_i)}$ 都归纳覆盖 $(c, d)$，因此考虑这样的 $c$ 和 $d$。
    由 \cref{classical-Heine-Borel}
    可知，存在一个有限子族 $[i_1, \ldots, i_n]$ 已经点态覆盖了 $[c,
    d]$，从而也点态覆盖 $(c,d)$。设 $\epsilon : \Qp$ 为关于 $(q_{i_1}, r_{i_1}), \ldots, (q_{i_n}, r_{i_n})$ 的一个Lebesgue 数
    \index{Lebesgue number}
    ，如 \cref{ex:finite-cover-lebesgue-number} 所述。取正数 $k : \N$ 使得 $2 (d - c)/k
    < \min(1, \epsilon)$。对于 $0 \leq i \leq k$，令
    %
    \begin{equation*}
      c_k \defeq ((k - i) c + i d) / k.
    \end{equation*}
    %
    区间 $(c_0, c_2)$, $(c_1, c_3)$, \dots, $(c_{k-2}, c_k)$ 通过反复运用 \cref{defn:inductive-cover} 中的传递性及~\cref{defn:inductive-cover-interval-1}
    归纳覆盖
    $(c,d)$。由于它们的宽度都小于 $\epsilon$，每一个这样的区间都被包含在某个 $(q_i, r_i)$ 之中；
    于是利用传递性和单调性即可断定，$\intfam{i}{I}{(q_i, r_i)}$ 归纳覆盖了 $(c, d)$。 \qedhere
  \end{enumerate}
\end{proof}

前一定理的要旨在于，就经典数学而言，点态覆盖与归纳覆盖之间并无区别。
特别地，由于在同伦类型论中假定排中律是一致的，我们不可能构造出一个不是点态覆盖的归纳覆盖。
或者换一种说法，点态覆盖与归纳覆盖之间的区别，不在于它们覆盖了什么，而在于覆盖的\emph{证明}。

我们大可以像这样继续写下去，再写出一本书来，不过还是就此止步吧；希望以上讨论已为分析学可在同伦类型论中发展这一论断提供了充分的依据。
好奇的读者可以查阅 \cref{ex:mean-value-theorem} 以了解构造性版本的中间值定理。

\index{acceptance|)}

\index{mathematics!classical|)}%
\index{mathematics!constructive|)}%

\section{超现实数}
\label{sec:surreals}

\index{surreal numbers|(}%

本节我们将考察另一个高阶归纳-归纳类型的例子，它汇集了我们之前的许多线索：Conway 的域 $\NO$，即\emph{超现实数（surreal numbers）}~\cite{conway:onag}。
超现实数是（Dedekind ）实数（\cref{sec:dedekind-reals}）与序数（\cref{sec:ordinals}）的自然共同推广。
Conway 在运用排中律和选择公理的经典\index{mathematics!classical}数学框架下，将超现实数定义为一对超现实数的\emph{集合}，记为 $\surr L R$，并满足 $L$ 中的每个元素都严格小于 $R$ 中的每个元素。
这显然像是一个归纳定义，但若视其为归纳定义，会面临三个问题。

首先，该定义需要用到超现实数之间的（严格）序关系，因此这个关系必须与超现实数的类型 $\NO$ 同时定义。
（Conway 通过先定义\emph{游戏（game）}\index{game!Conway}避开了这个问题，游戏类似于超现实数，但省略了 $L$ 与 $R$ 之间的相容性条件。）
正如 Cauchy 实数中的关系 $\closesym$ 一样，这个同时定义\emph{先验地}既可以是归纳-归纳的，也可以是归纳-递归的。出于与为 $\closesym$ 做出该选择时相同的理由，我们在此也选择归纳-归纳定义。

此外，我们将分别（并相互归纳地）定义超现实数的严格序 $<$ 与非严格序 $\le$。
Conway 用 $\le$ 来定义 $<$，这种方式在经典意义下是合理的，但在构造性意义下则不然。
\index{mathematics!constructive}%
再者，一个否定形式的 $<$ 定义，将使其不适合作为高阶归纳类型构造子的假设（参见 \cref{sec:strictly-positive}）。

其次，Conway 说 $\surr L R$ 中的 $L$ 和 $R$ 应该是“超现实数的集合”，但若将其朴素地理解为一个谓词 $\NO\to\prop$，则它不是正的，因此不能用作归纳构造子的输入。
不过，这本身也不是对 Conway 原意的良好类型论翻译，因为在集合论中，超现实数构成一个真类，而集合 $L$ 和 $R$ 是真正的（小）集合，并非 $\NO$ 的任意子类。
在类型论中，这意味着 $\NO$ 将相对于一个宇宙 $\UU$ 来定义，但其本身属于下一个更高的宇宙 $\UU'$，就像序数的集合 $\ord$、基数的集合 $\card$、累积层次 $V$，甚至（在没有命题大小重整时的）Dedekind 实数一样。
\index{propositional!resizing}%
这样一来，我们就要求超现实数的“集合” $L$ 和 $R$ 是 $\UU$-小的，因此用某个以 $\UU$-小类型为指标的\emph{族}来表示它们是很自然的。
（这与我们在 \cref{sec:cumulative-hierarchy} 中对累积层次所做的完全一样。）
也就是说，超现实数的构造子将具有类型
\[ \prd{\LL,\RR:\UU} (\LL\to\NO) \to (\RR\to \NO) \to (\text{some condition}) \to \NO \]
这确实是严格正的。\index{strict!positivity}

最后，在给出 $\NO$ 及其序关系的相互定义后，Conway 规定，若 $x\le y$ 且 $y\le x$，则两个超现实数 $x$ 和 $y$ 是\emph{相等}的。
这自然可以理解为对一个“预超现实数”的集合按等价关系取商。
%(In set-theoretic foundations, one has to us an additional trick to deal with large equivalence classes.)
然而，在没有选择公理的情况下，这样的商构造会遇到与通常构造 Cauchy 实数时的商相同的问题：一对\emph{超现实数}的族将不再能给出一个新的超现实数 $\surr L R$，因为我们不一定能将 $L$ 和 $R$ “提升”为预超现实数的族。
当然，我们可以用与处理 Cauchy 实数相同的方式来解决这个问题，即使用一个\emph{高阶}归纳-归纳定义。

\begin{defn}\label{defn:surreals}
  \define{超现实数}的类型 $\NO$，
  \indexdef{surreal numbers}%
  \indexsee{number!surreal}{surreal numbers}%
  连同关系 $\mathord<:\NO\to\NO\to\type$ 和 $\mathord\le:\NO\to\NO\to\type$，由以下的高阶归纳-归纳方式定义。
  类型 $\NO$ 具有以下构造子。
  \begin{itemize}
  \item 对于任何 $\LL,\RR:\UU$ 以及函数 $\LL\to \NO$ 和 $\RR\to \NO$，若满足条件 $\fall{L:\LL}{R:\RR} x^L<x^R$，则存在一个超现实数 $x$。对于 $L:\LL$ 和 $R:\RR$，我们将这两个函数的值分别记作 $x^L$ 和 $x^R$。
  \item 对于任何满足 $x\le y$ 且 $y\le x$ 的 $x,y:\NO$，我们有 $\noeq(x,y):x=y$。
  \end{itemize}
  我们将第一个构造子的输入称为一个\define{分割}。
  \indexdef{cut!of surreal numbers}%
  如果 $x$ 是由一个分割构造出的超现实数，那么记号 $x^L$ 将隐含 $L:\LL$，类似地 $x^R$ 隐含 $R:\RR$。
  这样我们通常可以避免为索引类型 $\LL$ 和 $\RR$ 显式命名，当讨论涉及多个不同的分割时，这会很方便。
  遵循 Conway 的术语，我们称 $x^L$ 为 $x$ 的一个\emph{左选项（left option）}\indexdef{option of a surreal number}，称 $x^R$ 为一个\emph{右选项（right option）}。

  路径构造子意味着不同的分割可以定义出同一个超现实数。
  因此，谈论任意一个超现实数 $x$ 的左选项或右选项是无意义的，除非我们还知道 $x$ 是由某个特定分割定义的。
  因此，在下文中，我们会比如用“给定一个定义超现实数 $x$ 的分割”这样的说法，以区别于“给定一个超现实数 $x$”。

  关系 $\le$ 具有以下构造子。
  \index{non-strict order}%
  \index{order!non-strict}%
  \begin{itemize}
  \item 给定分别定义两个超现实数 $x$ 和 $y$ 的分割，如果对于所有 $L$ 有 $x^L<y$，并且对于所有 $R$ 有 $x<y^R$，那么 $x\le y$。
  \item 命题截断：
    对于任何 $x,y:\NO$，如果 $p,q:x\le y$，那么 $p=q$。
  \end{itemize}
  关系 $<$ 具有以下构造子。
  \index{strict!order}%
  \index{order!strict}%
  \begin{itemize}
    % Don't technically need x in the first one and y in the second one to be defined by cuts?
  \item 给定分别定义两个超现实数 $x$ 和 $y$ 的分割，如果存在某个 $L$ 使得 $x\le y^L$，那么 $x<y$。
  \item 给定分别定义两个超现实数 $x$ 和 $y$ 的分割，如果存在某个 $R$ 使得 $x^R\le y$，那么 $x<y$。
  \item 命题截断：对于任何 $x,y:\NO$，如果 $p,q:x<y$，那么 $p=q$。
  \end{itemize}
\end{defn}

\noindent
我们将其与 Conway 的定义进行比较：

\begin{itemize}\footnotesize
\item[-] 如果 $L,R$ 是任意两个数集，并且 $L$ 中没有任何成员 $\ge$ $R$ 中的任何成员，那么就存在一个数 $\surr L R$。
  所有的数都是以这种方式构造出来的。
\item[-] $x\ge y$ 当且仅当（没有 $x^R\le y$，且 $x\le$ 没有 $y^L$）。
\item[-] $x=y$ 当且仅当（$x \ge y$ 且 $y\ge x$）。
\item[-] $x>y$ 当且仅当（$x\ge y$ 且 $y\not\ge x$）。
\end{itemize}


若所有对象都是 [超现实] 数而非``游戏''\index{game!Conway}，则在 $x>y$ 的定义中纳入 $x\ge y$ 便非必要。
因此，Conway 的 $<$ 不过就是他的 $\ge$ 的否定，从而他关于 $\surr L R$ 成为超现实数的条件与我们的相同。
对 Conway 的 $\le$ 取否定，再消去双重否定，便得到了我们对 $<$ 的定义；之后我们便可以用 $<$ 来重新表述他的 $\le$，而无需使用否定。

我们立刻就能用许多超现实数来填充 $\NO$。
如同 Conway 那样，我们记
\symlabel{surreal-cut}
\[\surr{x,y,z,\dots}{u,v,w,\dots}\]
表示由分割所定义的超现实数，其中 $\LL\to\NO$ 和 $\RR\to\NO$ 是由 $x,y,z,\dots$ 和 $u,v,w,\dots$ 所描述的类型族。
当然，如果 $\LL$ 或 $\RR$ 是 $\emptyt$，我们就在记号中留空对应的部分。
这里与表示子集的标准记号 $\setof{x:A | P(x)}$ 发生了令人遗憾的冲突，但本节中我们不会使用后者。


\begin{itemize}
\item 我们递归地定义 $\iota_{\nat}:\nat\to\NO$ 如下：
  \begin{align*}
    \iota_{\nat}(0) &\defeq \surr{}{},\\
    \iota_{\nat}(\suc(n)) &\defeq \surr{\iota_{\nat}(n)}{}.
  \end{align*}
  也就是说，$\iota_{\nat}(0)$ 由分割 $\emptyt\to\NO$ 和 $\emptyt\to\NO$ 定义。
  类似地，$\iota_{\nat}(\suc(n))$ 由 $\unit\to\NO$（选出 $\iota_{\nat}(n)$）和 $\emptyt\to\NO$ 定义。
\item 类似地，我们利用符号情形递归原理（\cref{thm:sign-induction}）定义 $\iota_{\Z}:\Z\to\NO$：
  \begin{align*}
    \iota_{\Z}(0) &\defeq \surr{}{},\\
    \iota_{\Z}(n+1) &\defeq \surr{\iota_{\Z}(n)}{} & &\text{$n\ge 0$,}\\
    \iota_{\Z}(n-1) &\defeq \surr{}{\iota_{\Z}(n)} & &\text{$n\le 0$.}
  \end{align*}
\item 所谓\define{2进有理数（dyadic rational）}
  \indexdef{rational numbers!dyadic}%
  \indexsee{dyadic rational}{rational numbers, dyadic}%
  是指一对 $(a,n)$，其中 $a:\Z$ 且 $n:\nat$，并且满足若 $n>0$ 则 $a$ 是奇数。
  我们将其记作 $a/2^n$，并视其为对应的有理数。
  若用 $\Q_D$ 表示所有2进有理数的集合，我们通过对 $n$ 归纳来定义 $\iota_{\Q_D}:\Q_D\to\NO$：
  \begin{align*}
    \iota_{\Q_D}(a/2^0) &\defeq \iota_{\Z}(a),\\
    \iota_{\Q_D}(a/2^n) &\defeq \surr{\iota_{\Q_D}(a/2^n - 1/2^n)}{\iota_{\Q_D}(a/2^n + 1/2^n)},
    \quad \text{for $n>0$.}
  \end{align*}
  这里我们用到了如下事实：如果 $n>0$ 且 $a$ 是奇数，那么 $a/2^n \pm 1/2^n$ 是一个分母小于 $a/2^n$ 的2进有理数。
\item 我们定义 $\iota_{\RD}:\RD\to\NO$,\label{reals-into-surreals} 其中 $\RD$ 是 \cref{sec:dedekind-reals} 中（任意版本的）Dedekind 实数，如下：
  \begin{align*}
    \iota_{\RD}(x) &\defeq
    \surr{q\in\Q_D \text{ such that } q<x}{q\in\Q_D \text{ such that } x<q}.
  \end{align*}
  与前几种情形不同，当我们将2进有理数视为Dedekind 实数时，上式延拓 $\iota_{\Q_D}$ 这一点并非显然。
  这可由简单性定理（\cref{thm:NO-simplicity}）推出。
\item 回顾 \cref{sec:ordinals} 中由关系 $<$ 良序的\emph{序数（ordinal）}\index{ordinal} 类型 $\ord$，其中 $A<B$ 意味着对某个 $b:B$ 有 $A = \ordsl B b$。
  我们通过关于 $\ord$ 的良基递归（\cref{thm:wfrec}）定义 $\iota_{\ord}:\ord\to\NO$\label{ord-into-surreals}：
  \begin{equation*}
    \iota_{\ord}(A) \defeq
    \surr{\iota_{\ord}(\ordsl A a) \text{ for all } a:A}{}.
  \end{equation*}
  同样由简单性定理可知，将 $\iota_{\ord}$ 限制在有限序数上与 $\iota_{\nat}$ 一致。
  （但我们提醒读者，与上面的例子不同，除非将其限制到更小的一类序数上，否则 $\iota_{\ord}$ 并非构造性地单的；参见 \cref{ex:ord-into-surreals,ex:hiit-plump}。）
\item 从 Conway 那里再取几个有趣的例子：
  \begin{align*}
    \omega &\defeq \surr{0,1,2,3,\dots}{} \qquad\text{(also an ordinal)}\\
    -\omega &\defeq \surr{}{\dots,-3,-2,-1,0}\\
    1/\omega &\defeq \textstyle\surr{0}{1,\frac12,\frac14,\frac18,\dots}\\
    \omega-1 &\defeq \surr{0,1,2,3,\dots}{\omega}\\
    \omega/2 &\defeq \surr{0,1,2,3,\dots}{\dots,\omega-2,\omega-1,\omega}.
  \end{align*}
\end{itemize}

在识别由不同分割所给出的超现实数时，下面这个简单的观察很有用。

\begin{thm}[Conway 的简单性定理]\label{thm:NO-simplicity}
  \index{simplicity theorem}%
  \index{theorem!Conway's simplicity}%
  假设 $x$ 和 $z$ 是由分割定义的超现实数，并且满足以下条件。
  \begin{itemize}
  \item 对于所有的 $L$ 和 $R$，有 $x^L < z < x^R$。
  \item 对于 $z$ 的每一个左选项 $z^L$，存在一个左选项 $x^{L'}$ 使得 $z^L\le x^{L'}$。
  \item 对于 $z$ 的每一个右选项 $z^R$，存在一个右选项 $x^{R'}$ 使得 $x^{R'}\le z^R$。
  \end{itemize}
  那么 $x=z$。
\end{thm}

\begin{proof}
  应用 $\NO$ 的路径构造子，我们必须证明 $x\le z$ 且 $z\le x$。
  前者需要证明对所有 $L$ 有 $x^L<z$（这正是我们的假设），并且对所有 $R$ 有 $x<z^R$。
  但根据假设，对于任意 $z^R$，存在一个 $x^{R'}$ 满足 $x^{R'}\le z^R$，从而 $x<z^R$，如所需。
  因此 $x\le z$；$z\le x$ 的证明是对称的。
\end{proof}

\index{induction principle!for surreal numbers}
然而，要想对超现实数进行更深入的讨论，我们需要它们的归纳原理。
适用于 $(\NO,\le,<)$ 的相互归纳原理涉及三个类型族：
\begin{align*}
  A &: \NO\to\type\\
  B &: \prd{x,y:\NO}{a:A(x)}{b:A(y)} (x\le y) \to \type\\
  C &: \prd{x,y:\NO}{a:A(x)}{b:A(y)} (x<y) \to \type.
\end{align*}
与 Cauchy 实数的归纳原理类似，将 $B$ 和 $C$ 视为类型 $A(x)$ 与 $A(y)$ 之间的关系族，会更容易理解。
\symlabel{NO-recursion}
因此，我们将 $B(x,y,a,b,\xi)$ 写作 $(x,a) \ble^\xi (y,b)$，将 $C(x,y,a,b,\xi)$ 写作 $(x,a) \blt^\xi (y,b)$。
类似地，我们通常省略 $\xi$，因为它居于一个命题，因而无关紧要；我们也常常省略 $x$ 和 $y$，直接简写为 $a\ble b$ 或 $a\blt b$。
采用这些记号后，归纳原理的前提如下。

\begin{itemize}
\item 对于定义超现实数 $x$ 的任何分割，以及
  \begin{enumerate}
  \item 对每个 $L$，给定元素 $a^L:A(x^L)$，和
  \item 对每个 $R$，给定元素 $a^R:A(x^R)$，使得
  \item 对所有 $L$ 和 $R$ 有 $(x^L,a^L) \blt (x^R,a^R)$
  \end{enumerate}
  存在一个指定的元素 $f_a:A(x)$。
  我们称这样的数据为定义 $x$ 的分割上的一个\define{依值分割（dependent cut）}
  \indexdef{cut!of surreal numbers!dependent}%
  \indexdef{dependent!cut}。%
\item 对于任意 $x,y:\NO$ 及 $a:A(x)$ 和 $b:A(y)$，若 $x\le y$ 且 $y\le x$，并且还有 $(x,a) \ble (y,b)$
  与 $(y,b) \ble (x,a)$，
  则有 $\dpath{A}{\noeq}{a}{b}$。
\item 给定定义两个超现实数 $x$ 和 $y$ 的分割，以及 $x$ 上的依值分割 $a$ 和 $y$ 上的依值分割 $b$，使得对所有 $L$ 有 $x^L<y$ 且 $(x^L,a^L)\blt (y,f_b)$，
  并且对所有 $R$ 有 $x<y^R$ 且 $(x,f_a) \blt (y^R,b^R)$，
  则有 $(x,f_a) \ble (y,f_b)$。
\item $\ble$ 取值于命题。
\item 给定定义两个超现实数 $x$ 和 $y$ 的分割、$x$ 上的依值分割 $a$ 和 $y$ 上的依值分割 $b$，以及一个满足 $x\le y^{L_0}$ 和 $(x,f_a) \ble (y^{L_0},b^{L_0})$ 的 $L_0$，
  则有 $(x,f_a) \blt (y,f_b)$。
\item 给定定义两个超现实数 $x$ 和 $y$ 的分割、$x$ 上的依值分割 $a$ 和 $y$ 上的依值分割 $b$，以及一个满足 $x^{R_0}\le y$ 和 $(x^{R_0},a^{R_0}),\ble (y,f_b)$ 的 ${R_0}$，
  则有 $(x,f_a) \blt (y,f_b)$。
\item $\blt$ 取值于命题。
\end{itemize}

在这些前提下，我们推导出一个函数 $f:\prd{x:\NO} A(x)$，使得
\begin{align}
  f(x) &\;\jdeq\; f_{f[x]} \label{eq:noind1}\\
  (x\le y) &\;\Rightarrow\; (x,f(x)) \ble (y,f(y)) \notag\\
  (x< y) &\;\Rightarrow\; (x,f(x)) \blt (y,f(y)). \notag
\end{align}
在点构造子的计算规则~\eqref{eq:noind1} 中，$x$ 是由一个分割定义的超现实数，而 $f[x]$ 表示通过应用 $f$ 得到的 $x$ 上的依值分割（并利用 $f$ 将 $<$ 映到 $\blt$ 这一事实）。
通常，我们会使用模式匹配记法，其中 $f$ 在分割 $\surr{x^L}{x^R}$ 上的定义可能会用到符号 $f(x^L)$ 和 $f(x^R)$，以及它们构成一个依值分割的假设。

与 Cauchy 实数的情况类似，通过平凡化 $A$、$\ble$ 和 $\blt$ 中的某些部分，我们可以得到一些特例。
将 $\ble$ 和 $\blt$ 取为常值 \unit，我们就得到了 \define{$\NO$-归纳法（$\NO$-induction）}，为简洁起见，我们只针对命题来陈述它：

\begin{itemize}
\item 给定 $P:\NO\to\prop$，若对每个由分割定义的超现实数 $x$，只要对所有 $L$ 和 $R$ 有 $P(x^L)$ 和 $P(x^R)$ 成立，就有 $P(x)$ 成立，则 $P(x)$ 对所有 $x:\NO$ 成立。
\end{itemize}

这可以与 Conway 的下述评论相对照：

\begin{quote}\footnotesize
  一般来说，当我们希望对所有数 $x$ 确立一个命题 $P(x)$ 时，我们将用归纳法来证明它：由所有 $P(x^L)$ 和 $P(x^R)$ 成立来推导 $P(x)$。
  我们认为``所有数都是以这种方式构造的''这一说法为这一做法的合法性提供了依据。
\end{quote}

利用 $\NO$-归纳法，我们可以证明

\begin{thm}[Conway 定理 0]\label{thm:NO-refl-opt}\
  \index{theorem!Conway's 0}%
  \begin{enumerate}
  \item 对任意 $x:\NO$，有 $x\le x$。\label{item:NO-le-refl}
  \item 对任意由分割定义的 $x:\NO$，对所有 $L$ 和 $R$，有 $x^L <x$ 和 $x<x^R$。\label{item:NO-lt-opt}
  \end{enumerate}
\end{thm}

\begin{proof}
  首先注意，若 $x\le x$，则每当 $x$ 作为某个分割 $y$ 的左选项出现时，由 $<$ 的第一个构造子可得 $x<y$；类似地，每当 $x$ 作为某个分割 $y$ 的右选项出现时，由 $<$ 的第二个构造子可得 $y<x$。
  特别地，有 \ref{item:NO-le-refl}$\Rightarrow$\ref{item:NO-lt-opt}。

  我们通过对 $x$ 作 $\NO$-归纳法来证明 \ref{item:NO-le-refl}。
  因此，假设 $x$ 是由一个分割定义的，且对所有 $L$ 和 $R$ 有 $x^L\le x^L$ 和 $x^R \le x^R$。
  但根据我们上面的观察，这些假设意味着对所有 $L$ 和 $R$ 有 $x^L<x$ 和 $x<x^R$，从而由 $\le$ 的构造子可得 $x\le x$。
\end{proof}

\begin{cor}\label{thm:NO-set}
  $\NO$ 是一个 $0$-类型。
%  (As with $V$, it might be confusing to say that it is a ``set''.)
\end{cor}

\begin{proof}
  纯粹关系 $R(x,y)\defeq (x\le y) \land (y\le x)$ 通过 $\NO$ 的道路构造子蕴含相等性，并且由 \cref{thm:NO-refl-opt}\ref{item:NO-le-refl} 可知它包含对角线。
  因此，\cref{thm:h-set-refrel-in-paths-sets} 适用。
\end{proof}

相比之下，Conway 的定理 1（$\le$ 的传递性）用我们的定义来证明则要困难一些；参见 \cref{thm:NO-unstrict-transitive}。

% Of course, we also have:

% \begin{lem}
%   Every surreal number is merely defined by a cut.
% \end{lem}
% \begin{proof}
%   Obvious by $\NO$-induction.
% \end{proof}

我们还需要联合递归原理，即 \define{$(\NO,\le,<)$-递归（$({\NO,\le,<})$-recursion）}。
为方便起见，我们如下陈述。
假设 $A$ 是一个配备了关系 $\mathord\ble:A\to A\to\prop$ 和 $\mathord\blt:A\to A\to\prop$ 的类型。
那么我们可以通过以下步骤来定义 $f:\NO\to A$。

\begin{enumerate}
\item 对于由分割定义的任意 $x$，假设 $f(x^L)$ 和 $f(x^R)$ 已定义且对所有 $L$ 和 $R$ 满足 $f(x^L)\blt f(x^R)$，我们必须定义 $f(x)$。（我们称此为递归的\emph{主条款}。）\label{item:NO-rec-primary}
\item 证明 $\ble$ 是\emph{反对称的}\index{relation!antisymmetric}：若 $a\ble b$ 且 $b\ble a$，则 $a=b$。
\item 对于由分割定义的 $x,y$，若对所有 $L$ 有 $x^L<y$ 且对所有 $R$ 有 $x<y^R$，并且归纳地假设对所有 $L$ 有 $f(x^L)\blt f(y)$，对所有 $R$ 有 $f(x)\blt f(y^R)$，以及对所有 $L$ 和 $R$ 还有 $f(x^L)\blt f(x^R)$ 和 $f(y^L)\blt f(y^R)$，则我们必须证明 $f(x)\ble f(y)$。
\item 对于由分割定义的 $x,y$ 以及某个满足 $x\le y^{L_0}$ 的 $L_0$，并且归纳地假设 $f(x)\ble f(y^{L_0})$，以及对所有 $L$ 和 $R$ 还有 $f(x^L)\blt f(x^R)$ 和 $f(y^L)\blt f(y^R)$，则我们必须证明 $f(x)\blt f(y)$。
\item 对于由分割定义的 $x,y$ 以及某个满足 $x^{R_0}\le y$ 的 $R_0$，并且归纳地假设 $f(x^{R_0})\ble f(y)$，以及对所有 $L$ 和 $R$ 还有 $f(x^L)\blt f(x^R)$ 和 $f(y^L)\blt f(y^R)$，则我们必须证明 $f(x)\blt f(y)$。\label{item:NO-rec-last}
\end{enumerate}

后三个条款可以更简洁地表述为：我们必须证明 $f$（如第一个条款所定义）将 $\le$ 映为 $\ble$，将 $<$ 映为 $\blt$。
我们将这些性质称为\emph{$f$ 保持不等式}。
此外，在证明 $f$ 保持不等式时，我们可以假设所考虑的 $\le$ 或 $<$ 的具体实例是由其某个构造子生成的，且可使用归纳假设：$f$ 保持了出现在该构造子输入中的所有不等式。

如果我们完成了上述~\ref{item:NO-rec-primary}--\ref{item:NO-rec-last}，就得到 $f:\NO\to A$，它在分割上的计算如~\ref{item:NO-rec-primary} 所指定，且保持所有不等式：
%
\begin{narrowmultline*}
  \fall{x,y:\NO}\Big((x\le y) \to (f(x)\ble f(y))\Big) \land
  \narrowbreak
  \Big((x< y) \to (f(x)\blt f(y))\Big).
\end{narrowmultline*}
%
与 Cauchy 实数的 $(\RC,\closesym)$-递归类似，这个递归原理对于定义 $\NO$ 上的函数至关重要，因为我们无法先在``预超现实数''上定义函数，事后再证明它尊重相等的概念。

\begin{eg}
  我们来定义\emph{取负}函数 $\NO\to\NO$。
  我们应用联合递归原理，取 $A\defeq\NO$，$(x\ble y)\defeq (y\le x)$，以及 $(x\blt y)\defeq (y< x)$。
  显然此 $\ble$ 是反对称的。

  对于定义中的主条款，假设 $x$ 由分割定义，且 $-x^L$ 与 $-x^R$ 已定义并满足对所有 $L$ 和 $R$ 有 $-x^L \blt -x^R$。
  按定义，这意味着对所有 $L$ 和 $R$ 有 $-x^R < -x^L$，因此我们可以通过分割 $\surr{-x^R}{-x^L}$ 来定义 $-x$。
  此记号遵循 Conway 的约定，指的是这样一个分割：其左选项以类型 $\RR$ 为指标——$\RR$ 正是 $x$ 的右选项的指标类型；其右选项以类型 $\LL$ 为指标——$\LL$ 正是 $x$ 的左选项的指标类型；相应的族 $\RR\to\NO$ 和 $\LL\to\NO$ 由 $x$ 的对应族与取负复合而成。

  现在我们必须验证 $f$ 保持不等式。
  \begin{itemize}
  \item 对于 $x\le y$，我们可以假设对所有 $L$ 有 $x^L<y$ 且对所有 $R$ 有 $x < y^R$，并证明 $-y\le -x$。
    但由归纳假设可得 $-y < -x^L$ 且 $-y^R < -x$，再由 $-y$、$-x$ 的定义及 $\le$ 的构造子，即得所需结果。
  \item 对于 $x<y$，在第一种情形中，即它由某个 $x\le y^{L_0}$ 生成时，我们可以归纳地假设 $-y^{L_0} \le -x$，此时由 $<$ 的构造子可得 $-y<-x$。
  \item 类似地，若 $x<y$ 由 $x^{R_0}\le y$ 生成，归纳假设为 $-y \le -x^R$，同样可得 $-y<-x$。
  \end{itemize}
\end{eg}

然而，要在此基础上做进一步的讨论，我们需要像在 \cref{thm:RC-sim-characterization} 中对 Cauchy 实数所做的那样，更明确地刻画关系 $\le$ 和 $<$。
同样地，为了使归纳得以进行，我们必须同时证明这些关系的若干本质性质。

\begin{thm}\label{defn:No-codes}
  存在关系 $\mathord\preceq:\NO\to\NO\to\prop$ 和 $\mathord\prec:\NO\to\NO\to\prop$，使得如果 $x$ 和 $y$ 是由分割定义的超现实数，则
  \begin{align*}
    (x\preceq y) &\defeq
    \big(\fall{L} x^L\prec y\big) \land \big(\fall{R} x\prec y^R\big)\\
    (x\prec y) &\defeq
    \big(\exis{L} x\preceq y^L\big) \lor \big(\exis{R} x^R \preceq y\big).
  \end{align*}
  此外，我们有
  \begin{equation}\label{eq:NO-codes-unstrict}
    (x\prec y) \to (x\preceq y)
  \end{equation}
  以及所有合理的传递性，使 $\prec$ 和 $\preceq$ 成为 $\le$ 和 $<$ 上的``双模''\index{bimodule}：
  \begin{equation}\label{eq:NO-codes-transitivity}
    \begin{array}{c@{\hspace{1cm}}c}
      (x \le y) \to (y\preceq z) \to (x\preceq z) &
      (x \preceq y) \to (y\le z) \to (x\preceq z) \\
      (x \le y) \to (y\prec z) \to (x\prec z) &
      (x \preceq y) \to (y< z) \to (x\prec z) \\
      (x < y) \to (y\preceq z) \to (x\prec z) &
      (x \prec y) \to (y\le z) \to (x\prec z).
  \end{array}
  \end{equation}
\end{thm}

\begin{proof}
  我们通过对 $x, y$ 的双重 $(\NO,\le,<)$-归纳来定义 $\preceq$ 和 $\prec$。
  第一次归纳是一个简单的递归，其值域是 $(\NO\to\prop)\times (\NO\to\prop)$ 的子集 $A$，该子集由一对谓词组成，其中一个蕴含另一个，并且满足“右侧传递性”，即 \eqref{eq:NO-codes-unstrict} 和 \eqref{eq:NO-codes-transitivity} 的右列，其中 $(x\preceq \blank)$ 和 $(x\prec \blank)$ 被替换为两个给定的谓词。
  如同在 \cref{defn:RC-approx} 的证明中，我们将这些谓词视为二元关系的一半来书写，记为 $y\mapsto (\hle y)$ 和 $y\mapsto (\hlt y)$，并用 $\hlname$ 表示这一对关系。
  % The precise definition of $A$ is
  % \begin{align*}
  %   A\defeq \bigg\{ \hlname : (\NO\to\prop)\times (\NO\to\prop) \;\bigg|\;\\
  %   \begin{split}
  %     \fall{y,z:\NO}
  %     &\Big( (\hle y) \to (y\le z) \to (\hle z) \Big)\\
  %     \land\; &\Big( (\hle y) \to (y< z) \to (\hlt z) \Big)\\
  %     \land\; &\Big( (\hlt y) \to (y\le z) \to (\hlt z) \Big)\\
  %     \land\; &\Big( (\hlt y) \to (y< z) \to (\hlt z) \Big) \bigg\}
  %   \end{split}
  % \end{align*}
  我们在 $A$ 上装备以下两个关系：
  \begin{align*}
    (\hlname \ble \hlbname) &\defeq
    \fall{y:\NO} \Big( (\hleb y) \to (\hle y) \Big) \land
    \Big( (\hltb y) \to (\hlt y) \Big),\\
    (\hlname \blt \hlbname) &\defeq
    \fall{y:\NO} \Big( (\hleb y) \to (\hlt y) \Big).
    %\land \Big( (\hltb y) \to (\hlt y) \Big)
  \end{align*}
  注意 $\ble$ 是反对称的，因为如果 $\hlname \ble \hlbname$ 且 $\hlbname \ble \hlname$，那么对于所有 $y$ 有 $(\hleb y) \Leftrightarrow (\hle y)$ 且 $(\hltb y) \Leftrightarrow (\hlt y)$，因此由命题的泛等性和函数外延性可得 $\hlname=\hlbname$。
  此外，说一个函数 $\NO\to A$ 保持不等式，精确地说就是：当将其视为 $\NO$ 上的一对二元关系时，它满足“左侧传递性”（\eqref{eq:NO-codes-transitivity} 的左列）。

  现在处理递归的主情况，我们假设给定一个由分割定义的 $x$，以及对于所有 $L$ 和 $R$ 的关系 $(x^L \prec \blank)$, $(x^R \prec \blank)$, $(x^L \preceq \blank)$, 和 $(x^R \preceq \blank)$。其中严格的关系蕴含非严格的关系，满足右侧传递性，并且使得
  \begin{equation}\label{eq:NO-prec-outer-IH}
    \fall{L,R}{y:\NO}\Big( (x^R\preceq y) \to (x^L \prec y) \Big).
    % \land\Big( (x^R \prec y) \to (x^L \prec y) \Big)
  \end{equation}
  现在我们需要对所有 $y$ 定义 $(x\prec y)$ 和 $(x\preceq y)$。
  这里与 \cref{defn:RC-approx} 相反，我们使用一个嵌套的归纳而非嵌套的递归，以便能够归纳地使用关于不等式 $x^L<x$ 和 $x<x^R$ 的左侧传递性。
  定义 $A':\NO\to\type$，其中 $A'(y)$ 是 $\prop\times\prop$ 的子集，由两个命题组成，记作 $\tle y$ 和 $\tlt y$（对应的元素记为 $\tlname:A'(y)$），使得
  \begin{gather}
    (\tlt y) \to (\tle y)\\
    \fall{L} (\tle y)\to (x^L\prec y) \label{eq:NO-prec-IHL}\\
    \fall{R} (x^R \preceq y) \to (\tlt y) \label{eq:NO-prec-IHR}.
  \end{gather}
  使用类似于 $\ble$ 和 $\blt$ 的记号，我们在 $A'$ 上装备两个关系，对于 $\tlname:A'(y)$ 和 $\tlbname:A'(z)$ 定义如下：
  \begin{align*}
    (\tlname \bble \tlbname) &\defeq
    \Big((\tle y) \to (\tleb z)\Big) \land \Big((\tlt y) \to (\tltb z)\Big)\\
    (\tlname \bblt \tlbname) &\defeq
    \Big((\tle y) \to (\tltb z)\Big). % \land \Big(\tlt \to \tltb\Big).
  \end{align*}
  % (These are the type families $B$ and $C$ in the general induction principle.)
  同样，$\bble$ 在适当的意义下显然是反对称的。
  此外，一个能保持不等式的函数 $\prd{y:\NO} A'(y)$ 恰好是这样一对谓词：其中一个蕴含另一个，满足右侧传递性，并且关于不等式 $x^L<x$ 和 $x<x^R$ 满足左侧传递性。
  因此，这个内层归纳将提供我们完成外层递归主情况所需的内容。

  对于内层归纳的主情况，我们还假设给定一个由分割定义的 $y$，以及对于所有 $L$ 和 $R$ 的性质 $(x\prec y^L)$, $(x\prec y^R)$, $(x\preceq y^L)$, 和 $(x\preceq y^R)$。其中严格的关系蕴含非严格的关系，关于 $x^L<x$ 和 $x<x^R$ 满足左侧传递性，关于 $y^L<y^R$ 满足右侧传递性。
  % \begin{equation}
  %   \fall{L,R}\Big((x \preceq y^L) \to (x \prec y^R)\Big) % \land \Big((x \prec y^L) \to (x\prec y^R)\Big).
  %   \label{eq:NO-prec-inner-IH}
  % \end{equation}
  我们现在可以给出定理陈述中指定的定义：
  \begin{align}
    (x\preceq y) &\defeq
    (\fall{L} x^L\prec y) \land (\fall{R} x\prec y^R), \label{eq:NO-preceq-def}\\
    (x\prec y) &\defeq
    (\exis{L} x\preceq y^L) \lor (\exis{R} x^R \preceq y).\label{eq:NO-prec-def}
  \end{align}
  为了让这一定义出 $A'(y)$ 的一个元素，我们必须首先证明 $(x\prec y) \to (x\preceq y)$。
  假设 $x\prec y$ 有两种情况。
  一方面，如果存在 $L_0$ 使得 $x\preceq y^{L_0}$，那么通过关于 $y^{L_0}<y^R$ 的右侧传递性，我们得到对于所有 $R$ 有 $x\prec y^R$。
  此外，通过关于 $x^L<x$ 的左侧传递性，对于任何 $L$ 我们有 $x^L \prec y^{L_0}$，从而由右侧传递性得到 $x^L\prec y$。
  因此，$x\preceq y$。

  另一方面，如果存在 $R_0$ 使得 $x^{R_0}\preceq y$，那么由 \eqref{eq:NO-prec-outer-IH}，我们得到对于所有 $L$ 有 $x^L \prec y$。
  并且通过关于 $x<x^{R_0}$ 和 $y<y^R$ 的左侧与右侧传递性，对于任何 $R$ 我们有 $x\prec y^R$。
  因此，$x\preceq y$。

  我们还需要证明这些定义关于 $x^L<x$ 和 $x<x^R$ 满足左侧传递性。
  但如果 $x\preceq y$，根据定义对于所有 $L$ 有 $x^L\prec y$；而如果 $x^R\preceq y$，同样根据定义有 $x\prec y$。

  因此，\eqref{eq:NO-preceq-def} 和 \eqref{eq:NO-prec-def} 确实定义了 $A'(y)$ 的一个元素。
  我们现在需要验证，这个定义作为进入 $A'$ 的依值函数是保持不等式的，即这些关系满足右侧传递性。
  记住在每种情况下，我们可以归纳地假设它们对于不等式构造器中产生的所有不等式都满足右侧传递性。
  \begin{itemize}
  \item 假设 $x\preceq y$ 且 $y\le z$，后者源于对于所有 $L$ 和 $R$ 有 $y^L<z$ 和 $y<z^R$。
    那么将内层递归的归纳假设应用于 $y<z^R$，得到对于任何 $R$ 有 $x\prec z^R$。
    此外，根据定义 $x\preceq y$ 意味着对于任何 $L$ 有 $x^L \prec y$，因此通过外层递归的归纳假设我们得到 $x^L \prec z$。
    于是，$x\preceq z$。
  \item 假设 $x\preceq y$ 且 $y<z$。
    首先，假设 $y<z$ 源于 $y\le z^{L_0}$。
    那么将内层归纳假设应用于 $y\le z^{L_0}$，得到 $x \preceq z^{L_0}$，因此 $x\prec z$。

    其次，假设 $y<z$ 源于 $y^{R_0}\le z$。
    那么根据定义，$x\preceq y$ 意味着 $x\prec y^{R_0}$，然后将内层归纳假设应用于 $y^{R_0}\le z$，得到 $x\prec z$。
  \item 假设 $x\prec y$ 且 $y\le z$，后者源于对于所有 $L$ 和 $R$ 有 $y^L<z$ 和 $y<z^R$。
    根据定义，$x\prec y$ 意味着仅仅存在 $R_0$ 使得 $x^{R_0}\preceq y$ 或存在 $L_0$ 使得 $x\preceq y^{L_0}$。
    如果 $x^{R_0}\preceq y$，那么外层归纳假设给出 $x^{R_0}\preceq z$，因此 $x\prec z$。
    如果 $x\preceq y^{L_0}$，那么将内层归纳假设应用于 $y^{L_0}<z$（这由 $y\le z$ 的构造子保证）得到 $x\prec z$。
  % \item Suppose $x\prec y$ and $y<z$.
  %   First, suppose $y<z$ arises from $y\le z^{L_0}$.
  %   Then the inner inductive hypothesis for $y\le z^{L_0}$ yields $x\prec z^{L_0}$, hence $x\preceq z^{L_0}$; thus $x\prec z$.

  %   Second, suppose $y<z$ arises from $y^{R_0}\le z$.
  %   Then by definition, $x\prec y$ implies there merely exists $R_1$ with $x^{R_1}\preceq y$ or $L_1$ with $x\preceq y^{L_1}$.
  %   If $x^{R_1}\preceq y$, then the outer inductive hypothesis implies $x^{R_1}\prec z$, hence $x^{R_1}\preceq z$, and thus $x\prec z$.
  %   And if $x\preceq y^{L_1}$, then the inner inductive hypothesis applied to $y^{L_1}<y^{R_0}$ (which comes from $y$ being defined as a cut) and $y^{R_0}\le z$ yields $x\prec z$.
  \end{itemize}
  这就完成了内层归纳。
  因此，对于任何由分割定义的 $x$，我们由 \eqref{eq:NO-preceq-def} 和 \eqref{eq:NO-prec-def} 定义了 $(x\prec \blank)$ 和 $(x\preceq \blank)$，并且它们满足右侧传递性。

  为了完成外层递归，我们需要验证这些定义满足左侧传递性。
  在对 z 进行 $\NO$-归纳之后，我们最终会得到与上面为右侧传递性描述的三种情况本质上相同的三种情况。
  因此，我们省略它们。
\end{proof}

\begin{thm}\label{thm:NO-encode-decode}
  对于任意 $x,y:\NO$，我们有 $(x<y)=(x\prec y)$ 与 $(x\le y)=(x\preceq y)$。
\end{thm}

\begin{proof}
  从左到右，我们使用 $(\NO,\le,<)$-归纳，其中 $A(x)\defeq\unit$，并用 $\preceq$ 和 $\prec$ 来提供关系 $\ble$ 和 $\blt$。
  在所有构造子情形中，$x$ 和 $y$ 都由分割定义，因此 $\preceq$ 和 $\prec$ 的定义可以直接计算，归纳前提适用。

  从右到左，我们使用 $\NO$-归纳，假设 $x$ 和 $y$ 由分割定义。
  但此时 $\preceq$ 和 $\prec$ 的定义以及归纳前提，恰好提供了 $\le$ 和 $<$ 的相关构造子所需的数据。
  % From right to left, we first prove by $\NO$-induction on $x$ that for any $y:\NO$ we have $(x\prec y) \to (x<y)$ and $(x\preceq y) \to (x\le y)$.
  % Thus, we assume this to be true for all $x^L$ and $x^R$ in a cut, and show it for the resulting $x:\NO$.
  % Next, we prove by $\NO$-induction on $y$ that $(x\prec y) \to (x<y)$ and $(x\preceq y) \to (x\le y)$, hence we assume it to be true for all $y^L$ and $y^R$ in a cut, and show it for the resulting $y:\NO$.
  % Now since $x$ and $y$ are both defined by cuts, $x\preceq y$ means that $x^L\prec y$ and $x\prec y^R$ for all $L$ and $R$.
  % By the inductive hypotheses, this gives $x^L<y$ and $x<y^R$, hence $x\le y$ by the constructor of $\le$.
  % Similarly, $x\prec y$ yields merely an $R_0$ with $x^{R_0}\preceq y$ or an $L_0$ with $x\preceq y^{L_0}$.
  % Hence merely $x^{R_0}\le y$ or $x\le y^{L_0}$ by the inductive hypothesis, so $x<y$ by a constructor.
\end{proof}

\begin{cor}\label{thm:NO-unstrict-transitive}
  $\NO$ 上的关系 $\le$ 和 $<$ 满足
  \[ \fall{x,y:\NO} (x<y) \to (x\le y) \]
  且是传递的：
  \index{transitivity!of . for surreals@of $<$ for surreals}
  \index{transitivity!of . for surreals@of $\leq$ for surreals}
  \begin{gather*}
    (x\le y) \to (y\le z) \to (x\le z)\\
    (x\le y) \to (y< z) \to (x< z)\\
    (x< y) \to (y\le z) \to (x< z).
  \end{gather*}
\end{cor}

与 Cauchy 实数类似，在定义 $\NO$ 上所有运算时，联合 $(\NO,\le,<)$-递归原理仍然是必不可少的。

\begin{eg}\label{eg:surreal-addition}
\index{addition!of surreal numbers}%
我们通过先对第一个变元递归、再对第二个变元归纳，来定义 $\mathord+:\NO\to\NO\to\NO$。
对于外层递归，我们取余域为 $\NO\to\NO$ 的一个子集，该子集由那些对所有 $x,y$ 满足 $(x<y) \to (g(x)<g(y))$ 和 $(x\le y) \to (g(x)\le g(y))$ 的函数 $g$ 组成。
对于这样的 $g,h$，我们定义 $(g\ble h)\defeq \fall{x:\NO} g(x)\le h(x)$ 与 $(g\blt h)\defeq \fall{x:\NO} g(x)< h(x)$。
显然 $\ble$ 是反对称的。

对于递归的主条款，我们假设 $x$ 由分割定义，函数 $(x^L+\blank)$ 和 $(x^R+\blank)$ 已定义、保序，且满足 $x^L+y<x^R+y$，然后定义 $(x+\blank)$。
如同在 \cref{defn:No-codes} 中那样，我们不进行内层递归，而是向类型族 $A:\NO\to\type$ 进行内层归纳，其中 $A(y)$ 是 $\NO$ 中满足以下条件的 $z$ 构成的子集：对所有 $x^L$ 有 $x^L + y < z$，且对所有 $x^R$ 有 $x^R + y > z$。
我们为 $A$ 赋予从 $\NO$ 诱导的关系 $\le$ 和 $<$，因此反对称性是显然的。
对于内层归纳的主条款，我们进一步假设 $y$ 也由分割定义，且每个 $x+y^L$ 和 $x+y^R$ 已定义并满足 $x^L+y^L < x+y^L$、$x^L+y^R < x+y^R$、$x+y^L < x^R + y^L$ 以及 $x+y^R < x^R+y^R$（这些来自对 $A(y)$ 元素施加的额外条件），同时也满足 $x+y^L < x+y^R$（因为 $A(y)$ 中的元素 $x+y^L$ 和 $x+y^R$ 构成了一个依值分割）。
现在我们给出 Conway 的定义：
\[ x+y \defeq \surr{x^L+y, x+y^L}{x^R+y,x+y^R}. \]
换句话说，$x+y$ 的左选项是所有形如 $x^L+y$（对某个左选项 $x^L$）或 $x+y^L$（对某个左选项 $y^L$）的数。
我们必须证明这些左选项中的每一个都小于每一个右选项：
\begin{itemize}
\item $x^L+y < x^R+y$，由外层归纳前提。
\item $x^L+y < x^L + y^R < x + y^R$，第一个因为 $(x^L+\blank)$ 保序，第二个因为 $x+y^R : A(y^R)$。
\item $x+y^L < x^R+ y^L < x^R + y$，第一个因为 $x+y^L : A(y^L)$，第二个因为 $(x^R+\blank)$ 保序。
\item $x+y^L < x+y^R$，由内层归纳前提（具体来说，由我们有一个依值分割这一事实）。
\end{itemize}
我们还必须证明如此定义的 $x+y$ 属于 $A(y)$，即 $x^L + y < x+y$ 且 $x+y < x^R + y$；但这由 \cref{thm:NO-refl-opt}\ref{item:NO-lt-opt} 即得。

接下来我们必须验证 $(x+\blank)$ 的定义保序：
\begin{itemize}
\item 若 $y\le z$ 是由已知对所有 $L$ 和 $R$ 有 $y^L<z$ 且 $y<z^R$ 而得到的，则内层归纳前提给出 $x+y^L<x+z$ 和 $x+y < x+z^R$，而外层归纳前提给出 $x^L+y \le x^L+z$ 和 $x^R+ y \le x^R+z$。
  此外，由于 $x^R+y$ 依定义是 $x+y$ 的一个右选项，我们有 $x+y < x^R+y$。
  类似地，$x^L+z$ 是 $x+z$ 的一个左选项，因此 $x^L+z < x+z$。
  于是，利用传递性，我们得到 $x^L+y < x+z$ 和 $x+y < x^R+z$；因此由 $\le$ 的构造子可得 $x+y \le x+z$。
\item 若 $y<z$ 来自某个 $L_0$ 使得 $y\le z^{L_0}$，则由归纳有 $x+y \le x+z^{L_0}$，因此 $x+y<x+z$，因为 $x+z^{L_0}$ 是 $x+z$ 的一个右选项。
\item 类似地，若 $y<z$ 来自 $y^{R_0}\le z$，则 $x+y<x+z$，因为 $x+y^{R_0}\le x+z$。
\end{itemize}
这就完成了内层归纳。
对于外层递归，我们必须验证 $+$ 在左侧也保序。
进行一次 $\NO$-归纳后，论证过程完全相同。
\end{eg}

\index{acceptance|(}%
\index{mathematics!formalized}%
在~\cite{conway:onag} 第零部分的附录中，Conway 讨论了如何在 ZFC 集合论中形式化超现实数：沿序数迭代，并对每个等价类取秩最低的代表元所构成的集合；或者用``符号展开''来表示数。
他随后指出

\begin{quote}\footnotesize
  这些构造出奇复杂的性质，更多地揭示了 ZF 框架内形式化的本质，而非我们的数系本身……
\end{quote}

接着，他主张建立一种关于``容许的构造种类''的一般理论，其中应当包括：

\begin{enumerate}\footnotesize
\item 对象可以以任何具有合理构造性的方式从先前的对象中创建。\label{item:conway1}
\item 所创建对象之间的相等关系可以是任意的等价关系。\label{item:conway2}
\end{enumerate}

\noindent
条件~\ref{item:conway1} 可以很自然地理解为对\emph{归纳定义（inductive definition）}一般原则的证成，例如在 \cref{sec:strictly-positive,sec:generalizations} 中阐述的那些。
特别地，构造子的严格正性（strict positivity）条件，可以看作是``具有合理构造性''这一要求的形式化表述。
条件~\ref{item:conway2} 则提示我们，应将此扩展到各种\emph{高阶}归纳定义（higher inductive definitions），在其中我们可以引入道路构造子，以任意合理的方式规定对象之间的相等。
例如，Conway 在下一段中写道：

\begin{quote}\footnotesize
  ……我们也可以，比方说，自由地创造一个新对象 $(x,y)$，并将其称为 $x$ 和 $y$ 的有序对。
  我们还可以创造一个与 $(x,y)$ 不同但与之共存的有序对 $[x,y]$……
  而如果我们想把 $(x,y)$ 做成无序对，我们可以通过等价关系来定义相等，即：当且仅当 $x=z,y=t$ \emph{或} $x=t,y=z$ 时，有 $(x,y)=(z,t)$。
\end{quote}

通过特定形式的构造子来引入具有新名称的新对象——这种自由，恰恰是我们在归纳定义理论中所拥有的。
正如我们在 \cref{sec:appetizer-univalence} 中处理自然数的两个副本 $\nat$ 和 $\nat'$ 时一样，如果我们写下与 Cartesian 积类型 $A\times B$ 完全相同的定义，我们会得到一个不同的积类型 $A\times' B$，其典范元素可自由地写作 $[x,y]$。
并且，我们可以通过添加适当的道路构造子，使其成为无序对的类型。% (and perhaps 0-truncating).

诚然，Conway 的重点并非专门抱怨 ZF，而是一并反对所有的基础理论：

\begin{quote}\footnotesize
  ……本提议并非要用某个特定理论来替代 ZF……
  相反，提议的是：我们赋予自己自由，去创造任意此类数学理论，但要证明一个元定理，它一劳永逸地确保任何此类理论都可以用任何标准的基础理论来形式化。
\end{quote}

或许有人会回应说，事实上，泛等基础并非 Conway 心目中的那些``标准基础理论''之一，而是这样一种\emph{元理论}：在其中我们可以表达创造新理论的能力，也可以据此证明 Conway 的元定理。
例如，超现实数就是 Conway 心目中的``数学理论''之一，而我们已看到，它们可以在泛等基础内部构造出来并证明其合理性。
类似地，Conway 早前曾指出：

\begin{quote}\footnotesize
  ……集合论就是这样一种理论，集合通过对应于通常公理的构造过程从更早的集合构造出来，而相等关系即具有相同成员。
\end{quote}

这一描述与 \cref{sec:cumulative-hierarchy} 中对集合论累积层次的高阶归纳构造紧密契合。
Conway 的元定理便对应于我们多次提及的那个事实：我们可以在 ZFC 内部构造一个泛等基础的模型（这超出了本书的范围）。

然而，泛等基础本身是如此丰富和强大，若仅仅将它降格为一种用于构造类集合理论的元理论，那将是愚蠢的。
我们看到，即使在集合（0-类型）的层面，泛等基础中的高阶归纳类型也能通过万有性质直接产生对象的构造（\cref{sec:free-algebras}），例如 Cauchy 完备化的构造性理论（\cref{sec:cauchy-reals}）。
但最重要的是，在基础系统中直接为同伦论和范畴论建模的潜力（\cref{cha:homotopy,cha:category-theory}），赋予了泛等基础任何集合论基础都无法比拟的优势。
\index{acceptance|)}%

\index{surreal numbers|)}%

\sectionNotes

在Dedekind 实数上定义代数运算，尤其是乘法，既有些棘手又相当繁琐。
有好几种方法可以建立算术运算：每种都有其优势，但它们似乎都需要一些技术性工作。
例如，Richman~\cite{Richman:reals} 先在正分割上定义乘法，然后再代数地延拓到所有Dedekind 分割；
而 Conway~\cite{conway:onag} 则观察到，超现实数的乘法定义同样可以很好地适用于Dedekind 实数。

我们对Dedekind 实数的处理借鉴了~\cite{BauerTaylor09} 中的许多想法，该文献在 Abstract Stone Duality 的背景下构造了Dedekind 实数。
\index{Abstract Stone Duality}%
这是一种（受限的）简单类型 $\lambda$-演算，带有一个特定的对象 $\Sigma$，该对象分类开集，并由对偶性也分类闭集。
在~\cite{BauerTaylor09} 中也可以找到算术运算基本性质的详细证明。

$\RC$ 是最小的 Cauchy 完备的 Archimedean 的有序域，这一事实（已在 \cref{RC-initial-Cauchy-complete} 中证明）表明，我们的 Cauchy 实数很可能与 Escard{\'o}--Simpson 实数~\cite{EscardoSimpson:01} 相一致。
\index{real numbers!Escardo--Simpson@Escard\'o--Simpson}%
验证这是否确实如此\index{open!problem}，将是一个有趣的问题。
Escard{\'o}--Simpson 实数的概念，更准确地说，是相应的闭区间概念，之所以有趣，是因为它可以在任何具有有限积的范畴中陈述。

在增加了``正则扩张公理''的构造性集合论中，人们也可以尝试通过超限迭代、对 Cauchy 序列的极限取闭包来定义 Cauchy 完备化。
验证这种构造是否与我们的构造一致，同样是一个有趣的问题。

构造性数学中一个广为流传的结论是：Cauchy 实数与Dedekind 实数的一致需要依赖选择公理。
但较少人知道的是，可数选择公理就足够了。
回想一下，\define{依赖选择公理}
\indexdef{axiom!of choice!dependent}%
\index{axiom!of choice!countable}%
\index{total!relation}%
说的是：对于 $A$ 上的一个全关系 $R$（这里全关系是指 $\fall{x : A} \exis{y : A} R(x,y)$），以及任意 $a : A$，都仅存在 $f : \N \to A$ 使得 $f(0) = a$，并且对于所有 $n : \N$ 都有 $R(f(n), f(n+1))$。
我们的 \cref{when-reals-coincide} 使用了典型的技巧，将一个依赖选择公理的应用转化为使用可数选择公理。
具体来说，我们先用一次可数选择公理，提前做好所有可能出现的那些选择，然后利用这个选择函数来避免做出依赖性的选择。

各种紧性概念在构造性背景下的复杂关系在 \cite{bridges2002compactness} 中有所讨论。Palmgren~\cite{Palmgren:FT} 则对逐点分析和无点拓扑进行了很好的比较。

超现实数由~\cite{conway:onag} 定义，使用了一种归纳定义，但并未在任何基础系统下明确证明其合理性。
正因如此，一些后来的作者倾向于使用符号展开式或其他更明确的表示法，以便更明显地编码到集合论中。
在类型论中表示它们的想法最初由 Hancock 考虑，而 Setzer 和 Forsberg~\cite{forsbergfinite} 指出，超现实数及其不等式关系 $<$ 和 $\le$ 自然地构成了一个归纳-归纳定义。
这里介绍的\emph{高阶}归纳-归纳版本，内置了超现实数的正确相等概念，是新的。

\sectionExercises

\begin{ex}\label{ex:alt-dedekind-reals}
  通过先定义平方再利用 \cref{mult-from-square}，给出Dedekind 实数的另一种定义。验证由此得到的是一个交换环。
\end{ex}

\begin{ex} \label{ex:RD-extended-reals}
  %
  假设我们去掉 \cref{defn:dedekind-reals} 中的有界性条件 \ref{defn:dedekind-reals-inhabited}。
  那么我们就得到了 \define{扩充实数}
  \indexdef{real numbers!extended}%
  \indexdef{extended real numbers}%
  它们包含 $-\infty \defeq (\emptyt, \Q)$ 和 $\infty \defeq (\Q, \emptyt)$。哪些分割上的算术运算定义对扩充实数仍有意义？我们会得到什么样的代数结构？
\end{ex}

\begin{ex} \label{ex:RD-lower-cuts}
  %
  通过考虑单侧分割，我们分别得到\define{下Dedekind 实数（lower Dedekind reals）}和\define{上Dedekind 实数（upper Dedekind reals）}。
  \indexdef{real numbers!Dedekind!upper}%
  \indexdef{real numbers!Dedekind!lower}%
  \indexdef{lower Dedekind reals}%
  \indexdef{upper Dedekind reals}%
  \index{cut!Dedekind}%
  例如，一个下实数可由一个谓词 $L : \Q \to \Omega$ 给出，它满足以下条件：
  %
  \begin{enumerate}
  \item \emph{有实例的：} $\exis{q : \Q} L(q)$；
  \item \emph{舍入的：} $L(q) = \exis{r : \Q} q < r \land L(r)$。
    \index{rounded!Dedekind cut}
  \end{enumerate}
  %
  （我们也可以要求 $\exis{r : \Q} \lnot L(r)$，以排除分割 $\infty \defeq
  \Q$。）在下实数上可以定义哪些算术运算？特别是，加法逆元的情况如何？
\end{ex}

\begin{ex} \label{ex:RD-interval-arithmetic}
  %
  \index{interval!arithmetic}%
  假设我们去掉 \cref{defn:dedekind-reals} 中的位于性条件。
  那么我们就得到了\define{区间域（interval domain）}
  \indexdef{interval!domain}%
  $\mathbb{I}$，因为此时分割允许存在``间隙''，而这些间隙就是区间。通过
  %
  \begin{narrowmultline*}
    ((L, U) \sqsubseteq (L', U'))
    \defeq \narrowbreak
    (\fall{q : \Q} L(q) \Rightarrow L'(q)) \land
    (\fall{q : \Q} U(q) \Rightarrow U'(q)).
  \end{narrowmultline*}
  %
  定义 $\mathbb{I}$ 上的偏序 $\sqsubseteq$。
  $\mathbb{I}$ 中的极大元有哪些？定义``端点''运算，它将区间域中的一个元素映为其下端点与上端点。这些端点是实数、下实数还是上实数（参见
  \cref{ex:RD-lower-cuts}）？分割上的哪些算术运算定义对区间域仍然有意义？
\end{ex}

\begin{ex} \label{ex:RD-lt-vs-le}
  证明对于所有 $x, y : \RD$，
  %
  \begin{equation*}
    \lnot (x < y) \Rightarrow y \leq x
  \end{equation*}
  %
  以及
  %
  \begin{equation*}
    \eqv{(x \leq y)}{\Parens{\prd{\epsilon : \Qp} x < y + \epsilon}}.
  \end{equation*}
  %
  成立。
  $\lnot (x \leq y)$ 是否蕴含 $y < x$？
\end{ex}

\begin{ex} \label{ex:reals-non-constant-into-Z}
  \mbox{}
  %
  \begin{enumerate}
  \item
    假设排中律，构造一个从 $\RD$ 到 $\Z$ 的非常值映射。
  \item
    假设 $f : \RD \to \Z$ 是一个映射，满足 $f(0) = 0$ 且对所有 $x >
    0$ 有 $f(x) \neq 0$。试由此推导出有限全知原理~\eqref{eq:lpo}。
\index{limited principle of omniscience}%
  \end{enumerate}
\end{ex}

\begin{ex} \label{ex:traditional-archimedean}
  \index{ordered field!archimedean}%
  证明在有序域 $F$ 中，$\Q$ 的稠密性与传统的 Archimedean 公理是等价的：
  %
  \begin{equation*}
    (\fall{x, y : F} x < y \Rightarrow \exis{q : \Q} x < q < y)
    \Leftrightarrow
    (\fall{x : F} \exis{k : \Z} x < k).
  \end{equation*}
\end{ex}

\begin{ex} \label{RC-Lipschitz-on-interval} 假设 $a, b : \Q$ 且 $f : \setof{ q : \Q |
    a \leq q \leq b } \to \RC$ 是 Lipschitz 的，常数为~$L$。证明存在 $f$ 的唯一延拓 $\bar{f} : [a,b] \to \RC$，且该延拓是 Lipschitz 的，常数为~$L$。提示：不必为闭区间重做 \cref{RC-extend-Q-Lipschitz}，可以注意到存在一个收缩 $r : \RC \to [-n,n]$，然后将 \cref{RC-extend-Q-Lipschitz} 应用于 $f \circ r$ 即可。
\end{ex}

\begin{ex} \label{ex:metric-completion}
  \index{completion!of a metric space}%
  将 $\RC$ 的构造推广，以构造任意度量空间的 Cauchy 完备化。首先，思考哪一种实数概念作为度量空间距离函数\index{distance}的陪域最为自然。这有影响吗？接下来，详细完成两种构造：
  %
  \begin{enumerate}
  \item 遵循 Cauchy 实数的构造，将度量空间的完备化定义为一个在柯西序列极限下封闭的归纳-归纳类型。\index{Cauchy!sequence}
  \item 使用由 Lawvere（劳维尔）~\cite{lawvere:metric-spaces}\index{Lawvere} 和 Richman（里奇曼）~\cite{Richman00thefundamental} 提出的构造，其中度量空间 $(M, d)$ 的完备化由\define{位置（locations）}的类型给出。
    \indexdef{location}%
    一个位置是一个函数 $f : M \to \R$，使得
    %
    \begin{enumerate}
    \item $f(x) \geq |f(y) - d(x,y)|$ 对所有 $x, y : M$ 成立；
    \item $\inf_{x \in M} f(x) = 0$，即 $\fall{\epsilon : \Qp} \exis{x : M} |f(x)| < \epsilon$ 与 $\fall{x : M} f(x) \geq 0$。
    \end{enumerate}
    %
    这里的想法是，$f$ 看起来像是在测量到某个点的距离。
  \end{enumerate}
  %
  \index{universal!property!of metric completion}%
  最后，证明度量空间完备化的如下泛性质：从一个度量空间到 Cauchy 完备度量空间的局部一致连续映射，可以唯一地延拓到该完备化上。（我们称一个映射是\define{局部一致连续的（locally uniformly continuous）}
  \indexdef{function!locally uniformly continuous}%
  \indexdef{locally uniformly continuous map}%
  如果它在每个开球上都是一致连续的。）
\end{ex}

\index{metric space|)}%

\begin{ex} \label{ex:reals-apart-neq-MP}
  \define{Markov 原理（Markov's principle）}
  \indexdef{axiom!Markov's principle}%
  \indexdef{Markov's principle}%
  说的是：对于所有 $f : \nat \to \bool$，
  %
  \begin{equation*}
    (\lnot \lnot \exis{n : \nat} f(n) = \btrue)
    \Rightarrow
    \exis{n : \nat} f(n) = \btrue.
  \end{equation*}
  %
  这是双重否定律~\eqref{eq:ldn} 的一个特例。证明
  $\fall{x, y: \RD} x \neq y \Rightarrow x \apart y$ 蕴含 Markov 原理。反之是否也成立？
\end{ex}

\begin{ex} \label{ex:reals-apart-zero-divisors}
  \index{apartness}%
  验证实数满足下述“无零因子”性质：
  $x y \apart 0 \Leftrightarrow x \apart 0 \land y \apart 0$。
\end{ex}

\begin{ex} \label{ex:finite-cover-lebesgue-number}
  %
  设 $(q_1, r_1), \ldots, (q_n, r_n)$ 逐点覆盖 $(a, b)$。则存在 $\epsilon : \Qp$，使得只要 $a < x < y < b$ 且 $|x - y| < \epsilon$，就仅仅存在 $i$ 使得 $q_i < x < r_i$ 和 $q_i < y < r_i$。这样的 $\epsilon$ 称为该覆盖的一个 \define{Lebesgue 数}
  \indexdef{Lebesgue number}。%
\end{ex}

\begin{ex} \label{ex:mean-value-theorem}
  %
  证明下述中值定理的近似版本：
  %
  \begin{quote}
    \emph{
      若 $f : [0,1] \to \R$ 是一致连续的，且 $f(0) < 0 < f(1)$，则对每个 $\epsilon : \Qp$，都仅仅存在 $x : [0,1]$ 使得 $|f(x)| < \epsilon$。
    }
  \end{quote}
  %
  提示：不要尝试使用二分法，因为它会引入选择公理。而是用一个分段线性映射来逼近 $f$。如何构造一个分段线性映射？
\end{ex}

\begin{ex}\label{ex:knuth-surreal-check}
  检查~\cite{knuth74:_surreal_number} 中的所有内容是否都能使用 \cref{sec:surreals} 中定义的高阶归纳-归纳超现实数来完成。
\end{ex}

\begin{ex}\label{ex:reals-into-surreals}
  回忆在第~\pageref{reals-into-surreals} 页定义的函数 $\iota_{\RD}:\RD\to\NO$。
  \begin{enumerate}
  \item 证明 $\iota_{\RD}$ 是单的。
  \item $\iota_{\RD}$ 可以显然地扩展到扩充实数（\cref{ex:RD-extended-reals}）和区间域（\cref{ex:RD-interval-arithmetic}）。它们是单的吗？
  \end{enumerate}
\end{ex}

\begin{ex}\label{ex:ord-into-surreals}
  证明在第~\pageref{ord-into-surreals} 页定义的函数 $\iota_{\ord}:\ord\to\NO$ 是单的当且仅当 \LEM{} 成立。
\end{ex}

\begin{ex}\label{ex:hiit-plump}
  模仿 $\NO$ 的定义（但仅使用左选项），定义一个配备二元关系 $\le$ 和 $<$ 的类型 $\mathsf{POrd}$。
  \begin{enumerate}
  \item 构造一个映射 $j:\mathsf{POrd} \to \NO$，并证明它是一个嵌入。
  \item 证明在关系 $<$ 下，$\mathsf{POrd}$ 是一个序数（在更高一层的宇宙中，如 $\ord$）。
  \item 假设命题大小重整，证明 $\mathsf{POrd}$ 等价于 \cref{ex:plump-ordinals} 中 \ord 的子集 \[\setof{A:\ord | \mathsf{isPlump}(A)}\]由此推断，$\iota_{\ord}:\ord\to\NO$ 限制在丰满序数上时是单的。
  \end{enumerate}
  在没有命题大小重整的情况下，我们依然可以称 $\mathsf{POrd}$ 的元素（或它们在 $\NO$ 中的像）为 \define{丰满序数}。\index{ordinal!plump}\index{plump!ordinal}
\end{ex}

\begin{ex}\label{ex:pseudo-ordinals}
  如果一个超现实数等于一个没有右选项的分割 $\surr{x^L}{}$（但其左选项本身可以包含右选项），则称之为 \define{伪序数}\index{pseudo-ordinal}\index{ordinal!pseudo-}。证明“每个伪序数都是丰满序数”这一陈述等价于 \LEM{}。
\end{ex}

\begin{ex}\label{ex:double-No-recursion}
  注意到 \cref{defn:No-codes} 和 \cref{eg:surreal-addition} 都使用了相似的模式来定义函数 $\NO \to \NO \to B$：一个外层的 $\NO$-递归，其陪域是保序函数 $\NO\to B$ 的集合；随后是一个内层的 $\NO$-归纳，进入一个类型族 $A:\NO\to\type$，其中 $A(y)$ 是 $B$ 的子集，用以确保不等式 $x^L<x$ 和 $x<x^R$ 也得到保持。请表述并证明一个“双重 $\NO$-递归”的一般原理，用以概括这些证明。
\end{ex}

\index{real numbers|)}%

%%% Local Variables:
%%% mode: latex
%%% TeX-master: "hott-online"
%%% End: